%% LyX 2.4.3 created this file.  For more info, see https://www.lyx.org/.
%% Do not edit unless you really know what you are doing.
\documentclass[journal,article,submit,pdftex,moreauthors]{Definitions/mdpi}
\usepackage[utf8]{inputenc}
\usepackage{array}
\usepackage{longtable}
\usepackage{float}
\usepackage{booktabs}
\usepackage{url}
\usepackage{multirow}
\usepackage{varwidth}
\usepackage{pdflscape}
\usepackage{amsbsy}
\usepackage{amstext}
\usepackage{graphicx}

\makeatletter

%%%%%%%%%%%%%%%%%%%%%%%%%%%%%% LyX specific LaTeX commands.

\Title{A Comparative Evaluation of SPARQ Against Leading Metaheuristics}

\TitleCitation{A Comparative Evaluation of SPARQ Against Leading Metaheuristics}

\Author{Vasileios Charilogis$^{1}$, Ioannis G. Tsoulos$^{2,*}$ and Anna
Maria Gianni$^{3}$}

\AuthorNames{Vasileios Charilogis, Ioannis G. Tsoulos and Anna Maria Gianni.}

\AuthorCitation{Charilogis, V., Tsoulos, I.G., Gianni, A. M.}


\address{$^{1}$\quad{}Department of Informatics and Telecommunications,
University of Ioannina, 47150 Kostaki Artas, Greece, v.charilog@uoi.gr\\
$^{2}$\quad{}Department of Informatics and Telecommunications, University
of Ioannina, 47150 Kostaki Artas, Greece, itsoulos@uoi.gr\\
$^{3}$\quad{}Department of Informatics and Telecommunications, University
of Ioannina, 47150 Kostaki Artas, Greece, am.gianni@uoi.gr}


\corres{Correspondence: itsoulos@uoi.gr}


\abstract{SPARQ is an advanced global optimization algorithm, the direct evolution
of an earlier method called ARQ2. It searches efficiently for the
best solution to problems where the space of options is too vast to
check exhaustively, such as tuning antenna arrays or planning fuel-efficient
spacecraft trajectories. Like its predecessor, it works with a population
of candidate solutions that improve generation after generation, alternating
between two complementary search strategies. What sets SPARQ apart
is that it makes nearly every part of this process adaptive. Its population
shrinks intelligently as the search matures. Its internal settings
draw from a memory of many past successful configurations, not a single
average. Its escape-from-stagnation mechanisms come in graduated strength,
from a gentle nudge to a deeper partial restart. It also adds capabilities
its predecessor never had, including a dedicated phase that locally
polishes the current best solution using its own memory of productive
directions. Every addition is kept only where it has been proven to
genuinely help. The result is markedly more reliable than its predecessor,
while staying grounded in the same battle-tested core.}


\keyword{SPARQ; Continuous optimization; Evolutionary algorithm; Adaptive
parameter control; Metaheuristics; Differential evolution; Algorithm
benchmarking; Population-based search;}

\newcommand*\LyXZeroWidthSpace{\hspace{0pt}}
\DeclareTextSymbolDefault{\textquotedbl}{T1}
%% Because html converters don't know tabularnewline
\providecommand{\tabularnewline}{\\}
%% Variable width box for table cells
\newenvironment{cellvarwidth}[1][t]
    {\begin{varwidth}[#1]{\linewidth}}
    {\@finalstrut\@arstrutbox\end{varwidth}}
\floatstyle{ruled}
\newfloat{algorithm}{tbp}{loa}
\providecommand{\algorithmname}{Algorithm}
\floatname{algorithm}{\protect\algorithmname}

%%%%%%%%%%%%%%%%%%%%%%%%%%%%%% User specified LaTeX commands.
\usepackage{rotating}
%  LaTeX support: latex@mdpi.com
%  For support, please attach all files needed for compiling as well as the log file, and specify your operating system, LaTeX version, and LaTeX editor.
%=================================================================
% For posting an early version of this manuscript as a preprint, you may use "preprints" as the journal and change "submit" to "accept". The document class line would be, e.g., \documentclass[preprints,article,accept,moreauthors,pdftex]{mdpi}. This is especially recommended for submission to arXiv, where line numbers should be removed before posting. For preprints.org, the editorial staff will make this change immediately prior to posting.
%--------------------
% Class Options:
%--------------------
%----------
% journal
%----------
% Choose between the following MDPI journals:
% acoustics, actuators, addictions, admsci, adolescents, aerospace, agriculture, agriengineering, agronomy, ai, algorithms, allergies, alloys, analytica, animals, antibiotics, antibodies, antioxidants, applbiosci, appliedchem, appliedmath, applmech, applmicrobiol, applnano, applsci, aquacj, architecture, arts, asc, asi, astronomy, atmosphere, atoms, audiolres, automation, axioms, bacteria, batteries, bdcc, behavsci, beverages, biochem, bioengineering, biologics, biology, biomass, biomechanics, biomed, biomedicines, biomedinformatics, biomimetics, biomolecules, biophysica, biosensors, biotech, birds, bloods, blsf, brainsci, breath, buildings, businesses, cancers, carbon, cardiogenetics, catalysts, cells, ceramics, challenges, chemengineering, chemistry, chemosensors, chemproc, children, chips, cimb, civileng, cleantechnol, climate, clinpract, clockssleep, cmd, coasts, coatings, colloids, colorants, commodities, compounds, computation, computers, condensedmatter, conservation, constrmater, cosmetics, covid, crops, cryptography, crystals, csmf, ctn, curroncol, currophthalmol, cyber, dairy, data, dentistry, dermato, dermatopathology, designs, diabetology, diagnostics, dietetics, digital, disabilities, diseases, diversity, dna, drones, dynamics, earth, ebj, ecologies, econometrics, economies, education, ejihpe, electricity, electrochem, electronicmat, electronics, encyclopedia, endocrines, energies, eng, engproc, ent, entomology, entropy, environments, environsciproc, epidemiologia, epigenomes, est, fermentation, fibers, fintech, fire, fishes, fluids, foods, forecasting, forensicsci, forests, foundations, fractalfract, fuels, futureinternet, futureparasites, futurepharmacol, futurephys, futuretransp, galaxies, games, gases, gastroent, gastrointestdisord, gels, genealogy, genes, geographies, geohazards, geomatics, geosciences, geotechnics, geriatrics, hazardousmatters, healthcare, hearts, hemato, heritage, highthroughput, histories, horticulturae, humanities, humans, hydrobiology, hydrogen, hydrology, hygiene, idr, ijerph, ijfs, ijgi, ijms, ijns, ijtm, ijtpp, immuno, informatics, information, infrastructures, inorganics, insects, instruments, inventions, iot, j, jal, jcdd, jcm, jcp, jcs, jdb, jeta, jfb, jfmk, jimaging, jintelligence, jlpea, jmmp, jmp, jmse, jne, jnt, jof, joitmc, jor, journalmedia, jox, jpm, jrfm, jsan, jtaer, jzbg, kidney, kidneydial, knowledge, land, languages, laws, life, liquids, literature, livers, logics, logistics, lubricants, lymphatics, machines, macromol, magnetism, magnetochemistry, make, marinedrugs, materials, materproc, mathematics, mca, measurements, medicina, medicines, medsci, membranes, merits, metabolites, metals, meteorology, methane, metrology, micro, microarrays, microbiolres, micromachines, microorganisms, microplastics, minerals, mining, modelling, molbank, molecules, mps, msf, mti, muscles, nanoenergyadv, nanomanufacturing, nanomaterials, ncrna, network, neuroglia, neurolint, neurosci, nitrogen, notspecified, nri, nursrep, nutraceuticals, nutrients, obesities, oceans, ohbm, onco, oncopathology, optics, oral, organics, organoids, osteology, oxygen, parasites, parasitologia, particles, pathogens, pathophysiology, pediatrrep, pharmaceuticals, pharmaceutics, pharmacoepidemiology, pharmacy, philosophies, photochem, photonics, phycology, physchem, physics, physiologia, plants, plasma, pollutants, polymers, polysaccharides, poultry, powders, preprints, proceedings, processes, prosthesis, proteomes, psf, psych, psychiatryint, psychoactives, publications, quantumrep, quaternary, qubs, radiation, reactions, recycling, regeneration, religions, remotesensing, reports, reprodmed, resources, rheumato, risks, robotics, ruminants, safety, sci, scipharm, seeds, sensors, separations, sexes, signals, sinusitis, skins, smartcities, sna, societies, socsci, software, soilsystems, solar, solids, sports, standards, stats, stresses, surfaces, surgeries, suschem, sustainability, symmetry, synbio, systems, taxonomy, technologies, telecom, test, textiles, thalassrep, thermo, tomography, tourismhosp, toxics, toxins, transplantology, transportation, traumacare, traumas, tropicalmed, universe, urbansci, uro, vaccines, vehicles, venereology, vetsci, vibration, viruses, vision, waste, water, wem, wevj, wind, women, world, youth, zoonoticdis
%---------
% article
%---------
% The default type of manuscript is "article", but can be replaced by:
% abstract, addendum, article, book, bookreview, briefreport, casereport, comment, commentary, communication, conferenceproceedings, correction, conferencereport, entry, expressionofconcern, extendedabstract, datadescriptor, editorial, essay, erratum, hypothesis, interestingimage, obituary, opinion, projectreport, reply, retraction, review, perspective, protocol, shortnote, studyprotocol, systematicreview, supfile, technicalnote, viewpoint, guidelines, registeredreport, tutorial
% supfile = supplementary materials
%----------
% submit
%----------
% The class option "submit" will be changed to "accept" by the Editorial Office when the paper is accepted. This will only make changes to the frontpage (e.g., the logo of the journal will get visible), the headings, and the copyright information. Also, line numbering will be removed. Journal info and pagination for accepted papers will also be assigned by the Editorial Office.
%------------------
% moreauthors
%------------------
% If there is only one author the class option oneauthor should be used. Otherwise use the class option moreauthors.
%---------
% pdftex
%---------
% The option pdftex is for use with pdfLaTeX. If eps figures are used, remove the option pdftex and use LaTeX and dvi2pdf.
%=================================================================
% MDPI internal commands - do not modify
\firstpage{1}
\setcounter{page}{\@firstpage}
\pubvolume{1}
\issuenum{1}
\articlenumber{0}
\pubyear{2024}
\copyrightyear{2024}
%\externaleditor{Academic Editor: Firstname Lastname} % For journal Automation, please change Academic Editor to "Communicated by"
\datereceived{}
\daterevised{ } % Comment out if no revised date
\dateaccepted{}
\datepublished{}
%\datecorrected{} % Corrected papers include a "Corrected: XXX" date in the original paper.
%\dateretracted{} % Corrected papers include a "Retracted: XXX" date in the original paper.
\hreflink{} % If needed use \linebreak
%\doinum{}
%------------------------------------------------------------------
% The following line should be uncommented if the LaTeX file is uploaded to arXiv.org
%\pdfoutput=1
%=================================================================
% Add packages and commands here. The following packages are loaded in our class file: fontenc, inputenc, calc, indentfirst, fancyhdr, graphicx, epstopdf, lastpage, ifthen, lineno, float, amsmath, setspace, enumitem, mathpazo, booktabs, titlesec, etoolbox, tabto, xcolor, soul, multirow, microtype, tikz, totcount, changepage, attrib, upgreek, cleveref, amsthm, hyphenat, natbib, hyperref, footmisc, url, geometry, newfloat, caption
%=================================================================
%% Please use the following mathematics environments: Theorem, Lemma, Corollary, Proposition, Characterization, Property, Problem, Example, ExamplesandDefinitions, Hypothesis, Remark, Definition, Notation, Assumption
%% For proofs, please use the proof environment (the amsthm package is loaded by the MDPI class).
%=================================================================
% The fields PACS, MSC, and JEL may be left empty or commented out if not applicable
%\PACS{J0101}
%\MSC{}
%\JEL{}
%%%%%%%%%%%%%%%%%%%%%%%%%%%%%%%%%%%%%%%%%%
% Only for the journal Diversity
%\LSID{\url{http://}}
%%%%%%%%%%%%%%%%%%%%%%%%%%%%%%%%%%%%%%%%%%
% Only for the journal Applied Sciences:
%\featuredapplication{Authors are encouraged to provide a concise description of the specific application or a potential application of the work. This section is not mandatory.}
%%%%%%%%%%%%%%%%%%%%%%%%%%%%%%%%%%%%%%%%%%
%%%%%%%%%%%%%%%%%%%%%%%%%%%%%%%%%%%%%%%%%%
% Only for the journal Data:
%\dataset{DOI number or link to the deposited data set in cases where the data set is published or set to be published separately. If the data set is submitted and will be published as a supplement to this paper in the journal Data, this field will be filled by the editors of the journal. In this case, please make sure to submit the data set as a supplement when entering your manuscript into our manuscript editorial system.}
%\datasetlicense{license under which the data set is made available (CC0, CC-BY, CC-BY-SA, CC-BY-NC, etc.)}
%%%%%%%%%%%%%%%%%%%%%%%%%%%%%%%%%%%%%%%%%%
% Only for the journal Toxins
%\keycontribution{The breakthroughs or highlights of the manuscript. Authors can write one or two sentences to describe the most important part of the paper.}
%%%%%%%%%%%%%%%%%%%%%%%%%%%%%%%%%%%%%%%%%%
% Only for the journal Encyclopedia
%\encyclopediadef{Instead of the abstract}
%\entrylink{The Link to this entry published on the encyclopedia platform.}
%%%%%%%%%%%%%%%%%%%%%%%%%%%%%%%%%%%%%%%%%%
%%%%%%%%%%%%%%%%%%%%%%%%%%%%%%%%%%%%%%%%%%
% Only for the journal Advances in Respiratory Medicine
%\addhighlights{yes}
%\renewcommand{\addhighlights}{%
%\noindent This is an obligatory section in “Advances in Respiratory Medicine”, whose goal is to increase the discoverability and readability of the article via search engines and other scholars. Highlights should not be a copy of the abstract, but a simple text allowing the reader to quickly and simplified find out what the article is about and what can be cited from it. Each of these parts should be devoted up to 2~bullet points.\vspace{3pt}\\
%\textbf{What are the main findings?}
% \begin{itemize}[labelsep=2.5mm,topsep=-3pt]
% \item First bullet.
% \item Second bullet.
% \end{itemize}\vspace{3pt}
%\textbf{What is the implication of the main finding?}
% \begin{itemize}[labelsep=2.5mm,topsep=-3pt]
% \item First bullet.
% \item Second bullet.
% \end{itemize}
%}
%%%%%%%%%%%%%%%%%%%%%%%%%%%%%%%%%%%%%%%%%%
% Added by lyx2lyx
\usepackage{array}
\usepackage[noend]{algpseudocode}
\algrenewcommand\algorithmicindent{1.0em}

\makeatother

\begin{document}
\maketitle

\section{Introduction}

Optimization problems in which the goal is to locate the best possible
configuration of a set of continuous variables, without relying on
gradient information and without any guarantee that the objective
function is smooth, convex, or even continuous, arise throughout engineering,
physics, and applied economics. Formally, such problems are instances
of global optimization:

\[
\min_{\mathbf{x}\in\Omega}f(\mathbf{x}),\qquad\Omega:=[\ell_{1},u_{1}]\times\dots\times[\ell_{D},u_{D}]\subset\mathbb{R}^{D},
\]

where $f:\Omega\rightarrow\mathbb{R}$ is accessed only through pointwise
evaluations, and derivative information may be unavailable. Given
an evaluation budget $N_{fe}$\LyXZeroWidthSpace , an algorithm produces
a sequence of iterates $\{\mathbf{x}_{t}\}_{t=1}^{N_{fe}}$ and tracks
the best-so-far value

\[
f_{\min}(t):=\min_{1\le s\le t}f(\mathbf{x}_{s})
\]

The task is to locate the global minimum (or maximum) of $f$ over
$\Omega$ despite the presence of numerous local optima, ill-conditioning,
and, in many real applications, a strict limit $N_{fe}$ on how many
times the objective function can be evaluated. The difficulty of this
task grows sharply with the dimension $D$, and no single search strategy
can be expected to perform best on every possible problem instance
\citep{problem}.

In response, the optimization literature has developed several broad
families of methods. Exact and gradient-based methods offer strong
convergence guarantees but require differentiability assumptions that
most real-world objective functions do not satisfy. Classical local
random search methods, most notably the algorithm of Solis and Wets
\citep{SolisAndWets}, instead perturb a single candidate solution
and accept improving moves, offering simplicity at the cost of a tendency
to stall in local optima. To escape this limitation, the field turned
to population-based metaheuristics, which explore the search space
with many candidate solutions simultaneously. Early and influential
examples include simulated annealing \citep{SA}, genetic algorithms
\citep{GA}, together with more recent refinements to their initialization
strategies such as k-means-based seeding \citep{GiantArmadillo},
particle swarm optimization \citep{PSO}, and the Covariance Matrix
Adaptation Evolution Strategy (CMA-ES) \citep{CMAES}, later extended
to large-scale problems via limited-memory variants such as LM-CMA-ES
\citep{LMCMAES}. More recently, a growing number of bio-inspired
metaheuristics have been proposed and subsequently refined through
hybridization with local search and modern termination criteria, as
in the case of the Giant-Armadillo Optimization method \citep{InitialDistribution}.

Among population-based metaheuristics, Differential Evolution (DE),
introduced by Storn and Price \citep{DE} and since surveyed extensively
\citep{DE2}, has proven especially effective and has spawned an active
line of successive refinements. JADE \citep{JADE} introduced the
current-to-pbest mutation strategy together with an external archive
of recently displaced solutions, both of which became standard components
of later variants. SHADE \citep{SHADE} replaced JADE's single adaptive
parameter pair with a historical memory of successful settings, and
L-SHADE \citep{LSHADE} combined this with a linear population size
reduction scheme building on earlier work on population sizing \citep{DEd}.
jSO \citep{JSO} refined this further with progress-dependent parameter
schedules, and subsequent work introduced rank-based selective pressure
in choosing difference vectors \citep{LSHADE2}, culminating in NL-SHADE-RSP
\citep{NL-SHADE-RSP}, which combined non-linear population size reduction
with rank-based selection and adaptive archiving. A parallel line
of research explored hybridizing DE with other search paradigms and
enriching its crossover operator: CoBiDE \citep{CoBiDE} introduced
covariance-guided, bimodal parameter sampling; individual-dependent
DE variants (IDE) proposed alternative generation-based scheduling
\citep{DEp}, eigenvector-based (\textquotedbl eigen\textquotedbl )
crossover, aligning the crossover operation with the population's
natural axes of variation rather than the raw coordinate axes, was
introduced in \citep{SPS-L-SHADE-EIG} and later integrated into a
cooperative multi-algorithm framework, EA4Eig, that combines CoBiDE,
IDE, CMA-ES, and jSO under a shared eigen-crossover mechanism \citep{EI4Eig}.
Separately, other general-purpose mechanisms have been proposed for
escaping stagnation and premature convergence, including opposition-based
learning \citep{obl}, \citep{obde}, Lévy-flight perturbations based
on the numerical simulation method of Mantegna \citep{Mantegna},
and, from the domain of sequential decision-making, Thompson sampling
as a principled method for adaptively allocating trials between competing
strategies \citep{Thompson}.

The method presented in this article, SPARQ, is the latest member
of a family of Differential-Evolution-based optimizers developed internally,
beginning with ARQ \citep{ARQ} and its direct successor ARQ2 \citep{ARQ2},
both of which combine a current-to-pbest DE mutation strategy with
an alternative individual-dependent evolution strategy (IDE) selected
via an adaptive roulette mechanism, restricted tournament replacement,
an external archive, and outlier quarantine. SPARQ builds directly
on ARQ2's architecture while replacing nearly every one of its fixed
or singly-adapted mechanisms with an adaptive counterpart drawn from
the DE literature surveyed above, among them a SHADE-style historical
memory in place of a single running parameter average, non-linear
population size reduction, rank-based selective pressure in the style
of NL-SHADE-RSP, a jSO-style progress-dependent schedule for its elitist
pressure, eigen-coordinate crossover, a Thompson-sampling bandit in
place of frequency-based strategy selection, and Lévy-flight-based
quarantine perturbations. On top of this synthesis, SPARQ introduces
two mechanisms with no counterpart in ARQ2 or the algorithms above:
a stagnation-gated local polishing phase based on the classical Solis-Wets
adaptive random search, and a short-term memory of recently successful
search directions (\textquotedbl trajectory echo\textquotedbl )
that is replayed during this local phase.

The remainder of this article is organized as follows. Section \ref{subsec:ARQandARQ2}
describes the architecture inherited from ARQ2 in detail. Section
\ref{sec:sparq} presents each of SPARQ's adaptive and novel mechanisms
in turn, together with the motivation behind each. Section 4 describes
the experimental configuration and reports and discusses the results,
comparing SPARQ against nine other state-of-the-art optimizers on
the standard-dimensionality benchmark and against five large-scale
optimizers on the high-dimensional benchmark, together with an analysis
of computational scalability with dimension. Section \ref{sec:conclusions}
concludes.

\section{Background: The ARQ and ARQ2 Predecessor Methods\label{subsec:ARQandARQ2}}

\subsection{ARQ: A Cohesive Optimization Design with Outlier Quarantine\label{subsec:ARQ}}

ARQ \citep{ARQ} is a differential-evolution-based optimizer built
around a single, cohesive update cycle that combines four complementary
mechanisms.

Archive-assisted pbest/1/bin trial generation. Each new candidate
is generated by attracting a target vector toward one of the top-p
fraction of the population, plus a difference term drawn partly from
the live population and partly from an external archive of recently
displaced solutions:

\[
v=x+F\,(x_{pbest}-x)+F\,(r_{1}-r_{2}),\qquad r_{1}\in P_{t},\ r_{2}\in P_{t}\cup A_{t}.
\]

This is followed by binomial crossover and projection back onto the
feasible box.

Success-history parameter adaptation. The scale factor $F$ and crossover
rate $CR$ are drawn from a Cauchy and a Normal distribution, respectively,
centered on running means $\mu_{F}$\LyXZeroWidthSpace , $\mu_{CR}$\LyXZeroWidthSpace .
After each cycle, these means are updated using gain-weighted statistics,
a weighted Lehmer mean for F and a weighted arithmetic mean for $CR$,
so that parameter settings that produced larger fitness gains are
trusted more.

Restricted Tournament Replacement (RTR). A trial vector first competes
against its own parent, if it fails to improve, it is instead compared
against its nearest neighbor (in bounds-normalized distance) within
a small random pool drawn from the population, and replaces that neighbor
if superior. This localizes selection pressure and helps preserve
population diversity.

Outlier quarantine and targeted micro-restart. ARQ's central contribution
is an event-driven mechanism that detects fitness outliers using a
robust interquartile-range threshold $\theta=Q3+a\cdot IQR$ and gently
relocates a fraction of them toward a robust center (the mean of the
best half of the population) under a small perturbation, accepting
the repair only if it improves fitness. A complementary micro-restart,
triggered after a fixed number of non-improving iterations, perturbs
a small fraction of the worst individuals around the current incumbent.
Together, these two mechanisms convert destabilizing outliers into
controlled exploration stimuli, narrowing the gap between best-case
and average-case performance without sacrificing peak results.

\subsection{ARQ2: A Roulette-Guided Hybrid Extension of ARQ\label{subsec:ARQ2}}

ARQ2 \citep{ARQ2} preserves this entire stability-oriented backbone
but extends it into an asymmetric hybrid architecture in which two
internal search branches, ARQ and an auxiliary IDE (Individual-dependent
Differential Evolution) branch, operate on a shared population and
a shared external archive.

Roulette-based branch scheduling. At each iteration, a controller
selects which of the two branches to execute, based on branch \textquotedbl credits\textquotedbl{}
that are updated in proportion to each branch's recent number of successful
replacements. Branches with higher recent credit are selected more
often, but a reset mechanism restores both branches to a common baseline
credit whenever either branch's selection probability falls below
a threshold $\delta$, preventing one branch from permanently dominating
and collapsing the hybrid into a single-strategy method. An initial
bootstrap phase enforces the ARQ branch for a fixed number of early
iterations, before branch competition becomes informative.

The IDE branch. When selected, this branch constructs its donor vectors
using its own, per-individual memories of $F$ and $CR$, and a progress-dependent
elite-guided pool whose size and guidance strength shift over the
course of the search, favoring stronger elite guidance early on, and
a more balanced elite/random mixture later. It further switches between
two donor forms (guidance-based and current-based) depending on a
phase threshold that adapts if the branch's success rate stalls.

An enriched ARQ branch. When the ARQ branch is selected, its mutation
rule is enriched with a jSO-style, progress-dependent attraction coefficient
$K(\pi)$ that modulates how strongly each new candidate is pulled
toward an elite exemplar, growing more elitist as the search progresses.

Asymmetric stabilization. Importantly, the quarantine and micro-restart
mechanisms inherited from ARQ remain attached only to the ARQ branch.
This preserves the original design's corrective role without duplicating
the same machinery inside the auxiliary IDE branch, while still letting
IDE contribute a structurally distinct search dynamic that helps on
multimodal landscapes.

Empirically, this hybridization was shown to primarily strengthen
repeated-run reliability: across 36 benchmark problems, ARQ2 improved
on ARQ's mean-value performance on 21 problems (matching it on 11,
and trailing on only 4), while remaining highly competitive in best-case
terms \citep{ARQ2}.

\section{The Proposed Method: SPARQ\label{sec:sparq}}

\subsection{The mechanisms of the SPARQ\label{subsec:mechanisms}}

SPARQ is built directly on top of ARQ2's architecture (Section \ref{subsec:ARQ2}):
it retains the shared population, the shared external archive, and
the two-branch structure (ARQ vs. IDE), but it replaces nearly every
one of ARQ2's fixed or singly-adapted mechanisms with an adaptive
counterpart, and adds several mechanisms with no counterpart in ARQ2
at all. This section describes each addition in turn.

\textbf{Population State and Non-Linear Population Size Reduction
(NLPSR)}\label{subsec:NLPSR}

Unlike ARQ2, which runs with a fixed population size $N$ throughout,
SPARQ starts from an initial size $N_{init}$ (by default scaled with
dimension, $N_{init}=18D$) and shrinks it toward a floor $N_{min}$
as the search progresses, following

\[
N(t)=\text{round}\Big(N_{init}+(N_{min}-N_{init})\cdot p(t)^{,e(t)}\Big),\qquad e(t)=1-(1-\alpha_{\text{NL}})\cdot p(t),
\]
where $p(t)\in[0,1]$ is the fraction of the evaluation budget consumed
so far and $\alpha_{\text{NL}}\in(0,1]$ controls the shape of the
reduction: $\alpha_{\text{NL}}=1$ recovers a purely linear LSHADE-style
schedule, while $\alpha_{\text{NL}}<1$produces a near-linear reduction
early in the search that accelerates as the budget is exhausted, in
the spirit of NL-SHADE-RSP. At the end of every iteration, if the
current population exceeds the target size, the worst individuals
are discarded, freeing up evaluations for the remaining ones. The
cached eigenbasis (Section \ref{subsec:EigenCrossover}) is invalidated
whenever the population shrinks, since it must be recomputed on the
new, smaller sample.

For the comparative evaluation reported in Section \ref{sec:ExperimentalResults},
however, $N_{init}$\LyXZeroWidthSpace{} was fixed at 100 for SPARQ
and for every other population-based competitor in the comparison,
overriding the dimension-scaled default $N_{init}=18D$. This choice
ensures that all compared methods start from an identical population
size rather than from method-specific, dimension-dependent values,
so that observed performance differences reflect the search mechanisms
themselves rather than differences in initial population sizing. The
dimension-scaled schedule $N_{init}=18D$ remains SPARQ's default,
recommended setting outside of this controlled comparison.

\textbf{SHADE-Style Historical Memory for F and CR}\label{subsec:3.2. SHADEStyle}

ARQ2 keeps a single, continuously-updated pair of running means ($\mu_{F},\mu_{CR}$).
SPARQ instead maintains a circular memory of size $H$, $\{(M_{F}^{k},M_{CR}^{k})\}_{k=1}^{H}$,
one slot of which (the \textquotedbl terminal\textquotedbl{} slot,
following jSO) is permanently pinned to a high value ($M_{F}^{H}=M_{CR}^{H}=0.9$)
rather than being overwritten. Before generating each trial, a memory
slot $r$ is drawn uniformly at random, and
\[
CR\sim\mathcal{N}(M_{CR}^{r},0.1),\qquad F\sim\text{Cauchy}(M_{F}^{r},0.1),
\]
both resampled/clamped to their valid ranges ($CR\in[0,1]CR\in[0,1]CR\in[0,1],F\in[F_{lo},F_{hi}]$).
Two jSO-style schedules further constrain these draws early in the
search: $CR$ is floored at 0.7 before 20\% progress and at 0.6 before
50\% progress, and $F$ is capped at $0.7+0.3\,p(t)$ before 60\%
progress. After each generation, the (non-terminal) memory slots are
updated cyclically using the same gain-weighted Lehmer/arithmetic
mean rule as ARQ2, but only one slot advances per generation rather
than the single running pair, so that a diverse set of previously-successful
settings remains available for future sampling rather than being collapsed
into one blended average.

\textbf{Rank-Based Selective Pressure and jSO-Style Elitism Schedule}\label{subsec:RankBased}

ARQ2 draws the difference-vector companions $r_{1}$ (and, when not
drawn from the archive, $r_{2}$) uniformly at random from the population.
SPARQ instead biases this choice toward better-ranked individuals:
given the population sorted by fitness (rank 0 = best), an index at
rank $r$ is selected with probability proportional to $(N-r)^{k_{r}}$\LyXZeroWidthSpace ,
where the selective-pressure exponent itself increases linearly over
the run,
\[
k_{r}(t)=k_{r,init}+(k_{r,final}-k_{r,init})\cdot p(t),\qquad k_{r,init}=2,\ k_{r,final}=3,
\]
so that rank bias is moderate early on and sharpens later, in the
style of LSHADE-RSP. Independently, the size of the $p$best pool
from which the elite-attraction target $x_{pbest}$ is drawn also
narrows over the course of the run, following a jSO-style linear schedule
from $p_{max}=0.25$ down to $p_{min}=0.10$, so that the elitist
pull becomes progressively more selective as the search matures. A
further NL-SHADE-RSP-style refinement, $CR$-sorting, reassigns the
sampled $CR$ values within each generation's batch of parents so
that better-ranked parents receive the smallest values (conservative,
fine-grained crossover) and worse-ranked parents the largest (more
disruptive recombination).

\textbf{Eigen-Coordinate Binomial Crossover}\label{subsec:EigenCrossover}

ARQ2's binomial crossover operates coordinate-by-coordinate in the
raw problem coordinates, which is ill-suited to correlated, non-separable
landscapes whose productive directions run diagonally across several
variables at once. SPARQ periodically estimates the dominant axes
of the current population's spread: using the best 50\% of individuals
(by fitness), it computes their empirical covariance matrix, diagonalizes
it via Jacobi rotation to obtain an eigenbasis $B$ and per-axis scale
factors, and refreshes this basis every $\texttt{eig\_period}$ iterations
(or whenever the population is resorted or resized). With probability
$p_{eig}=0.40$, and whenever $D$ stays within a practical cap ($D\le100$,
since the Jacobi decomposition is $\mathcal{O}(D^{3})$ per sweep),
the binomial crossover of Equation (10) in ARQ2's formulation is instead
applied in the rotated coordinate system: both the parent and the
mutant are projected onto the eigenbasis ($x_{e}=B^{\top}x$), crossover
is performed there, and the result is rotated back ($u=B\,u_{e}$).
This lets the operator recombine along the population's natural axes
of variation rather than the arbitrary problem axes.

\textbf{Thompson Sampling Bandit for Strategy Selection}\label{subsec:Thompson}

ARQ2 selects between the ARQ and IDE branches using a frequency-based
\textquotedbl credit\textquotedbl{} roulette with a threshold-triggered
reset. SPARQ replaces this with a genuine two-armed Bayesian bandit:
each branch $k$ maintains a Beta posterior with parameters ($a_{k},b_{k}$),
initialized with a mild prior favoring ARQ ($a_{ARQ}=2$) to provide
brief warm-up bias before a short deterministic bootstrap (a small
fixed number of forced ARQ iterations) hands control to the bandit.
At each iteration, a sample is drawn from each arm's posterior via
two Gamma draws, and the arm with the larger sample is selected. After
a branch executes, its posterior is updated in proportion to its observed
success rate, normalized by the number of trial evaluations actually
attempted so that a large IDE sweep and a smaller ARQ sweep are weighted
comparably, between iterations, both posteriors are geometrically
decayed toward their uninformative prior (1,1) to accommodate a non-stationary
environment, in which a strategy that was previously effective may
cease to be so as the population and progress evolve.

\textbf{Lévy-Flight Quarantine}\label{subsec:L=0000E9vyFlight}

ARQ2's quarantine mechanism relocates fitness outliers (detected via
the interquartile threshold $\theta=Q_{3}+\alpha\cdot IQR$) toward
the centroid of the top half of the population using an ordinary Gaussian
perturbation. SPARQ instead perturbs each repaired outlier using a
symmetric Lévy-stable step, generated via the Mantegna (1994) algorithm:
a step $u/|v|^{1/\beta}$ is formed from two correlated Gaussian draws
$u$, $v$ with $\beta=1.5$, producing a distribution that is mostly
small but occasionally very large, mimicking the search pattern of
animals foraging under scarce resources. Applied per coordinate and
scaled by the box range, this gives quarantined individuals a genuine
chance of relocating to a distant, unexplored region rather than merely
oscillating near the population centroid.

\textbf{On-Demand Opposition-Based Basin Escape}\label{subsec:ODOBB}

This mechanism has no counterpart in ARQ2. It monitors two independent
signals: the number of consecutive iterations without any improvement
in the incumbent best ($no\_improve\ge stag\_trigger$), and the population's
normalized spread (the average, per-coordinate, box-normalized standard
deviation across the population) falling below a small collapse ratio.
When both conditions hold simultaneously, and a cooldown counter has
elapsed, a fraction of the worst individuals are relocated to points
computed as a quasi-opposition of their current position, either the
exact reflection through the center of the box, or a point interpolated
between that reflection and the current incumbent best, and the relocation
is accepted only if it improves fitness. This directly targets premature
convergence: rather than waiting for a general-purpose restart, it
intervenes as soon as the population has visibly collapsed onto a
single region while genuinely failing to progress, forcing a look
at the opposite side of the search space.

\textbf{Stagnation-Gated Elite Polish}\label{subsec:StagnationPolish}

This is the most substantial addition with no analogue in ARQ2. Whenever
the incumbent best has failed to improve by a meaningful relative
margin for a number of iterations (this trigger tightens to a single
stagnant iteration in the closing quarter of the budget), SPARQ enters
a short, budget-capped burst (never more than a fixed fraction, by
default 12\%, of the evaluations consumed so far) of local refinement
attempts centered on the current best $x_{best}$. Each attempt draws
one of four probe modes at random:
\begin{enumerate}
\item Single-coordinate probing (round-robin over dimensions, so that every
coordinate is eventually visited even in high dimension): either a
coordinate-wise opposition ($x_{j}\to\ell_{j}+u_{j}-x_{j}$) or an
adaptive Gaussian perturbation of that one coordinate
\item Anisotropic eigen-aligned steps, drawn as a Gaussian in the (possibly
slightly stale) eigenbasis of Section 3.4, scaled per-axis by the
square root of the corresponding eigenvalue, so that steps are long
along the population's natural \textquotedbl valley\textquotedbl{}
and short across it
\item Trajectory echo (described in Section \ref{subsec:TrajectoryEcho});
\item Solis-Wets adaptive random search (described in Section \ref{SolisWets}
below, integrated as the fourth probe mode).
\end{enumerate}
Each successful probe (one that improves $best_{f}$) is injected
into the population by overwriting its current worst member, propagating
the refinement into the DE core's own mutation pool, probe-mode-specific
step sizes are adapted by a 1/5-success rule (grown on success, shrunk
on failure). The whole mechanism is protected by three independent
safety nets: a hard evaluation-budget cap, an exponential cooldown
backoff whenever an entire burst fails to find any improvement, and
an efficiency arbitration that compares the gain-per-evaluation achieved
by the polish against the gain-per-evaluation achieved by the ordinary
evolutionary loop over the same interval, forcing the polish to stand
down whenever the core DE search is demonstrably the better investment
of evaluations. A landscape-level self-selection rule additionally
imposes a long cooldown whenever several consecutive bursts produce
only negligible relative gains, so that the mechanism naturally becomes
inactive on landscapes (such as those with a distant global optimum
reachable only through large exploratory jumps) where local refinement
of the current incumbent is not productive.

\textbf{Trajectory Echo: A Memory of Recently Successful Directions}\label{subsec:TrajectoryEcho}

Also with no counterpart in ARQ2: whenever the DE core (either the
ARQ branch's restricted-tournament acceptance or the IDE branch's
greedy replacement) accepts a strictly improving trial, the accepted
displacement (trial vector minus the individual it replaced) is pushed
into a small, fixed-size ring buffer of at most eight entries. As
one of the elite-polish probe modes (Section \ref{subsec:StagnationPolish}),
a candidate step is formed as a random linear combination of the buffered
directions, normalized so that its expected magnitude stays comparable
to a single one of its ingredients regardless of how many are currently
buffered. This lets the polish phase replay an echo of movements that
have already proven productive during this specific run, of particular
value on problems where genuine improvement requires shifting several
coordinates together in a coordinated way, a pattern that isotropic
or single-coordinate probing cannot easily discover. Like the other
probe modes, this mechanism carries its own circuit breaker: a run
of repeated failures disables it for the remainder of the search rather
than continuing to spend evaluations on a direction source that has
stopped paying off.

\textbf{Solis-Wets Adaptive Local Search}\label{SolisWets}

The fourth elite-polish probe mode is a single running trajectory,
based on the classical adaptive random search method of Solis and
Wets (1981), anchored at the current incumbent. A bias vector $b$
(initialized at zero) accumulates a memory of productive directions.
Each attempt draws a random perturbation $\delta$ and first tries
$x_{best}+b+\delta$; if that fails to improve, it tries the exact
reflection $x_{best}-b-\delta$ using the same random draw before
giving up. A success in the biased direction grows $b$ toward that
direction and expands the step scale, a success only in the reflected
direction shrinks $b$ and partially reverses it; and a failure in
both directions decays $b$ toward zero and contracts the step, the
classical Solis-Wets rule for avoiding wasted wandering. Because this
probe mode only ever evolves a single trajectory from the current
incumbent, rather than recombining information from two different
population members, it cannot produce the kind of structurally incoherent
points that a population-level recombination operator could produce
on permutation-symmetric landscapes.

\textbf{Hard-Stagnation Rejuvenation}\label{subsec:HardStagnationRejuvenation}

ARQ2's single micro-restart mechanism is replaced in SPARQ by a two-tier
system. A survival trigger fires whenever the population's fitness
spread has collapsed to near-zero (a looser condition than any typical
similarity-based stopping rule), and re-seeds only a small sliver
(about 5\%) of the population uniformly at \ref{subsec:HardStagnationRejuvenation}random,
at each iteration if needed; its purpose is purely to keep the population
from appearing fully converged to an external stopping rule while
genuine progress is still possible. A hard-stagnation trigger fires
only once the incumbent has failed to improve for a much longer interval
(a configurable multiple of the basin-escape threshold of Section
\ref{subsec:ODOBB}) and the population has also collapsed, when it
fires, it re-seeds a much larger fraction of the population (by default,
all but the top 25\%) uniformly at random, while explicitly preserving
the incumbent best itself, which is never part of the population and
is therefore never at risk. An escalation mechanism further tracks
whether repeated firings of either tier keep failing to produce meaningful
gains: a persistent survival-tier failure forces a hard-stagnation
escape early, and a persistent hard-stagnation failure escalates to
a near-complete reset (keeping only a small safety fraction), preventing
a strong-but-insufficient partial restart from being immediately \textquotedbl reconquered\textquotedbl{}
by an entrenched elite on unusually deceptive landscapes. The SHADE
memory of Section \ref{subsec:3.2. SHADEStyle} is deliberately preserved
rather than reset across a rejuvenation event, since it continues
to encode useful parameter statistics for the refreshed population.

\textbf{Overall Control Flow}\label{subsec:OverallControlFlow}

Each SPARQ iteration proceeds as follows. The archive is trimmed to
its capacity; a branch (ARQ or IDE) is selected, either by the initial
bootstrap or by the Thompson bandit of Section \ref{subsec:Thompson},
and executed, with the ARQ branch additionally followed by Lévy-flight
quarantine (Section \ref{subsec:L=0000E9vyFlight}). Two independent
stagnation counters are then updated, one tracking any improvement
at all (feeding the basin-escape and rejuvenation triggers), and one
tracking only meaningful relative improvement (feeding the elite-polish
trigger, so that large populations at high dimensionality, where some
individual improves the incumbent by a vanishing amount almost every
generation, do not starve the polish mechanism of activations). The
elite polish (Section \ref{subsec:StagnationPolish}), the basin escape
(Section \ref{subsec:ODOBB}), and the rejuvenation mechanism (Section
\ref{subsec:HardStagnationRejuvenation}) are then each given the
opportunity to fire, in that order, subject to their own independent
cooldowns and triggers. The population is then shrunk toward its current
NLPSR target size if needed (Section \ref{subsec:NLPSR}), and the
bandit posteriors are decayed before the next iteration begins.

\subsection{Algorithmic Summary\label{subsec:=000391lgorithmicSummary}}

\begin{algorithm}[H]
{\footnotesize INPUT}{\footnotesize\par}

{\footnotesize - $f$: objective function}{\footnotesize\par}

{\footnotesize - $\Omega\subset\mathbb{R}^{D}$: box-constrained search
domain, dimension D}{\footnotesize\par}

{\footnotesize - $N_{\mathrm{init}},\ N_{\mathrm{min}},\ \alpha_{\mathrm{NL}}$:
NLPSR: initial/minimum population size and shrink-shape exponent}{\footnotesize\par}

{\footnotesize - $T_{\mathrm{eval}}$: maximum number of function evaluations}{\footnotesize\par}

{\footnotesize - $H,\ M_{F},\ M_{CR},\ F_{lo},\ F_{hi}$: SHADE memory
size, circular memories, and F bounds}{\footnotesize\par}

{\footnotesize - $p_{\max},\ p_{\min}$: jSO-style pbest fraction schedule}{\footnotesize\par}

{\footnotesize - $k_{r}^{0},\ k_{r}^{1}$: LSHADE-RSP rank-selection
pressure schedule}{\footnotesize\par}

{\footnotesize - $p_{\mathrm{eig}},\ T_{\mathrm{eig}},\ \rho_{\mathrm{eig}}$:
eigen-crossover probability, refresh period, and top-fraction}{\footnotesize\par}

{\footnotesize - $\alpha_{A}$: archive-rate coefficient}{\footnotesize\par}

{\footnotesize - $\{(a_{k},b_{k})\}_{k=1}^{2},\ \gamma_{b},\ n_{bs}$:
Thompson bandit priors (ARQ/IDE), decay rate, bootstrap length}{\footnotesize\par}

{\footnotesize - $\alpha_{q},\ \rho_{q},\ \beta_{L},\ \sigma_{L}$:
Levy quarantine: IQR coefficient, repair fraction, Levy exponent,
step scale}{\footnotesize\par}

{\footnotesize - $\tau_{s},\ \varepsilon_{v},\ \rho_{b},\ \kappa_{b}$:
OBL basin escape: stagnation trigger, spread-collapse ratio, fraction,
cooldown}{\footnotesize\par}

{\footnotesize - $\tau_{p},\ \pi_{p},\ \rho_{p},\ B_{p}$: elite polish:
trigger, late-phase threshold, burst fraction, budget cap}{\footnotesize\par}

{\footnotesize - $m_{r},\ \kappa_{r},\ C_{r}$: rejuvenation: stagnation
factor, keep-fraction, cooldown}{\footnotesize\par}

{\footnotesize OUTPUT}{\footnotesize\par}

{\footnotesize - $x^{\ast},\ f^{\ast}$}{\footnotesize\par}

{\footnotesize INITIALIZATION}{\footnotesize\par}

{\footnotesize - Set $N\leftarrow N_{\mathrm{init}}$, initialize population
$X=\{x_{i}\}_{i=1}^{N}\subset\Omega$ and evaluate all individuals}{\footnotesize\par}

{\footnotesize - Set archive $A\leftarrow\varnothing$}{\footnotesize\par}

{\footnotesize - Identify the incumbent best pair $(x^{\ast},f^{\ast})$}{\footnotesize\par}

{\footnotesize - Set $f^{\mathrm{prev}}\leftarrow f^{\ast}$, $s\leftarrow0$,
$s_{p}\leftarrow0$}{\footnotesize\par}

{\footnotesize - Initialize the SHADE memory $(M_{F},M_{CR})$ with
the terminal slot fixed at 0.9}{\footnotesize\par}

{\footnotesize - Initialize the bandit state $(a_{1},b_{1})\leftarrow(2,1)$,
$(a_{2},b_{2})\leftarrow(1,1)$, and the bootstrap counter $n_{bs}$}{\footnotesize\par}

{\footnotesize - Initialize the individual IDE memories $\{F_{i}^{\mathrm{IDE}},CR_{i}^{\mathrm{IDE}}\}_{i=1}^{N}$}{\footnotesize\par}

{\footnotesize - Invalidate the eigen-basis and set $e\leftarrow N$}{\footnotesize\par}

{\footnotesize SPARQ main pseudocode}{\footnotesize\par}

{\footnotesize 01 While $e<T_{\mathrm{eval}}$ do}{\footnotesize\par}

{\footnotesize 02 \hspace{0.5cm}Trim $A$ to size $\lfloor\alpha_{A}N\rfloor$}{\footnotesize\par}

{\footnotesize 03 \hspace{0.5cm}If $n_{bs}>0$ then}{\footnotesize\par}

{\footnotesize 04 \hspace{0.5cm}\hspace{0.5cm}Select the ARQ branch;
set $n_{bs}\leftarrow n_{bs}-1$}{\footnotesize\par}

{\footnotesize 05 \hspace{0.5cm}Else}{\footnotesize\par}

{\footnotesize 06 \hspace{0.5cm}\hspace{0.5cm}Sample branch $k$
via Thompson sampling from $\{(a_{k},b_{k})\}_{k=1}^{2}$}{\footnotesize\par}

{\footnotesize 07 \hspace{0.5cm}Endif}{\footnotesize\par}

{\footnotesize 08 \hspace{0.5cm}If $k=\mathrm{IDE}$ then}{\footnotesize\par}

{\footnotesize 09 \hspace{0.5cm}\hspace{0.5cm}Apply the IDE branch
update (Algorithm \ref{alg:sparq-ide})}{\footnotesize\par}

{\footnotesize 10 \hspace{0.5cm}Else}{\footnotesize\par}

{\footnotesize 11 \hspace{0.5cm}\hspace{0.5cm}Refresh the eigen-basis
$B$ if stale (Algorithm \ref{alg:sparq-eigen})}{\footnotesize\par}

{\footnotesize 12 \hspace{0.5cm}\hspace{0.5cm}Apply the ARQ branch
update to $X$ and $A$ (Algorithm \ref{alg:sparq-arq})}{\footnotesize\par}

{\footnotesize 13 \hspace{0.5cm}\hspace{0.5cm}Apply the Levy-flight
quarantine correction using $\alpha_{q},\rho_{q},\beta_{L},\sigma_{L}$
(\ref{alg:sparq-quarantine})}{\footnotesize\par}

{\footnotesize 14 \hspace{0.5cm}Endif}{\footnotesize\par}

{\footnotesize 15 \hspace{0.5cm}Update the bandit posterior $(a_{k},b_{k})$
for branch $k$}{\footnotesize\par}

{\footnotesize 16 \hspace{0.5cm}If $f^{\ast}<f^{\mathrm{prev}}$ then}{\footnotesize\par}

{\footnotesize 17 \hspace{0.5cm}\hspace{0.5cm}Set $f^{\mathrm{prev}}\leftarrow f^{\ast}$,
$s\leftarrow0$}{\footnotesize\par}

{\footnotesize 18 \hspace{0.5cm}Else}{\footnotesize\par}

{\footnotesize 19 \hspace{0.5cm}\hspace{0.5cm}Set $s\leftarrow s+1$}{\footnotesize\par}

{\footnotesize 20 \hspace{0.5cm}Endif}{\footnotesize\par}

{\footnotesize 21 \hspace{0.5cm}Update the polish-stagnation counter
$s_{p}$ from the relative gain in $f^{\ast}$}{\footnotesize\par}

{\footnotesize 22 \hspace{0.5cm}If $s_{p}\geq\tau_{p}$ then apply
the stagnation-gated elite polish (Algorithm \ref{alg:sparq-polish})}{\footnotesize\par}

{\footnotesize 23 \hspace{0.5cm}If $s\geq\tau_{s}$ and the population
spread has collapsed then apply the OBL basin escape (Algorithm \ref{alg:sparq-obl})}{\footnotesize\par}

{\footnotesize 24 \hspace{0.5cm}Apply the hard-stagnation rejuvenation
check (Algorithm \ref{alg:sparq-rejuv})}{\footnotesize\par}

{\footnotesize 25 \hspace{0.5cm}Set $N$ to the current NLPSR target
size and shrink $X$ if needed (Algorithm \ref{alg:sparq-nlpsr})}{\footnotesize\par}

{\footnotesize 26 \hspace{0.5cm}Decay the bandit posteriors $(a_{k},b_{k})$
toward the uniform prior}{\footnotesize\par}

{\footnotesize 27 \hspace{0.5cm}Trim $A$ again to size $\lfloor\alpha_{A}N\rfloor$}{\footnotesize\par}

{\footnotesize 28 \hspace{0.5cm}Update the stopping state and report
the current incumbent}{\footnotesize\par}

{\footnotesize 29 \hspace{0.5cm}Set $e$ equal to the current number
of function evaluations}{\footnotesize\par}

{\footnotesize 30 Endwhile}{\footnotesize\par}

{\footnotesize 31 Return $x^{\ast},\ f^{\ast}$}{\footnotesize\par}

{\footnotesize\caption{Main workflow of SPARQ\label{alg:sparq-main}}
}{\footnotesize\par}
\end{algorithm}

Concretely, each pass through Algorithm 1 maps to the following steps.
The archive is trimmed at the start of the loop (line 2). Branch selection
happens in lines 3--7: a forced ARQ warm-up while the bootstrap counter
$n_{bs}>0$ (lines 3--4), after which the Thompson bandit samples
a branch from $\{(a_{k},b_{k})\}$ (lines 5--6). The chosen branch
is then executed in lines 8--14: the IDE update (Algorithm \ref{alg:sparq-ide})
if $k=\mathrm{IDE}$ (line 9), or, on the ARQ side, an eigen-basis
refresh if stale (Algorithm \ref{alg:sparq-eigen}, line 11), the
ARQ trial generation and RTR selection (Algorithm \ref{alg:sparq-arq},
line 12), and the Lévy-flight quarantine correction (Algorithm (\ref{alg:sparq-quarantine},
line 13). The bandit posterior for the executed branch is updated
immediately afterward (line 15). Two independent stagnation signals
are then maintained: the ordinary counter $s$, reset on any improvement
and incremented otherwise (lines 16--20), and the meaningful-gain
counter $s_{p}$\LyXZeroWidthSpace{} (line 21), which gates the elite
polish (Algorithm \ref{alg:sparq-polish}) once $s_{p}\geq\tau_{p}$\LyXZeroWidthSpace{}
(line 22). The remaining stabilization mechanisms fire in sequence
and are each individually gated: the OBL basin escape (Algorithm \ref{alg:sparq-obl})
once $s\geq\tau_{s}$\LyXZeroWidthSpace{} and the population spread
has collapsed (line 23), and the hard-stagnation rejuvenation check
(Algorithm \ref{alg:sparq-rejuv}), which evaluates its own survival/hard-stagnation
conditions internally (line 24). The population is then shrunk toward
its current NLPSR target size if needed (Algorithm \ref{alg:sparq-nlpsr},
line 25), the bandit posteriors are decayed toward their uninformative
prior (line 26), and the archive is trimmed a second time before the
stopping state is updated and the current incumbent reported (lines
27--28).

\begin{algorithm}[H]
{\footnotesize INPUT}{\footnotesize\par}

{\footnotesize - $X=\{x_{i}\}_{i=1}^{N}$: current population}{\footnotesize\par}

{\footnotesize - $f,\ \Omega$: objective function and feasible domain}{\footnotesize\par}

{\footnotesize - $g,\ g_{\max},\ g_{t},\ \Lambda,\ \lambda_{c}$: IDE
progress index, horizon, phase threshold, patience window, low-success
counter}{\footnotesize\par}

{\footnotesize - $\{F_{i}^{\mathrm{IDE}},CR_{i}^{\mathrm{IDE}}\}_{i=1}^{N}$:
individual IDE parameter memories}{\footnotesize\par}

{\footnotesize OUTPUT}{\footnotesize\par}

{\footnotesize - Updated $X$, incumbent pair $(x^{\ast},f^{\ast})$,
updated $g_{t},\ \lambda_{c}$}{\footnotesize\par}

{\footnotesize IDE branch update}{\footnotesize\par}

{\footnotesize 01 Advance $g$ (synced to global progress or incremented),
clamped to $g_{\max}$}{\footnotesize\par}

{\footnotesize 02 Sort $X$ in nondecreasing objective value}{\footnotesize\par}

{\footnotesize 03 Set $\varepsilon(g)\leftarrow0.1+0.9\cdot10^{5(g/g_{\max}-1)}$}{\footnotesize\par}

{\footnotesize 04 If $g<g_{t}$ then set $S_{th}\leftarrow0$; otherwise
set $S_{th}\leftarrow0.1$}{\footnotesize\par}

{\footnotesize 05 For each $i=1,\ldots,N$ do}{\footnotesize\par}

{\footnotesize 06 \hspace{0.5cm}Sample distinct indices $o,r_{1},r_{2},r_{3}\neq i$}{\footnotesize\par}

{\footnotesize 07 \hspace{0.5cm}Set $m(g)\leftarrow\max\{2,\mathrm{round}(\varepsilon(g)N)\}$}{\footnotesize\par}

{\footnotesize 08 \hspace{0.5cm}If $g\leq g_{t}$ then with probability
$0.9\,\varepsilon(g)$ choose $r_{1}^{\star}$ from the best $m(g)$
individuals; otherwise $r_{1}^{\star}\leftarrow r_{1}$}{\footnotesize\par}

{\footnotesize 09 \hspace{0.5cm}Else with probability $0.5$ choose
$r_{1}^{\star}$ from the best $m(g)$ individuals; otherwise $r_{1}^{\star}\leftarrow r_{1}$}{\footnotesize\par}

{\footnotesize 10 \hspace{0.5cm}If $g>g_{t}$ and a Bernoulli($0.5$)
trial succeeds then}{\footnotesize\par}

{\footnotesize 11 \hspace{0.5cm}\hspace{0.5cm}Set $v_{i}\leftarrow x_{i}+F_{i}^{\mathrm{IDE}}(x_{r_{1}^{\star}}-x_{o})+F_{i}^{\mathrm{IDE}}(x_{r_{2}}-x_{r_{3}})$}{\footnotesize\par}

{\footnotesize 12 \hspace{0.5cm}Else}{\footnotesize\par}

{\footnotesize 13 \hspace{0.5cm}\hspace{0.5cm}Set $v_{i}\leftarrow x_{o}+F_{i}^{\mathrm{IDE}}(x_{r_{1}^{\star}}-x_{o})+F_{i}^{\mathrm{IDE}}(x_{r_{2}}-x_{r_{3}})$}{\footnotesize\par}

{\footnotesize 14 \hspace{0.5cm}Endif}{\footnotesize\par}

{\footnotesize 15 \hspace{0.5cm}Repair $v_{i}$ to $\Omega$}{\footnotesize\par}

{\footnotesize 16 \hspace{0.5cm}Construct $y_{i}$ by binomial crossover
of $x_{i}$ and $v_{i}$ with rate $CR_{i}^{\mathrm{IDE}}$}{\footnotesize\par}

{\footnotesize 17 \hspace{0.5cm}Repair $y_{i}$ to $\Omega$, evaluate
$f(y_{i})$}{\footnotesize\par}

{\footnotesize 18 Endfor}{\footnotesize\par}

{\footnotesize 19 Define the success set $I_{s}=\{i:f(y_{i})<f(x_{i})\}$
and set $SR\leftarrow|I_{s}|/N$}{\footnotesize\par}

{\footnotesize 20 If $g<g_{t}$ then}{\footnotesize\par}

{\footnotesize 21 \hspace{0.5cm}If $SR\leq S_{th}$ then $\lambda_{c}\leftarrow\lambda_{c}+1$;
otherwise $\lambda_{c}\leftarrow0$}{\footnotesize\par}

{\footnotesize 22 \hspace{0.5cm}If $\lambda_{c}\geq\Lambda$ then
$g_{t}\leftarrow g$}{\footnotesize\par}

{\footnotesize 23 Endif}{\footnotesize\par}

{\footnotesize 24 For each $i\in I_{s}$ do record the accepted step
$y_{i}-x_{i}$ (trajectory echo), replace $x_{i}\leftarrow y_{i}$,
and update $(x^{\ast},f^{\ast})$ if improved}{\footnotesize\par}

{\footnotesize 25 Return the updated $X,\ (x^{\ast},f^{\ast}),\ g_{t},\ \lambda_{c}$}{\footnotesize\par}

{\footnotesize\caption{IDE branch update in SPARQ\label{alg:sparq-ide}}
}{\footnotesize\par}
\end{algorithm}

Each individual is updated using a progress-dependent mix of guidance
from an elite pool and random companions (lines 1--17), with the
elite-guidance probability and pool size shrinking as the search matures.
Successful trials replace their parents and feed the trajectory-echo
memory (lines 19, 24), while a separate success-rate check can advance
the branch's internal phase threshold early if progress stalls (lines
20--23).

\begin{algorithm}[H]
{\footnotesize INPUT}{\footnotesize\par}

{\footnotesize - $X$: population, sorted by fitness}{\footnotesize\par}

{\footnotesize - $D,\ N,\ \rho_{\mathrm{eig}},\ D_{\min}$: dimension,
population size, top-fraction, minimum dimension for activation}{\footnotesize\par}

{\footnotesize OUTPUT}{\footnotesize\par}

{\footnotesize - $B$ (eigenvectors), $\mathbf{s}$ (per-axis scales),
validity flag}{\footnotesize\par}

{\footnotesize Eigen-basis refresh}{\footnotesize\par}

{\footnotesize 01 If $D>100$ then mark the basis invalid and return}{\footnotesize\par}

{\footnotesize 02 Set $top\leftarrow\max(D+2,\mathrm{round}(\rho_{\mathrm{eig}}N))$,
clamped to $top\leq N$}{\footnotesize\par}

{\footnotesize 03 If $D<D_{\min}$ or $top<\max(D+2,4)$ then mark
the basis invalid and return}{\footnotesize\par}

{\footnotesize 04 Compute the mean $\bar{x}$ of the top $top$ individuals}{\footnotesize\par}

{\footnotesize 05 Compute the empirical covariance matrix $C$ of the
top $top$ individuals about $\bar{x}$}{\footnotesize\par}

{\footnotesize 06 Regularize the diagonal of $C$ to avoid singularity}{\footnotesize\par}

{\footnotesize 07 Diagonalize $C$ via Jacobi rotation to obtain eigenvectors
$B$ and eigenvalues $w$}{\footnotesize\par}

{\footnotesize 08 Set the per-axis scale $s_{j}\leftarrow\mathrm{clip}(\sqrt{w_{j}/w_{\max}},\,0.05,\,1)$
for each axis $j$}{\footnotesize\par}

{\footnotesize 09 Mark the basis valid and return $B,\ \mathbf{s}$}{\footnotesize\par}

{\footnotesize\caption{Eigen-basis refresh in SPARQ\label{alg:sparq-eigen}}
}{\footnotesize\par}
\end{algorithm}

The covariance matrix of the best half of the population is computed
and diagonalized (lines 4--7) to obtain the natural axes along which
the population is currently spread. These axes, together with per-axis
scale factors, are reused by both the ARQ crossover (Algorithm \ref{alg:sparq-arq})
and the elite polish (Algorithm 6) until the next refresh.

\begin{algorithm}[H]
{\footnotesize INPUT}{\footnotesize\par}

{\footnotesize - $X,\ A,\ N,\ D,\ \Omega$: population, archive, size,
dimension, domain}{\footnotesize\par}

{\footnotesize - $H,\ M_{F},\ M_{CR},\ F_{lo},\ F_{hi}$: SHADE memory
and F bounds}{\footnotesize\par}

{\footnotesize - $p_{\max},\ p_{\min},\ k_{r}^{0},\ k_{r}^{1}$: pbest
and rank-pressure schedules}{\footnotesize\par}

{\footnotesize - $B,\ p_{\mathrm{eig}}$: eigen-basis (Algorithm 3)
and crossover probability}{\footnotesize\par}

{\footnotesize - $n_{\mathrm{pool}}$: restricted-tournament pool size}{\footnotesize\par}

{\footnotesize OUTPUT}{\footnotesize\par}

{\footnotesize - Updated $X,\ A,\ (M_{F},M_{CR}),\ (x^{\ast},f^{\ast})$,
ARQ success statistics}{\footnotesize\par}

{\footnotesize ARQ branch update}{\footnotesize\par}

{\footnotesize 01 Sort $X$ by fitness to obtain the rank order $\mathrm{ord}$}{\footnotesize\par}

{\footnotesize 02 If the eigen-basis is stale then refresh it (Algorithm
3)}{\footnotesize\par}

{\footnotesize 03 Set $m$ parents to try (dimension-dependent agent
fraction) and draw them at random from $\mathrm{ord}$}{\footnotesize\par}

{\footnotesize 04 For each drawn parent $t=1,\ldots,m$ sample $(F_{t},CR_{t})$
from a random slot of $(M_{F},M_{CR})$}{\footnotesize\par}

{\footnotesize 05 If CR-sorting is enabled then reassign the sampled
$CR$ values so the best-ranked parents receive the smallest ones}{\footnotesize\par}

{\footnotesize 06 For each selected parent $i$ do}{\footnotesize\par}

{\footnotesize 07 \hspace{0.5cm}Draw $x_{pbest}$ uniformly from the
top $\lceil p(t)N\rceil$ individuals, $p(t)$ following the $p_{\max}\!\to\!p_{\min}$
schedule}{\footnotesize\par}

{\footnotesize 08 \hspace{0.5cm}Pick $r_{1}$ by rank-biased selection
with pressure $k_{r}(t)$, forbidding $i$}{\footnotesize\par}

{\footnotesize 09 \hspace{0.5cm}With probability $0.5$ draw $r_{2}$
from the archive $A$; otherwise pick $r_{2}$ by rank-biased selection,
forbidding $i,r_{1}$}{\footnotesize\par}

{\footnotesize 10 \hspace{0.5cm}Compute the jSO weight $K(F_{t})$
and form the mutant $v\leftarrow x_{i}+K(F_{t})(x_{pbest}-x_{i})+F_{t}(x_{r_{1}}-x_{r_{2}})$}{\footnotesize\par}

{\footnotesize 11 \hspace{0.5cm}Repair $v$ to $\Omega$}{\footnotesize\par}

{\footnotesize 12 \hspace{0.5cm}With probability $p_{\mathrm{eig}}$
apply eigen-coordinate binomial crossover using $B$; otherwise apply
classical binomial crossover with rate $CR_{t}$ to obtain the trial
$u$}{\footnotesize\par}

{\footnotesize 13 \hspace{0.5cm}Repair $u$ to $\Omega$ and evaluate
$f(u)$}{\footnotesize\par}

{\footnotesize 14 \hspace{0.5cm}If $f(u)<f(x_{i})$ then}{\footnotesize\par}

{\footnotesize 15 \hspace{0.5cm}\hspace{0.5cm}Archive $x_{i}$, set
$x_{i}\leftarrow u$, and record the success $(F_{t},CR_{t},\mathrm{gain})$}{\footnotesize\par}

{\footnotesize 16 \hspace{0.5cm}Else}{\footnotesize\par}

{\footnotesize 17 \hspace{0.5cm}\hspace{0.5cm}Draw a pool $Q\subset X$
of size $n_{\mathrm{pool}}$ (excluding $i$) and find the nearest
$q^{\star}\in Q$ to $u$}{\footnotesize\par}

{\footnotesize 18 \hspace{0.5cm}\hspace{0.5cm}If $f(u)<f(q^{\star})$
then archive $q^{\star}$, set $q^{\star}\leftarrow u$, inherit its
IDE parameters from $i$, and record the success}{\footnotesize\par}

{\footnotesize 19 \hspace{0.5cm}Endif}{\footnotesize\par}

{\footnotesize 20 \hspace{0.5cm}Update $(x^{\ast},f^{\ast})$ if improved}{\footnotesize\par}

{\footnotesize 21 Endfor}{\footnotesize\par}

{\footnotesize 22 Update the memory slot of $(M_{F},M_{CR})$ from
the recorded successes}{\footnotesize\par}

{\footnotesize 23 Return the updated $X,\ A,\ (M_{F},M_{CR}),\ (x^{\ast},f^{\ast})$}{\footnotesize\par}

{\footnotesize\caption{ARQ branch trial generation and RTR selection in SPARQ\label{alg:sparq-arq}}
}{\footnotesize\par}
\end{algorithm}

For each selected parent, a mutant is formed by pulling it toward
a rank-biased elite exemplar plus a rank-biased difference term (lines
7--10), then crossed over either in the eigen-basis or in the raw
coordinates (line 11). The trial first competes with its own parent
and, failing that, with its nearest neighbor in a small random pool
(lines 13--19), with all successes feeding the SHADE memory update
(line 22).

\begin{algorithm}[H]
{\footnotesize INPUT}{\footnotesize\par}

{\footnotesize - $X,\ N,\ D,\ \Omega$: population, size, dimension,
domain}{\footnotesize\par}

{\footnotesize - $\alpha_{q},\ \rho_{q},\ \beta_{L},\ \sigma_{L}$:
IQR coefficient, repair fraction, Levy exponent, step scale}{\footnotesize\par}

{\footnotesize OUTPUT}{\footnotesize\par}

{\footnotesize - Updated $X,\ A,\ (x^{\ast},f^{\ast})$}{\footnotesize\par}

{\footnotesize Levy-flight quarantine}{\footnotesize\par}

{\footnotesize 01 Compute $Q_{1},Q_{3}$ and $IQR=Q_{3}-Q_{1}$ of
the fitness values}{\footnotesize\par}

{\footnotesize 02 Set the threshold $\theta\leftarrow Q_{3}+\alpha_{q}\cdot IQR$}{\footnotesize\par}

{\footnotesize 03 Compute the centroid $c$ of the best half of $X$
by fitness}{\footnotesize\par}

{\footnotesize 04 Set the outlier set $O\leftarrow\{i:f(x_{i})\geq\theta\}$;
if $O=\varnothing$ then return}{\footnotesize\par}

{\footnotesize 05 Draw a random subset $O_{\rho}\subset O$ with $|O_{\rho}|=\lfloor\rho_{q}|O|\rfloor$}{\footnotesize\par}

{\footnotesize 06 For each $i\in O_{\rho}$ do}{\footnotesize\par}

{\footnotesize 07 \hspace{0.5cm}For each coordinate $j$ draw a Levy
step $\eta_{j}$ via the Mantegna algorithm with exponent $\beta_{L}$}{\footnotesize\par}

{\footnotesize 08 \hspace{0.5cm}Set $x_{i}'\leftarrow c+\sigma_{L}\cdot(\mathbf{u}-\mathbf{\ell})\odot\boldsymbol{\eta}$}{\footnotesize\par}

{\footnotesize 09 \hspace{0.5cm}Repair $x_{i}'$ to $\Omega$ and
evaluate $f(x_{i}')$}{\footnotesize\par}

{\footnotesize 10 \hspace{0.5cm}If $f(x_{i}')<f(x_{i})$ then archive
$x_{i}$, set $x_{i}\leftarrow x_{i}'$, and re-seed its IDE parameters}{\footnotesize\par}

{\footnotesize 11 \hspace{0.5cm}Update $(x^{\ast},f^{\ast})$ if improved}{\footnotesize\par}

{\footnotesize 12 Endfor}{\footnotesize\par}

{\footnotesize 13 Return the updated $X,\ A,\ (x^{\ast},f^{\ast})$}{\footnotesize\par}

{\footnotesize\caption{Levy-flight quarantine in SPARQ\label{alg:sparq-quarantine}}
}{\footnotesize\par}
\end{algorithm}

Individuals whose fitness exceeds a robust outlier threshold (line
2) are relocated toward the population's best-half centroid using
a heavy-tailed Lévy step rather than an ordinary Gaussian one (lines
7--8), giving them an occasional large jump to a genuinely different
region. A relocation is only kept if it actually improves on the original
individual (line 9).

\begin{algorithm}[H]
{\footnotesize INPUT}{\footnotesize\par}

{\footnotesize - $x^{\ast},\ f^{\ast},\ X,\ N,\ D,\ \Omega$: incumbent,
population, size, dimension, domain}{\footnotesize\par}

{\footnotesize - $B,\ \mathbf{s}$: eigen-basis and per-axis scales
(Algorithm 3)}{\footnotesize\par}

{\footnotesize - $\rho_{p},\ B_{p}$: burst-size fraction and evaluation-budget
cap}{\footnotesize\par}

{\footnotesize - $\sigma_{p},\ \sigma_{c},\ b_{sw},\ \rho_{sw},\ \sigma_{e}$:
eigen-step, coordinate-step, Solis-Wets bias/step, echo scale}{\footnotesize\par}

{\footnotesize OUTPUT}{\footnotesize\par}

{\footnotesize - Updated $x^{\ast},\ f^{\ast},\ X$, adapted step sizes
and cooldown state}{\footnotesize\par}

{\footnotesize Stagnation-gated elite polish}{\footnotesize\par}

{\footnotesize 01 Set $f_{\mathrm{before}}\leftarrow f^{\ast}$ and
compute the evolutionary gain-per-evaluation since the last activation}{\footnotesize\par}

{\footnotesize 02 Set the burst size $k\leftarrow\max(6,\mathrm{round}(\rho_{p}N))$,
enlarged late in the run to cover $D$ coordinates}{\footnotesize\par}

{\footnotesize 03 Set $\mathrm{fails}\leftarrow0$, $\mathrm{wins}\leftarrow0$}{\footnotesize\par}

{\footnotesize 04 For $t=1,\ldots,k$ do}{\footnotesize\par}

{\footnotesize 05 \hspace{0.5cm}If the evaluation budget or the cap
$B_{p}$ is exceeded then break}{\footnotesize\par}

{\footnotesize 06 \hspace{0.5cm}Draw a probe mode $m\sim U(0,1)$}{\footnotesize\par}

{\footnotesize 07 \hspace{0.5cm}If $m<0.30$ then probe a single round-robin
coordinate (opposition or adaptive Gaussian step)}{\footnotesize\par}

{\footnotesize 08 \hspace{0.5cm}Else if $m<0.65$ and $B$ is available
then draw an anisotropic Gaussian step in the eigen-basis, scaled
by $\mathbf{s}$}{\footnotesize\par}

{\footnotesize 09 \hspace{0.5cm}Else if $m<0.80$ and the echo buffer
is non-empty then draw a random linear combination of the buffered
directions (trajectory echo)}{\footnotesize\par}

{\footnotesize 10 \hspace{0.5cm}Else propose $x^{\ast}+b_{sw}+\delta$,
and if it fails, the reflection $x^{\ast}-b_{sw}-\delta$ (Solis-Wets)}{\footnotesize\par}

{\footnotesize 11 \hspace{0.5cm}Repair the candidate to $\Omega$
and evaluate its fitness}{\footnotesize\par}

{\footnotesize 12 \hspace{0.5cm}If the candidate improves $f^{\ast}$
then}{\footnotesize\par}

{\footnotesize 13 \hspace{0.5cm}\hspace{0.5cm}Update $(x^{\ast},f^{\ast})$,
grow the step size of the active probe mode, replace the population's
worst individual, reset $\mathrm{fails}$, increment $\mathrm{wins}$}{\footnotesize\par}

{\footnotesize 14 \hspace{0.5cm}Else shrink the step size of the active
probe mode and increment $\mathrm{fails}$}{\footnotesize\par}

{\footnotesize 15 \hspace{0.5cm}If $\mathrm{fails}\geq10$ then break}{\footnotesize\par}

{\footnotesize 16 Endfor}{\footnotesize\par}

{\footnotesize 17 If $\mathrm{wins}=0$ then double the cooldown backoff;
otherwise reset it}{\footnotesize\par}

{\footnotesize 18 If the polish's gain-per-evaluation is below the
evolutionary loop's then extend the cooldown}{\footnotesize\par}

{\footnotesize 19 If the relative gain over the burst is negligible
for several consecutive activations then extend the cooldown further}{\footnotesize\par}

{\footnotesize 20 Return the updated $x^{\ast},\ f^{\ast},\ X$}{\footnotesize\par}

{\footnotesize\caption{Stagnation-gated elite polish in SPARQ\label{alg:sparq-polish}}
}{\footnotesize\par}
\end{algorithm}

While the incumbent best is stagnating, a short burst of local probes
is applied directly to it, drawn from four alternative modes, coordinate-wise,
eigen-aligned, trajectory-echo, and Solis-Wets (lines 6--10), and
any improvement is written back into the population's worst slot (line
12). Three independent safeguards (lines 16--18) shut the mechanism
down whenever it proves less productive than ordinary evolution.

\begin{algorithm}[H]
{\footnotesize INPUT}{\footnotesize\par}

{\footnotesize - $X,\ x^{\ast},\ N,\ D,\ \Omega$: population, incumbent,
size, dimension, domain}{\footnotesize\par}

{\footnotesize - $\tau_{s},\ \varepsilon_{v},\ \rho_{b},\ \kappa_{b},\ s$:
stagnation trigger, spread-collapse ratio, escape fraction, cooldown,
stagnation counter}{\footnotesize\par}

{\footnotesize OUTPUT}{\footnotesize\par}

{\footnotesize - Updated $X,\ (x^{\ast},f^{\ast}),\ s$, cooldown state}{\footnotesize\par}

{\footnotesize OBL basin escape}{\footnotesize\par}

{\footnotesize 01 If the cooldown counter is positive then decrement
it and return}{\footnotesize\par}

{\footnotesize 02 If not ($s\geq\tau_{s}$ and the normalized population
spread $<\varepsilon_{v}$) then return}{\footnotesize\par}

{\footnotesize 03 Sort the population by fitness}{\footnotesize\par}

{\footnotesize 04 Set $count\leftarrow\max(1,\lfloor\rho_{b}N\rfloor)$}{\footnotesize\par}

{\footnotesize 05 For each of the $count$ worst individuals $x_{i}$
do}{\footnotesize\par}

{\footnotesize 06 \hspace{0.5cm}For each coordinate $j$ compute the
opposite point $\ell_{j}+u_{j}-x_{i,j}$}{\footnotesize\par}

{\footnotesize 07 \hspace{0.5cm}With probability $0.5$ use the pure
opposition; otherwise interpolate between the opposition and $x^{\ast}$}{\footnotesize\par}

{\footnotesize 08 \hspace{0.5cm}Repair the candidate to $\Omega$
and evaluate its fitness}{\footnotesize\par}

{\footnotesize 09 \hspace{0.5cm}If it improves $x_{i}$ then archive
$x_{i}$, replace it, and re-seed its IDE parameters}{\footnotesize\par}

{\footnotesize 10 \hspace{0.5cm}Update $(x^{\ast},f^{\ast})$ if improved}{\footnotesize\par}

{\footnotesize 11 Endfor}{\footnotesize\par}

{\footnotesize 12 If any replacement was applied then reset the cooldown
to $\kappa_{b}$ and set $s\leftarrow0$}{\footnotesize\par}

{\footnotesize 13 Return the updated $X,\ (x^{\ast},f^{\ast}),\ s$}{\footnotesize\par}

{\footnotesize\caption{On-demand OBL basin escape in SPARQ\label{alg:sparq-obl}}
}{\footnotesize\par}
\end{algorithm}

This only fires when the search has stopped improving and the population
has visibly collapsed onto one region (line 2). A fraction of the
worst individuals are then reflected to the opposite side of the search
domain (lines 6--7), and the stagnation counter is reset if this
yields any actual improvement (line 10).

\begin{algorithm}[H]
{\footnotesize INPUT}{\footnotesize\par}

{\footnotesize - $X,\ x^{\ast},\ N,\ D,\ \Omega$: population, incumbent,
size, dimension, domain}{\footnotesize\par}

{\footnotesize - $\tau_{s},\ m_{r},\ \kappa_{r},\ C_{r},\ s$: stagnation
trigger, rejuvenation factor, keep-fraction, cooldown, stagnation
counter}{\footnotesize\par}

{\footnotesize OUTPUT}{\footnotesize\par}

{\footnotesize - Updated $X,\ (x^{\ast},f^{\ast})$, cooldown and escalation
state}{\footnotesize\par}

{\footnotesize Hard-stagnation rejuvenation}{\footnotesize\par}

{\footnotesize 01 If the cooldown counter is positive then decrement
it and return}{\footnotesize\par}

{\footnotesize 02 Compute the relative fitness spread $f_{\mathrm{rel}}$
across the population}{\footnotesize\par}

{\footnotesize 03 Set $\mathrm{collapsed}\leftarrow(f_{\mathrm{rel}}<\varepsilon_{1})\ \text{or}\ (\mathrm{spread}<\varepsilon_{2})$}{\footnotesize\par}

{\footnotesize 04 Set $\mathrm{hard\_stag}\leftarrow(s\geq m_{r}\tau_{s})$
and (progress within its cutoff) and (spread bounded)}{\footnotesize\par}

{\footnotesize 05 If $\mathrm{collapsed}$ and not $\mathrm{hard\_stag}$
then update the weak-reseed streak; escalate to $\mathrm{hard\_stag}$
once the streak exceeds its limit}{\footnotesize\par}

{\footnotesize 06 Else reset the weak-reseed streak}{\footnotesize\par}

{\footnotesize 07 Set $\mathrm{survival}\leftarrow\mathrm{collapsed}$
and not $\mathrm{hard\_stag}$; if neither holds then return}{\footnotesize\par}

{\footnotesize 08 If $\mathrm{hard\_stag}$ then update the strong-reseed
failure streak; escalate to a near-full reset once it exceeds its
limit}{\footnotesize\par}

{\footnotesize 09 Sort the population by fitness}{\footnotesize\par}

{\footnotesize 10 Set the elite fraction $\mathrm{keep}$ to preserve:
a tiny sliver under a near-full reset, $\kappa_{r}$ under $\mathrm{hard\_stag}$,
or nearly all of $N$ under $\mathrm{survival}$}{\footnotesize\par}

{\footnotesize 11 For each individual ranked below $\mathrm{keep}$
do}{\footnotesize\par}

{\footnotesize 12 \hspace{0.5cm}Draw a uniformly random point in $\Omega$;
archive the old individual and replace it; re-seed its IDE parameters}{\footnotesize\par}

{\footnotesize 13 \hspace{0.5cm}Update $(x^{\ast},f^{\ast})$ if improved}{\footnotesize\par}

{\footnotesize 14 Endfor}{\footnotesize\par}

{\footnotesize 15 If $\mathrm{hard\_stag}$ then reset the stagnation
counter $s$ and set a long cooldown; otherwise set a short cooldown}{\footnotesize\par}

{\footnotesize 16 Return the updated $X,\ (x^{\ast},f^{\ast})$}{\footnotesize\par}

{\footnotesize\caption{Hard-stagnation rejuvenation in SPARQ\label{alg:sparq-rejuv}}
}{\footnotesize\par}
\end{algorithm}

A loose \textquotedbl survival\textquotedbl{} condition keeps re-seeding
a small sliver of the population indefinitely to avoid the appearance
of full convergence, while a much stricter \textquotedbl hard-stagnation\textquotedbl{}
condition, reached only after prolonged non-improvement, re-seeds
the large majority of the population (lines 2--10). Two escalation
counters (lines 5, 8) force a stronger response whenever repeated
firings of either condition keep failing to produce real gains.

\begin{algorithm}[H]
{\footnotesize INPUT}{\footnotesize\par}

{\footnotesize - $X,\ N,\ N_{\mathrm{init}},\ N_{\min},\ \alpha_{\mathrm{NL}}$:
population, current size, size bounds, shrink-shape exponent}{\footnotesize\par}

{\footnotesize - $p$: current progress ratio}{\footnotesize\par}

{\footnotesize OUTPUT}{\footnotesize\par}

{\footnotesize - Updated $X$ and $N$}{\footnotesize\par}

{\footnotesize NLPSR population-size reduction}{\footnotesize\par}

{\footnotesize 01 Compute the shrink-shape exponent $q\leftarrow1-(1-\alpha_{\mathrm{NL}})p$}{\footnotesize\par}

{\footnotesize 02 Compute the target size $N_{\mathrm{target}}\leftarrow\mathrm{round}(N_{\mathrm{init}}+(N_{\min}-N_{\mathrm{init}})p^{q})$,
clamped to $[N_{\min},N_{\mathrm{init}}]$}{\footnotesize\par}

{\footnotesize 03 If $N_{\mathrm{target}}\geq N$ then return unchanged}{\footnotesize\par}

{\footnotesize 04 Sort the population ascending by fitness}{\footnotesize\par}

{\footnotesize 05 Keep the best $N_{\mathrm{target}}$ individuals
(and their IDE memories); discard the rest}{\footnotesize\par}

{\footnotesize 06 Invalidate the eigen-basis and set $N\leftarrow N_{\mathrm{target}}$}{\footnotesize\par}

{\footnotesize 07 Return the updated $X,\ N$}{\footnotesize\par}

{\footnotesize\caption{NLPSR population-size reduction in SPARQ\label{alg:sparq-nlpsr}}
}{\footnotesize\par}
\end{algorithm}

The overall control flow of SPARQ, summarizing the sequence of decisions
executed within each iteration of the main loop, is illustrated in
Figure \ref{fig:floatChart}. The diagram highlights the probabilistic
selection between the ARQ and IDE branches, the mandatory Lévy-flight
quarantine attached to the ARQ branch, and the three graduated stagnation-driven
mechanisms, namely the stagnation-gated elite polish, the opposition-based
basin escape, and the hard-stagnation rejuvenation, together with
their respective activation gates.

\begin{figure}[H]
\begin{centering}
\includegraphics[scale=0.6]{flowchart}
\par\end{centering}
\caption{Overall control flow of the SPARQ algorithm, illustrating branch selection,
the ARQ-attached quarantine mechanism, and the graduated stagnation-driven
activation gates\label{fig:floatChart}}

\end{figure}

The algorithm Figure \ref{fig:floatChart} operates within a loop
governed exclusively by the evaluation budget ($e<T_{eval}$). At
the start of every pass through this loop, the external archive is
first trimmed to size $\lfloor\alpha_{A}\cdot N\rfloor$, preventing
it from growing unchecked as iterations accumulate.

The first substantive decision the system makes concerns which of
the two search branches, ARQ or IDE, will be executed. As long as
the bootstrap counter remains positive, the ARQ branch is enforced
outright, with no genuine choice involved, guaranteeing a brief warm-up
period before the bandit's statistics become reliable. Once this period
elapses, branch selection passes to a Thompson-sampling bandit, which
draws a random sample from each branch's posterior distribution and
favors whichever branch returns the larger value. This selection is
not a hard on/off switch but a probabilistic one: a branch that performs
consistently worse will be chosen less often, yet it is never permanently
excluded, since both posteriors decay toward the uninformative prior
at every iteration, allowing a branch that was previously effective
to be revisited later.

Depending on the outcome of this choice, execution diverges considerably.
If the IDE branch is selected, only its own update is executed, with
no accompanying mechanism. If the ARQ branch is selected, the eigen-basis
is first checked for staleness and refreshed if necessary, the ARQ
update itself is then executed, and the Lévy-flight quarantine correction
is applied immediately afterward as a mandatory step. This quarantine
mechanism is structurally attached exclusively to the ARQ branch and
is never triggered when IDE executes, regardless of how often the
latter is chosen. Both paths, however, converge on the same subsequent
step: updating the bandit's posterior for whichever branch was just
executed.

Two independent stagnation counters are then updated, each feeding
a different downstream mechanism. The first records any absence of
improvement whatsoever, however small, and feeds the basin-escape
and rejuvenation mechanisms. The second records only the absence of
meaningful, relative improvement, and feeds the elite-polish mechanism
exclusively. This separation is deliberate: in large populations or
at high dimensionality, some individual improves the incumbent by
a negligible amount in nearly every generation, so a shared counter
would almost never allow the polish phase to activate.

Three activation gates follow, always evaluated in the same order
within a given iteration, and not mutually exclusive. The first concerns
the elite polish, which is triggered once the second counter exceeds
its stagnation threshold; once activated, it runs a short burst of
local probes, alternating among single-coordinate probing, anisotropic
eigen-aligned steps, replay from the trajectory-echo memory, and the
classical Solis-Wets adaptive search, all protected by three independent
safeguards: a hard evaluation-budget cap, an exponential cooldown
whenever an entire burst fails to yield any improvement, and a gain-per-evaluation
comparison against the ordinary evolutionary loop, which forces the
polish to stand down whenever the core DE search proves the better
use of the same evaluations. The second gate concerns the opposition-based
basin-escape mechanism, which fires only when prolonged non-improvement
and a visible collapse in population spread occur simultaneously,
and only once its own cooldown has elapsed; if it produces no improvement,
the cooldown is renewed without resetting the stagnation counter.
The third gate, hard-stagnation rejuvenation, is not a simple binary
decision but a hierarchical one: it internally evaluates whether only
the milder survival condition holds, re-seeding a small fraction of
the population at random without materially disrupting progress, or
whether the stricter hard-stagnation condition holds, re-seeding the
large majority of the population while preserving only its elite members.
It further maintains its own escalation counters, so that repeated
failures of the mild reseeding force an earlier transition to the
harder condition, and repeated failures of the harder condition escalate
to a near-complete reset.

Two further mechanisms execute at every iteration with no gating condition
at all: population-size reduction via NLPSR, and the decay of the
bandit's posteriors toward the uniform prior. NLPSR is always invoked,
but effectively does nothing whenever the current target size is not
smaller than the current population, so it remains continuously active
without always being consequential. The iteration closes with a second
trimming of the archive, since the population size may have changed,
and with an update to the current incumbent best, before the loop
returns to its head until the evaluation budget is exhausted and the
final solution is returned.

The overall picture that emerges is one of two permanently available
search branches, whose relative usage is regulated probabilistically
rather than through hard activation or deactivation, surrounded by
an increasingly rare and increasingly disruptive layer of rescue mechanisms,
each gated by strictly graduated stagnation conditions and each equipped
with its own safeguards that switch it off automatically once it proves
ineffective.

\section{Experimental results\label{sec:ExperimentalResults}}

\subsection{Experimental Configuration\label{subsec:generalsettings}}

All methods evaluated in this study, the proposed SPARQ algorithm
together with the thirteen competing optimizers, were implemented
in ANSI C++ within the open-source OptimSolution \citep{optimsolution}
framework, developed by Vasileios Charilogis and publicly available
at \url{https://github.com/charilog/optimsolution} (accessed 18 July
2026). All experiments were compiled and executed under Debian 12.12
with GCC 13.4, on a workstation equipped with an AMD Ryzen 9 9950X3D
processor (16 cores, 32 threads) and 64 GB of DDR4 memory. To account
for the stochastic nature of the algorithms, each benchmark problem
was solved over 30 independent runs, every run initialized with a
distinct random seed. Across all runs and all methods, the evaluation
budget was fixed at 150,000 function evaluations, consistent with
the termination criterion adopted by the CEC 2011 real-world optimization
benchmark \citep{de_CEC2011} and allowing direct comparison with
prior work. The best and mean objective values obtained at termination,
averaged over the 30 independent runs per problem, are reported as
the primary performance measures throughout the results section. Table
\ref{tab:settings} summarizes the SPARQ parameter settings actually
used in these experiments, expressed in the notation introduced in
Algorithm $\ref{alg:sparq-main}$ and Sections $\ref{subsec:NLPSR}$--$\ref{subsec:HardStagnationRejuvenation}$,
while Table \ref{tab:constants} lists the internal constants that
remain fixed at compile time and are therefore not exposed through
the configuration interface.
\begin{flushleft}
{\scriptsize{}%
\begin{longtable}[c]{lll}
\caption{SPARQ parameter settings used in the experiments, following the notation
introduced in Algorithm \ref{alg:sparq-main} and Sections \ref{subsec:NLPSR}--\ref{subsec:HardStagnationRejuvenation}\label{tab:settings}}
\tabularnewline
\toprule 
{\scriptsize\textbf{Symbol}} & {\scriptsize\textbf{Value}} & {\scriptsize\textbf{Description}}\tabularnewline
\midrule
\endhead
\midrule 
\multicolumn{3}{l}{{\scriptsize General configuration}}\tabularnewline
{\scriptsize$N_{fe}$} & {\scriptsize 150,000} & {\scriptsize Evaluation budget per run}\tabularnewline
{\scriptsize Runs} & {\scriptsize 30} & {\scriptsize Independent repeated runs}\tabularnewline
{\scriptsize Initialization} & {\scriptsize uniform} & {\scriptsize Population initialization distribution}\tabularnewline
\multicolumn{3}{l}{{\scriptsize\textit{Population / NLPSR (Section~\ref{subsec:NLPSR})}}}\tabularnewline
{\scriptsize$N_{init}$} & {\scriptsize 100} & {\scriptsize Initial population size (a good choice would be $N_{init}=18D$)}\tabularnewline
{\scriptsize$N_{min}$} & {\scriptsize 4} & {\scriptsize Minimum population size}\tabularnewline
{\scriptsize$\alpha_{NL}$} & {\scriptsize 0.5} & {\scriptsize NLPSR shrink-shape exponent}\tabularnewline
\addlinespace
\multicolumn{3}{l}{{\scriptsize\textit{SHADE memory for $F$/$CR$ (Section~\ref{subsec:3.2. SHADEStyle})}}}\tabularnewline
{\scriptsize$H$} & {\scriptsize 6} & {\scriptsize SHADE circular memory size}\tabularnewline
{\scriptsize$M_{F}^{H}$} & {\scriptsize 0.9} & {\scriptsize Terminal memory slot for $F$}\tabularnewline
{\scriptsize$M_{CR}^{H}$} & {\scriptsize 0.9} & {\scriptsize Terminal memory slot for $CR$}\tabularnewline
{\scriptsize$F_{lo}$} & {\scriptsize 0.05} & {\scriptsize Lower bound for $F$ sampling}\tabularnewline
{\scriptsize$F_{hi}$} & {\scriptsize 1.40} & {\scriptsize Upper bound for $F$ sampling}\tabularnewline
\addlinespace
\multicolumn{3}{l}{{\scriptsize\textit{jSO schedule / rank-based pressure (Section~\ref{subsec:RankBased})}}}\tabularnewline
{\scriptsize$p_{max}$} & {\scriptsize 0.25} & {\scriptsize Maximum $p_{best}$ fraction}\tabularnewline
{\scriptsize$p_{min}$} & {\scriptsize 0.10} & {\scriptsize Minimum $p_{best}$ fraction}\tabularnewline
{\scriptsize$k_{r,init}$\textsuperscript{‡}} & {\scriptsize 2.0} & {\scriptsize Initial rank-selection pressure}\tabularnewline
{\scriptsize$k_{r,final}$\textsuperscript{‡}} & {\scriptsize 3.0} & {\scriptsize Final rank-selection pressure}\tabularnewline
\addlinespace
\multicolumn{3}{l}{{\scriptsize\textit{Eigen-coordinate crossover (Section~\ref{subsec:EigenCrossover},
Algorithm~\ref{alg:sparq-eigen})}}}\tabularnewline
{\scriptsize$p_{eig}$} & {\scriptsize 0.40} & {\scriptsize Eigen-crossover probability}\tabularnewline
{\scriptsize$T_{eig}$} & {\scriptsize 10} & {\scriptsize Basis refresh period (iterations)}\tabularnewline
{\scriptsize$\rho_{eig}$} & {\scriptsize 0.50} & {\scriptsize Top-fraction used for the basis}\tabularnewline
{\scriptsize$D_{min}$} & {\scriptsize 2} & {\scriptsize Minimum dimension for eigen-basis activation}\tabularnewline
\addlinespace
\multicolumn{3}{l}{{\scriptsize\textit{RTR selection / archive (Algorithm~\ref{alg:sparq-arq})}}}\tabularnewline
{\scriptsize$n_{pool}$} & {\scriptsize 14} & {\scriptsize RTR neighbor pool size}\tabularnewline
{\scriptsize --} & {\scriptsize 0.10} & {\scriptsize Dynamic pool fraction of $N$ (not named in the article)}\tabularnewline
{\scriptsize$\alpha_{A}$} & {\scriptsize 1.5} & {\scriptsize Archive size factor}\tabularnewline
\addlinespace
\multicolumn{3}{l}{{\scriptsize\textit{Thompson-sampling bandit (Section~\ref{subsec:Thompson})}}}\tabularnewline
{\scriptsize$\gamma_{b}$} & {\scriptsize 0.97} & {\scriptsize Bandit posterior decay rate}\tabularnewline
{\scriptsize$n_{bs}$} & {\scriptsize 2} & {\scriptsize Forced ARQ bootstrap iterations}\tabularnewline
\addlinespace
\multicolumn{3}{l}{{\scriptsize\textit{Lévy-flight quarantine (Section~\ref{subsec:L=0000E9vyFlight},
Algorithm~\ref{alg:sparq-quarantine})}}}\tabularnewline
{\scriptsize$\alpha_{q}$} & {\scriptsize 1.0} & {\scriptsize IQR outlier threshold coefficient}\tabularnewline
{\scriptsize$\rho_{q}$} & {\scriptsize 0.08} & {\scriptsize Fraction of outliers repaired}\tabularnewline
{\scriptsize$\beta_{L}$} & {\scriptsize 1.5} & {\scriptsize Lévy distribution exponent}\tabularnewline
{\scriptsize$\sigma_{L}$} & {\scriptsize 0.10} & {\scriptsize Lévy quarantine step scale}\tabularnewline
\addlinespace
\multicolumn{3}{l}{{\scriptsize\textit{OBL basin escape (Section~\ref{subsec:ODOBB},
Algorithm~\ref{alg:sparq-obl})}}}\tabularnewline
{\scriptsize$\tau_{s}$} & {\scriptsize 30} & {\scriptsize Stagnation threshold}\tabularnewline
{\scriptsize$\varepsilon_{v}$} & {\scriptsize$10^{-3}$} & {\scriptsize Population-spread collapse ratio}\tabularnewline
{\scriptsize$\rho_{b}$} & {\scriptsize 0.30} & {\scriptsize Fraction of individuals relocated}\tabularnewline
{\scriptsize$\kappa_{b}$} & {\scriptsize 80} & {\scriptsize Cooldown after escape}\tabularnewline
\addlinespace
\multicolumn{3}{l}{{\scriptsize\textit{Agent fraction (Algorithm~\ref{alg:sparq-arq},
line~3)}}}\tabularnewline
{\scriptsize --} & {\scriptsize 0.5} & {\scriptsize Fraction of parents tried per generation ($m$ in Alg.~4;
no dedicated symbol)}\tabularnewline
{\scriptsize --} & {\scriptsize 100} & {\scriptsize Dimension above which this fraction applies fully (not
named in the article)}\tabularnewline
\addlinespace
\multicolumn{3}{l}{{\scriptsize\textit{Stagnation-gated elite polish (Section~\ref{subsec:StagnationPolish},
Algorithm~\ref{alg:sparq-polish})}}}\tabularnewline
{\scriptsize$\tau_{p}$} & {\scriptsize 10} & {\scriptsize Iterations before a polish burst fires}\tabularnewline
{\scriptsize$\rho_{p}$} & {\scriptsize 0.10} & {\scriptsize Polish burst size (probes)}\tabularnewline
{\scriptsize$B_{p}$} & {\scriptsize 0.12} & {\scriptsize Maximum evaluation share for polish}\tabularnewline
{\scriptsize$\pi_{p}$} & {\scriptsize 0.80} & {\scriptsize Late-phase threshold (trigger tightens to 1)}\tabularnewline
{\scriptsize --} & {\scriptsize 0.85} & {\scriptsize Progress fraction after which the burst enlarges (not
separately named; grouped under $\pi_{p}$ in Algorithm~1)}\tabularnewline
\addlinespace
\multicolumn{3}{l}{{\scriptsize\textit{Hard-stagnation rejuvenation (Section~\ref{subsec:HardStagnationRejuvenation},
Algorithm~\ref{alg:sparq-obl})}}}\tabularnewline
{\scriptsize$m_{r}$} & {\scriptsize 4} & {\scriptsize Hard-stagnation threshold multiplier}\tabularnewline
{\scriptsize$\kappa_{r}$} & {\scriptsize 0.25} & {\scriptsize Elite fraction preserved}\tabularnewline
{\scriptsize$C_{r}$} & {\scriptsize 150} & {\scriptsize Cooldown after rejuvenation}\tabularnewline
{\scriptsize --} & {\scriptsize 8} & {\scriptsize Failed weak-reseed streak limit (not named in the article)}\tabularnewline
{\scriptsize --} & {\scriptsize 2} & {\scriptsize Failed hard-reseed streak limit (not named in the article)}\tabularnewline
{\scriptsize --} & {\scriptsize 0.90} & {\scriptsize Progress fraction beyond which disabled (not named in
the article)}\tabularnewline
\addlinespace
\multicolumn{3}{l}{{\scriptsize\textit{CR-sorting (Section~3.3)}}}\tabularnewline
{\scriptsize --} & {\scriptsize enabled} & {\scriptsize Rank-based CR reassignment}\tabularnewline
\end{longtable}}{\scriptsize\par}
\par\end{flushleft}

\begin{table}[H]
\centering{}{\footnotesize\caption{Internal constants of SPARQ fixed at compile time \label{tab:constants}}
}{\footnotesize{}%
\begin{tabular}{lll}
\toprule 
{\footnotesize\textbf{Symbol}} & {\footnotesize\textbf{Value}} & {\footnotesize\textbf{Description}}\tabularnewline
\midrule 
{\footnotesize$\sigma_{p}$} & {\footnotesize 0.02} & {\footnotesize Eigen-aligned probe step (Algorithm \ref{alg:sparq-polish})}\tabularnewline
{\footnotesize$\sigma_{c}$} & {\footnotesize 0.10} & {\footnotesize Single-coordinate probe step (Algorithm \ref{alg:sparq-polish})}\tabularnewline
{\footnotesize$\rho_{sw}$} & {\footnotesize 0.02} & {\footnotesize Solis--Wets probe step (Algorithm \ref{alg:sparq-polish})}\tabularnewline
{\footnotesize$\sigma_{e}$} & {\footnotesize 1.0} & {\footnotesize Trajectory-echo scale (Algorithm \ref{alg:sparq-polish})}\tabularnewline
{\footnotesize --} & {\footnotesize$10^{-9}$ / $0.25$} & {\footnotesize Min./max. clamp for adaptive probe steps (not named
in the article)}\tabularnewline
{\footnotesize --} & {\footnotesize$10^{-3}$} & {\footnotesize Backoff threshold for low relative gain (not named in
the article)}\tabularnewline
{\footnotesize --} & {\footnotesize 8} & {\footnotesize Trajectory-echo buffer size (described in Section \ref{subsec:TrajectoryEcho}
as ``eight entries'', no symbol)}\tabularnewline
{\footnotesize --} & {\footnotesize$10^{-6}$} & {\footnotesize Meaningful-gain threshold for the $s_{p}$ counter (not
named in the article)}\tabularnewline
\bottomrule
\end{tabular}}
\end{table}


\subsection{Comparative Results on Standard-Dimensionality Benchmark Problems}

This subsection presents the outcome of the comparative evaluation
on the standard-dimensionality portion of the benchmark suite. SPARQ
is compared against nine state-of-the-art optimizers: ARQ2 \citep{ARQ2},
EA4Eig \citep{EI4Eig}, jSO \citep{JSO}, LM-CMA-ES \citep{LMCMAES},
mLSHADE-RL \citep{mLSHADE-RL}, NL-SHADE-LBC \citep{NL-SHADE-LBC},
SFCDE \citep{SFCDE}, TridentDE \citep{TridentDE}, and UDE3 \citep{UDE3}.
The comparison covers forty-five benchmark problems. These include
classical synthetic test functions alongside real-world engineering
formulations, such as antenna array design, hydrothermal scheduling,
economic load dispatch, and several chemical-engineering problems.
Problem dimensionality in this subsection ranges from as few as five
variables to several hundred. For each problem, every method was executed
independently over a fixed number of repeated runs under an identical
evaluation budget. Three complementary tables report the results.
Table \ref{tab:comparison1} presents detailed statistics for each
method and problem, namely the best value found, the mean value, the
success rate, the standard deviation, and the execution time. Table
\ref{tab:bestRankin1} reports the ranking of all ten methods on each
problem according to the best value achieved. Table \ref{tab:meanRanking1}
reports the corresponding ranking according to the mean value, offering
a stricter measure of reliability across repeated runs. Together,
these three tables allow the performance of SPARQ to be assessed both
in absolute terms and relative to its competitors, and both in terms
of peak performance and average-case consistency.
\begin{center}
\begin{landscape}
\begin{longtable}[c]{|c|c|c|c|c|c|c|c|c|c|c|c|}
\caption{Detailed Performance Comparison of SPARQ Against Seven State-of-the-Art
Optimizers on Standard and Real-World Benchmark Problems\label{tab:comparison1}}
\tabularnewline
\hline 
{\scriptsize\textbf{PROBLEM}} & {\scriptsize\textbf{METRIC}} & {\scriptsize\textbf{arq2}} & {\scriptsize\textbf{ea4eig}} & {\scriptsize\textbf{jso}} & {\scriptsize\textbf{lmcmaes}} & {\scriptsize\textbf{mlshaderl}} & {\scriptsize\textbf{nlshadelbc}} & {\scriptsize\textbf{sfcde}} & {\scriptsize\textbf{sparq}} & {\scriptsize\textbf{tridentde}} & {\scriptsize\textbf{ude3}}\tabularnewline
\endhead
\hline 
\multirow{5}{*}{\begin{cellvarwidth}[t]
\centering
{\scriptsize\textbf{alpine1}}{\scriptsize\par}

{\scriptsize\textbf{{[}d=100{]}}}
\end{cellvarwidth}} & {\scriptsize\textbf{Best}} & {\scriptsize 6.6000e-11} & {\scriptsize 5.3800e-09} & {\scriptsize 1.7000e-11} & {\scriptsize 0} & {\scriptsize 1.2200e-10} & {\scriptsize 0.00410456} & {\scriptsize 1.6690e-08} & {\scriptsize 0} & {\scriptsize 0.000371752} & {\scriptsize 1.8615e-07}\tabularnewline
\cline{2-12}
 & {\scriptsize\textbf{Mean}} & {\scriptsize 0.000176474} & {\scriptsize 2.4673e-07} & {\scriptsize 1.7126e-09} & {\scriptsize 0} & {\scriptsize 1.0039e-08} & {\scriptsize 0.0223941} & {\scriptsize 0.00086093} & {\scriptsize 0} & {\scriptsize 0.0033323} & {\scriptsize 1.9353e-06}\tabularnewline
\cline{2-12}
 & {\scriptsize\textbf{Rate\%}} & {\scriptsize 3\%} & {\scriptsize 3\%} & {\scriptsize 3\%} & {\scriptsize 100\%} & {\scriptsize 3\%} & {\scriptsize 3\%} & {\scriptsize 3\%} & {\scriptsize 100\%} & {\scriptsize 3\%} & {\scriptsize 3\%}\tabularnewline
\cline{2-12}
 & {\scriptsize\textbf{SD}} & {\scriptsize 0.000615062} & {\scriptsize 6.8327e-07} & {\scriptsize 4.0570e-09} & {\scriptsize 0} & {\scriptsize 2.2429e-08} & {\scriptsize 0.0268403} & {\scriptsize 0.0047049} & {\scriptsize 0} & {\scriptsize 0.00380758} & {\scriptsize 3.4386e-06}\tabularnewline
\cline{2-12}
 & {\scriptsize\textbf{Time(s)}} & {\scriptsize 0.319} & {\scriptsize 6.438} & {\scriptsize 0.153} & {\scriptsize 0.1} & {\scriptsize 0.174} & {\scriptsize 0.165} & {\scriptsize 0.338} & {\scriptsize 0.164} & {\scriptsize 0.834} & {\scriptsize 0.677}\tabularnewline
\hline 
\multirow{5}{*}{\begin{cellvarwidth}[t]
\centering
{\scriptsize\textbf{antennaarray}}{\scriptsize\par}

{\scriptsize\textbf{{[}d=12{]} \citep{cec2011}}}
\end{cellvarwidth}} & {\scriptsize\textbf{Best}} & {\scriptsize 0.00680964} & {\scriptsize 0.00680964} & {\scriptsize 0.00680964} & {\scriptsize 0.233659} & {\scriptsize 0.00680964} & {\scriptsize 0.00680964} & {\scriptsize 0.00680964} & {\scriptsize 0.00680964} & {\scriptsize 0.00680964} & {\scriptsize 0.00680964}\tabularnewline
\cline{2-12}
 & {\scriptsize\textbf{Mean}} & {\scriptsize 0.00681202} & {\scriptsize 0.00680964} & {\scriptsize 0.00680964} & {\scriptsize 0.356811} & {\scriptsize 0.00680964} & {\scriptsize 0.00680964} & {\scriptsize 0.0155814} & {\scriptsize 0.00680964} & {\scriptsize 0.00680975} & {\scriptsize 0.00680964}\tabularnewline
\cline{2-12}
 & {\scriptsize\textbf{Rate\%}} & {\scriptsize 67\%} & {\scriptsize 100\%} & {\scriptsize 100\%} & {\scriptsize 3\%} & {\scriptsize 100\%} & {\scriptsize 7\%} & {\scriptsize 93\%} & {\scriptsize 83\%} & {\scriptsize 13\%} & {\scriptsize 100\%}\tabularnewline
\cline{2-12}
 & {\scriptsize\textbf{SD}} & {\scriptsize 7.4675e-06} & {\scriptsize 1.7644e-18} & {\scriptsize 1.7644e-18} & {\scriptsize 0.0550951} & {\scriptsize 1.7644e-18} & {\scriptsize 2.2063e-08} & {\scriptsize 0.0336259} & {\scriptsize 1.7978e-10} & {\scriptsize 4.1947e-07} & {\scriptsize 1.7644e-18}\tabularnewline
\cline{2-12}
 & {\scriptsize\textbf{Time(s)}} & {\scriptsize 6.472} & {\scriptsize 3.292} & {\scriptsize 2.445} & {\scriptsize 0.91} & {\scriptsize 2.613} & {\scriptsize 1.227} & {\scriptsize 8.19} & {\scriptsize 2.002} & {\scriptsize 30.615} & {\scriptsize 17.592}\tabularnewline
\hline 
\multirow{5}{*}{\begin{cellvarwidth}[t]
\centering
{\scriptsize\textbf{antennaula}}{\scriptsize\par}

{\scriptsize\textbf{{[}d=10{]}}}
\end{cellvarwidth}} & {\scriptsize\textbf{Best}} & {\scriptsize 0.156149} & {\scriptsize 0.156149} & {\scriptsize 0.156149} & {\scriptsize 0.156308} & {\scriptsize 0.156149} & {\scriptsize 0.156149} & {\scriptsize 0.156149} & {\scriptsize 0.156149} & {\scriptsize 0.156149} & {\scriptsize 0.156149}\tabularnewline
\cline{2-12}
 & {\scriptsize\textbf{Mean}} & {\scriptsize 0.156149} & {\scriptsize 0.156149} & {\scriptsize 0.156149} & {\scriptsize 0.156822} & {\scriptsize 0.156149} & {\scriptsize 0.156149} & {\scriptsize 0.156149} & {\scriptsize 0.156149} & {\scriptsize 0.156149} & {\scriptsize 0.156149}\tabularnewline
\cline{2-12}
 & {\scriptsize\textbf{Rate\%}} & {\scriptsize 100\%} & {\scriptsize 100\%} & {\scriptsize 100\%} & {\scriptsize 3\%} & {\scriptsize 100\%} & {\scriptsize 93\%} & {\scriptsize 100\%} & {\scriptsize 100\%} & {\scriptsize 13\%} & {\scriptsize 100\%}\tabularnewline
\cline{2-12}
 & {\scriptsize\textbf{SD}} & {\scriptsize 2.8230e-17} & {\scriptsize 2.8230e-17} & {\scriptsize 2.8230e-17} & {\scriptsize 0.00037492} & {\scriptsize 2.8230e-17} & {\scriptsize 1.2181e-10} & {\scriptsize 2.8230e-17} & {\scriptsize 2.8230e-17} & {\scriptsize 1.3557e-09} & {\scriptsize 2.8230e-17}\tabularnewline
\cline{2-12}
 & {\scriptsize\textbf{Time(s)}} & {\scriptsize 12.108} & {\scriptsize 6.358} & {\scriptsize 4.825} & {\scriptsize 1.626} & {\scriptsize 4.64} & {\scriptsize 2.424} & {\scriptsize 16.327} & {\scriptsize 3.991} & {\scriptsize 56.682} & {\scriptsize 32.536}\tabularnewline
\hline 
\multirow{5}{*}{\begin{cellvarwidth}[t]
\centering
{\scriptsize\textbf{bucherastrigin }}{\scriptsize\par}

{\scriptsize\textbf{{[}d=50{]}}}
\end{cellvarwidth}} & {\scriptsize\textbf{Best}} & {\scriptsize 3.03752} & {\scriptsize 10.6264} & {\scriptsize 7.9597} & {\scriptsize 0} & {\scriptsize 25.8689} & {\scriptsize 16.9143} & {\scriptsize 5.45194} & {\scriptsize 0} & {\scriptsize 42.7923} & {\scriptsize 37.8084}\tabularnewline
\cline{2-12}
 & {\scriptsize\textbf{Mean}} & {\scriptsize 21.4995} & {\scriptsize 18.457} & {\scriptsize 25.1062} & {\scriptsize 0} & {\scriptsize 39.0687} & {\scriptsize 71.4967} & {\scriptsize 24.0398} & {\scriptsize 0} & {\scriptsize 87.4512} & {\scriptsize 50.6661}\tabularnewline
\cline{2-12}
 & {\scriptsize\textbf{Rate\%}} & {\scriptsize 3\%} & {\scriptsize 3\%} & {\scriptsize 3\%} & {\scriptsize 100\%} & {\scriptsize 3\%} & {\scriptsize 3\%} & {\scriptsize 3\%} & {\scriptsize 100\%} & {\scriptsize 3\%} & {\scriptsize 3\%}\tabularnewline
\cline{2-12}
 & {\scriptsize\textbf{SD}} & {\scriptsize 10.5954} & {\scriptsize 4.53841} & {\scriptsize 17.0473} & {\scriptsize 0} & {\scriptsize 5.32883} & {\scriptsize 32.3878} & {\scriptsize 13.3618} & {\scriptsize 0} & {\scriptsize 31.335} & {\scriptsize 6.69681}\tabularnewline
\cline{2-12}
 & {\scriptsize\textbf{Time(s)}} & {\scriptsize 0.472} & {\scriptsize 2.25} & {\scriptsize 0.202} & {\scriptsize 0.097} & {\scriptsize 0.218} & {\scriptsize 0.17} & {\scriptsize 0.547} & {\scriptsize 0.241} & {\scriptsize 1.66} & {\scriptsize 1.182}\tabularnewline
\hline 
\multirow{5}{*}{\begin{cellvarwidth}[t]
\centering
{\scriptsize\textbf{cosinemixture}}{\scriptsize\par}

{\scriptsize\textbf{{[}d=50{]} \citep{cec2011}}}
\end{cellvarwidth}} & {\scriptsize\textbf{Best}} & {\scriptsize -5} & {\scriptsize -5} & {\scriptsize -5} & {\scriptsize -5} & {\scriptsize -5} & {\scriptsize -5} & {\scriptsize -5} & {\scriptsize -5} & {\scriptsize -5} & {\scriptsize -5}\tabularnewline
\cline{2-12}
 & {\scriptsize\textbf{Mean}} & {\scriptsize -5} & {\scriptsize -5} & {\scriptsize -5} & {\scriptsize -5} & {\scriptsize -4.99015} & {\scriptsize -4.95074} & {\scriptsize -4.78325} & {\scriptsize -5} & {\scriptsize -4.94089} & {\scriptsize -5}\tabularnewline
\cline{2-12}
 & {\scriptsize\textbf{Rate\%}} & {\scriptsize 100\%} & {\scriptsize 100\%} & {\scriptsize 100\%} & {\scriptsize 100\%} & {\scriptsize 93\%} & {\scriptsize 70\%} & {\scriptsize 27\%} & {\scriptsize 100\%} & {\scriptsize 73\%} & {\scriptsize 100\%}\tabularnewline
\cline{2-12}
 & {\scriptsize\textbf{SD}} & {\scriptsize 0} & {\scriptsize 0} & {\scriptsize 0} & {\scriptsize 0} & {\scriptsize 0.0374941} & {\scriptsize 0.0807896} & {\scriptsize 0.215154} & {\scriptsize 0} & {\scriptsize 0.113813} & {\scriptsize 0}\tabularnewline
\cline{2-12}
 & {\scriptsize\textbf{Time(s)}} & {\scriptsize 0.242} & {\scriptsize 0.563} & {\scriptsize 0.155} & {\scriptsize 0.149} & {\scriptsize 0.164} & {\scriptsize 0.167} & {\scriptsize 0.208} & {\scriptsize 0.19} & {\scriptsize 0.482} & {\scriptsize 0.468}\tabularnewline
\hline 
\multirow{5}{*}{\begin{cellvarwidth}[t]
\centering
{\scriptsize\textbf{ded1}}{\scriptsize\par}

{\scriptsize\textbf{{[}d=120{]} \citep{cec2011}}}
\end{cellvarwidth}} & {\scriptsize\textbf{Best}} & {\scriptsize 130643} & {\scriptsize 130643} & {\scriptsize 130643} & {\scriptsize 130774} & {\scriptsize 130643} & {\scriptsize 130855} & {\scriptsize 130746} & {\scriptsize 130643} & {\scriptsize 130841} & {\scriptsize 130650}\tabularnewline
\cline{2-12}
 & {\scriptsize\textbf{Mean}} & {\scriptsize 130645} & {\scriptsize 130644} & {\scriptsize 130643} & {\scriptsize 130892} & {\scriptsize 130644} & {\scriptsize 130903} & {\scriptsize 130788} & {\scriptsize 130644} & {\scriptsize 130905} & {\scriptsize 130659}\tabularnewline
\cline{2-12}
 & {\scriptsize\textbf{Rate\%}} & {\scriptsize 3\%} & {\scriptsize 3\%} & {\scriptsize 3\%} & {\scriptsize 3\%} & {\scriptsize 3\%} & {\scriptsize 3\%} & {\scriptsize 3\%} & {\scriptsize 3\%} & {\scriptsize 3\%} & {\scriptsize 3\%}\tabularnewline
\cline{2-12}
 & {\scriptsize\textbf{SD}} & {\scriptsize 1.95641} & {\scriptsize 0.709733} & {\scriptsize 0.163697} & {\scriptsize 65.8632} & {\scriptsize 0.895705} & {\scriptsize 29.644} & {\scriptsize 20.0706} & {\scriptsize 0.719647} & {\scriptsize 38.2176} & {\scriptsize 7.13931}\tabularnewline
\cline{2-12}
 & {\scriptsize\textbf{Time(s)}} & {\scriptsize 0.382} & {\scriptsize 0.236} & {\scriptsize 0.205} & {\scriptsize 0.13} & {\scriptsize 0.184} & {\scriptsize 0.164} & {\scriptsize 0.399} & {\scriptsize 0.165} & {\scriptsize 0.944} & {\scriptsize 0.778}\tabularnewline
\hline 
\multirow{5}{*}{\begin{cellvarwidth}[t]
\centering
{\scriptsize\textbf{ded2}}{\scriptsize\par}

{\scriptsize\textbf{{[}d=216{]} \citep{cec2011}}}
\end{cellvarwidth}} & {\scriptsize\textbf{Best}} & {\scriptsize 165137} & {\scriptsize 165128} & {\scriptsize 164964} & {\scriptsize 166873} & {\scriptsize 165290} & {\scriptsize 167727} & {\scriptsize 166356} & {\scriptsize 165202} & {\scriptsize 165953} & {\scriptsize 165108}\tabularnewline
\cline{2-12}
 & {\scriptsize\textbf{Mean}} & {\scriptsize 165656} & {\scriptsize 165599} & {\scriptsize 165626} & {\scriptsize 168467} & {\scriptsize 165639} & {\scriptsize 168779} & {\scriptsize 166961} & {\scriptsize 165672} & {\scriptsize 166519} & {\scriptsize 165791}\tabularnewline
\cline{2-12}
 & {\scriptsize\textbf{Rate\%}} & {\scriptsize 3\%} & {\scriptsize 3\%} & {\scriptsize 3\%} & {\scriptsize 3\%} & {\scriptsize 3\%} & {\scriptsize 3\%} & {\scriptsize 3\%} & {\scriptsize 3\%} & {\scriptsize 3\%} & {\scriptsize 3\%}\tabularnewline
\cline{2-12}
 & {\scriptsize\textbf{SD}} & {\scriptsize 398.646} & {\scriptsize 215.743} & {\scriptsize 245.073} & {\scriptsize 1248.51} & {\scriptsize 237.79} & {\scriptsize 428.514} & {\scriptsize 316.334} & {\scriptsize 265.268} & {\scriptsize 288.062} & {\scriptsize 248.323}\tabularnewline
\cline{2-12}
 & {\scriptsize\textbf{Time(s)}} & {\scriptsize 0.739} & {\scriptsize 0.427} & {\scriptsize 0.306} & {\scriptsize 0.283} & {\scriptsize 0.268} & {\scriptsize 0.185} & {\scriptsize 1.011} & {\scriptsize 0.245} & {\scriptsize 2.357} & {\scriptsize 1.627}\tabularnewline
\hline 
\multirow{5}{*}{\begin{cellvarwidth}[t]
\centering
{\scriptsize\textbf{eld1}}{\scriptsize\par}

{\scriptsize\textbf{{[}d=6{]} \citep{cec2011}}}
\end{cellvarwidth}} & {\scriptsize\textbf{Best}} & {\scriptsize 2967.25} & {\scriptsize 2967.25} & {\scriptsize 2967.25} & {\scriptsize 2967.25} & {\scriptsize 2967.25} & {\scriptsize 2967.25} & {\scriptsize 2967.25} & {\scriptsize 2967.25} & {\scriptsize 2967.25} & {\scriptsize 2967.25}\tabularnewline
\cline{2-12}
 & {\scriptsize\textbf{Mean}} & {\scriptsize 2967.52} & {\scriptsize 2967.54} & {\scriptsize 2967.58} & {\scriptsize 3190.82} & {\scriptsize 2975.72} & {\scriptsize 2967.6} & {\scriptsize 2969.06} & {\scriptsize 2969.13} & {\scriptsize 2978.99} & {\scriptsize 2974.84}\tabularnewline
\cline{2-12}
 & {\scriptsize\textbf{Rate\%}} & {\scriptsize 57\%} & {\scriptsize 53\%} & {\scriptsize 47\%} & {\scriptsize 7\%} & {\scriptsize 17\%} & {\scriptsize 43\%} & {\scriptsize 57\%} & {\scriptsize 53\%} & {\scriptsize 3\%} & {\scriptsize 30\%}\tabularnewline
\cline{2-12}
 & {\scriptsize\textbf{SD}} & {\scriptsize 0.310574} & {\scriptsize 0.312675} & {\scriptsize 0.312675} & {\scriptsize 289.445} & {\scriptsize 5.87434} & {\scriptsize 0.310574} & {\scriptsize 6.67985} & {\scriptsize 4.26807} & {\scriptsize 3.10799} & {\scriptsize 6.18104}\tabularnewline
\cline{2-12}
 & {\scriptsize\textbf{Time(s)}} & {\scriptsize 0.113} & {\scriptsize 0.145} & {\scriptsize 0.164} & {\scriptsize 0.177} & {\scriptsize 0.158} & {\scriptsize 0.17} & {\scriptsize 0.111} & {\scriptsize 0.164} & {\scriptsize 0.209} & {\scriptsize 0.283}\tabularnewline
\hline 
\multirow{5}{*}{\begin{cellvarwidth}[t]
\centering
{\scriptsize\textbf{eld2}}{\scriptsize\par}

{\scriptsize\textbf{{[}d=13{]} \citep{cec2011}}}
\end{cellvarwidth}} & {\scriptsize\textbf{Best}} & {\scriptsize 17864.4} & {\scriptsize 17864.7} & {\scriptsize 17924.4} & {\scriptsize 17936.7} & {\scriptsize 17866.9} & {\scriptsize 17924.4} & {\scriptsize 17874.9} & {\scriptsize 17874.5} & {\scriptsize 17867.4} & {\scriptsize 17879.7}\tabularnewline
\cline{2-12}
 & {\scriptsize\textbf{Mean}} & {\scriptsize 17881.5} & {\scriptsize 17895.1} & {\scriptsize 17954.8} & {\scriptsize 18072} & {\scriptsize 17881.4} & {\scriptsize 18005.1} & {\scriptsize 17917.2} & {\scriptsize 17923.1} & {\scriptsize 17938.9} & {\scriptsize 17885.4}\tabularnewline
\cline{2-12}
 & {\scriptsize\textbf{Rate\%}} & {\scriptsize 3\%} & {\scriptsize 7\%} & {\scriptsize 13\%} & {\scriptsize 3\%} & {\scriptsize 10\%} & {\scriptsize 3\%} & {\scriptsize 3\%} & {\scriptsize 3\%} & {\scriptsize 3\%} & {\scriptsize 10\%}\tabularnewline
\cline{2-12}
 & {\scriptsize\textbf{SD}} & {\scriptsize 21.6346} & {\scriptsize 31.1588} & {\scriptsize 17.4242} & {\scriptsize 88.4322} & {\scriptsize 8.64524} & {\scriptsize 58.2737} & {\scriptsize 22.7828} & {\scriptsize 32.4557} & {\scriptsize 50.907} & {\scriptsize 10.9639}\tabularnewline
\cline{2-12}
 & {\scriptsize\textbf{Time(s)}} & {\scriptsize 0.137} & {\scriptsize 0.134} & {\scriptsize 0.158} & {\scriptsize 0.173} & {\scriptsize 0.16} & {\scriptsize 0.17} & {\scriptsize 0.126} & {\scriptsize 0.16} & {\scriptsize 0.253} & {\scriptsize 0.332}\tabularnewline
\hline 
\multirow{5}{*}{\begin{cellvarwidth}[t]
\centering
{\scriptsize\textbf{eld3}}{\scriptsize\par}

{\scriptsize\textbf{{[}d=15{]} \citep{cec2011}}}
\end{cellvarwidth}} & {\scriptsize\textbf{Best}} & {\scriptsize 32367.6} & {\scriptsize 32375.3} & {\scriptsize 32458.9} & {\scriptsize 32760.6} & {\scriptsize 32367.6} & {\scriptsize 32597.3} & {\scriptsize 32401.5} & {\scriptsize 32375.3} & {\scriptsize 32367.6} & {\scriptsize 32367.6}\tabularnewline
\cline{2-12}
 & {\scriptsize\textbf{Mean}} & {\scriptsize 32399.4} & {\scriptsize 32425.9} & {\scriptsize 32554} & {\scriptsize 33110} & {\scriptsize 32405.4} & {\scriptsize 32798.9} & {\scriptsize 32488.3} & {\scriptsize 32440.5} & {\scriptsize 32494} & {\scriptsize 32434.2}\tabularnewline
\cline{2-12}
 & {\scriptsize\textbf{Rate\%}} & {\scriptsize 7\%} & {\scriptsize 3\%} & {\scriptsize 3\%} & {\scriptsize 3\%} & {\scriptsize 53\%} & {\scriptsize 3\%} & {\scriptsize 3\%} & {\scriptsize 7\%} & {\scriptsize 13\%} & {\scriptsize 33\%}\tabularnewline
\cline{2-12}
 & {\scriptsize\textbf{SD}} & {\scriptsize 29.7244} & {\scriptsize 40.967} & {\scriptsize 40.8494} & {\scriptsize 208.751} & {\scriptsize 74.2414} & {\scriptsize 69.4704} & {\scriptsize 53.4577} & {\scriptsize 47.9137} & {\scriptsize 100.649} & {\scriptsize 73.6464}\tabularnewline
\cline{2-12}
 & {\scriptsize\textbf{Time(s)}} & {\scriptsize 0.151} & {\scriptsize 0.151} & {\scriptsize 0.155} & {\scriptsize 0.171} & {\scriptsize 0.16} & {\scriptsize 0.17} & {\scriptsize 0.145} & {\scriptsize 0.156} & {\scriptsize 0.286} & {\scriptsize 0.355}\tabularnewline
\hline 
\multirow{5}{*}{\begin{cellvarwidth}[t]
\centering
{\scriptsize\textbf{eld4}}{\scriptsize\par}

{\scriptsize\textbf{{[}d=40{]} \citep{cec2011}}}
\end{cellvarwidth}} & {\scriptsize\textbf{Best}} & {\scriptsize 121064} & {\scriptsize 121096} & {\scriptsize 121152} & {\scriptsize 123426} & {\scriptsize 121064} & {\scriptsize 122082} & {\scriptsize 121647} & {\scriptsize 121096} & {\scriptsize 121075} & {\scriptsize 121075}\tabularnewline
\cline{2-12}
 & {\scriptsize\textbf{Mean}} & {\scriptsize 121188} & {\scriptsize 121335} & {\scriptsize 121390} & {\scriptsize 125460} & {\scriptsize 121181} & {\scriptsize 123934} & {\scriptsize 122353} & {\scriptsize 121399} & {\scriptsize 121403} & {\scriptsize 121214}\tabularnewline
\cline{2-12}
 & {\scriptsize\textbf{Rate\%}} & {\scriptsize 3\%} & {\scriptsize 3\%} & {\scriptsize 3\%} & {\scriptsize 3\%} & {\scriptsize 3\%} & {\scriptsize 3\%} & {\scriptsize 3\%} & {\scriptsize 3\%} & {\scriptsize 3\%} & {\scriptsize 3\%}\tabularnewline
\cline{2-12}
 & {\scriptsize\textbf{SD}} & {\scriptsize 106.955} & {\scriptsize 174.933} & {\scriptsize 194.4} & {\scriptsize 965.109} & {\scriptsize 85.3427} & {\scriptsize 816.891} & {\scriptsize 376.605} & {\scriptsize 202.602} & {\scriptsize 191.838} & {\scriptsize 116.902}\tabularnewline
\cline{2-12}
 & {\scriptsize\textbf{Time(s)}} & {\scriptsize 0.226} & {\scriptsize 0.503} & {\scriptsize 0.157} & {\scriptsize 0.154} & {\scriptsize 0.168} & {\scriptsize 0.166} & {\scriptsize 0.239} & {\scriptsize 0.192} & {\scriptsize 0.576} & {\scriptsize 0.504}\tabularnewline
\hline 
\multirow{5}{*}{\begin{cellvarwidth}[t]
\centering
{\scriptsize\textbf{eld5}}{\scriptsize\par}

{\scriptsize\textbf{{[}d=140{]} \citep{cec2011}}}
\end{cellvarwidth}} & {\scriptsize\textbf{Best}} & {\scriptsize 508614} & {\scriptsize 508614} & {\scriptsize 508614} & {\scriptsize 510807} & {\scriptsize 508614} & {\scriptsize 508877} & {\scriptsize 508715} & {\scriptsize 508614} & {\scriptsize 508659} & {\scriptsize 508618}\tabularnewline
\cline{2-12}
 & {\scriptsize\textbf{Mean}} & {\scriptsize 508619} & {\scriptsize 508614} & {\scriptsize 508614} & {\scriptsize 511210} & {\scriptsize 508616} & {\scriptsize 509014} & {\scriptsize 508786} & {\scriptsize 508615} & {\scriptsize 508702} & {\scriptsize 508635}\tabularnewline
\cline{2-12}
 & {\scriptsize\textbf{Rate\%}} & {\scriptsize 3\%} & {\scriptsize 3\%} & {\scriptsize 3\%} & {\scriptsize 3\%} & {\scriptsize 3\%} & {\scriptsize 3\%} & {\scriptsize 3\%} & {\scriptsize 3\%} & {\scriptsize 3\%} & {\scriptsize 3\%}\tabularnewline
\cline{2-12}
 & {\scriptsize\textbf{SD}} & {\scriptsize 9.98242} & {\scriptsize 0.106988} & {\scriptsize 0.0852147} & {\scriptsize 170.472} & {\scriptsize 3.64487} & {\scriptsize 79.002} & {\scriptsize 31.9054} & {\scriptsize 0.724255} & {\scriptsize 26.958} & {\scriptsize 15.0339}\tabularnewline
\cline{2-12}
 & {\scriptsize\textbf{Time(s)}} & {\scriptsize 0.416} & {\scriptsize 0.395} & {\scriptsize 0.375} & {\scriptsize 0.136} & {\scriptsize 0.18} & {\scriptsize 0.168} & {\scriptsize 0.4} & {\scriptsize 0.162} & {\scriptsize 0.861} & {\scriptsize 0.724}\tabularnewline
\hline 
\multirow{5}{*}{\begin{cellvarwidth}[t]
\centering
{\scriptsize\textbf{fmsynth}}{\scriptsize\par}

{\scriptsize\textbf{{[}d=6{]} \citep{cec2011}}}
\end{cellvarwidth}} & {\scriptsize\textbf{Best}} & {\scriptsize 0} & {\scriptsize 0} & {\scriptsize 0} & {\scriptsize 0.274042} & {\scriptsize 0.116158} & {\scriptsize 0.000131621} & {\scriptsize 9.6570e-06} & {\scriptsize 0} & {\scriptsize 0.116158} & {\scriptsize 6.6499e-08}\tabularnewline
\cline{2-12}
 & {\scriptsize\textbf{Mean}} & {\scriptsize 0.115033} & {\scriptsize 0.098221} & {\scriptsize 0.0972373} & {\scriptsize 0.331597} & {\scriptsize 0.121941} & {\scriptsize 0.10375} & {\scriptsize 0.108498} & {\scriptsize 0.105976} & {\scriptsize 0.126956} & {\scriptsize 0.115608}\tabularnewline
\cline{2-12}
 & {\scriptsize\textbf{Rate\%}} & {\scriptsize 3\%} & {\scriptsize 17\%} & {\scriptsize 10\%} & {\scriptsize 3\%} & {\scriptsize 83\%} & {\scriptsize 3\%} & {\scriptsize 3\%} & {\scriptsize 10\%} & {\scriptsize 73\%} & {\scriptsize 3\%}\tabularnewline
\cline{2-12}
 & {\scriptsize\textbf{SD}} & {\scriptsize 0.0228626} & {\scriptsize 0.0453469} & {\scriptsize 0.0442939} & {\scriptsize 0.0236314} & {\scriptsize 0.0171161} & {\scriptsize 0.0337617} & {\scriptsize 0.0329245} & {\scriptsize 0.0363381} & {\scriptsize 0.023476} & {\scriptsize 0.0245231}\tabularnewline
\cline{2-12}
 & {\scriptsize\textbf{Time(s)}} & {\scriptsize 0.47} & {\scriptsize 0.28} & {\scriptsize 0.226} & {\scriptsize 0.151} & {\scriptsize 0.219} & {\scriptsize 0.179} & {\scriptsize 0.657} & {\scriptsize 0.208} & {\scriptsize 1.859} & {\scriptsize 1.356}\tabularnewline
\hline 
\multirow{5}{*}{\begin{cellvarwidth}[t]
\centering
{\scriptsize\textbf{gallagher101}}{\scriptsize\par}

{\scriptsize\textbf{{[}d=16{]} \citep{cec2011}}}
\end{cellvarwidth}} & {\scriptsize\textbf{Best}} & {\scriptsize 0} & {\scriptsize 0} & {\scriptsize 0} & {\scriptsize 0.691857} & {\scriptsize 0.691857} & {\scriptsize 0.929997} & {\scriptsize 0} & {\scriptsize 0.929997} & {\scriptsize 0} & {\scriptsize 0}\tabularnewline
\cline{2-12}
 & {\scriptsize\textbf{Mean}} & {\scriptsize 1.81765} & {\scriptsize 1.88988} & {\scriptsize 2.34955} & {\scriptsize 9.20322} & {\scriptsize 1.88241} & {\scriptsize 4.38908} & {\scriptsize 3.25587} & {\scriptsize 3.11339} & {\scriptsize 2.9818} & {\scriptsize 2.21648}\tabularnewline
\cline{2-12}
 & {\scriptsize\textbf{Rate\%}} & {\scriptsize 13\%} & {\scriptsize 10\%} & {\scriptsize 3\%} & {\scriptsize 7\%} & {\scriptsize 3\%} & {\scriptsize 23\%} & {\scriptsize 7\%} & {\scriptsize 37\%} & {\scriptsize 13\%} & {\scriptsize 3\%}\tabularnewline
\cline{2-12}
 & {\scriptsize\textbf{SD}} & {\scriptsize 1.59835} & {\scriptsize 1.6648} & {\scriptsize 1.72389} & {\scriptsize 10.2299} & {\scriptsize 1.49672} & {\scriptsize 4.17408} & {\scriptsize 2.95444} & {\scriptsize 2.73578} & {\scriptsize 3.50141} & {\scriptsize 1.61107}\tabularnewline
\cline{2-12}
 & {\scriptsize\textbf{Time(s)}} & {\scriptsize 1.789} & {\scriptsize 0.857} & {\scriptsize 0.63} & {\scriptsize 0.256} & {\scriptsize 0.633} & {\scriptsize 0.324} & {\scriptsize 2.175} & {\scriptsize 0.565} & {\scriptsize 7.359} & {\scriptsize 4.419}\tabularnewline
\hline 
\multirow{5}{*}{\begin{cellvarwidth}[t]
\centering
{\scriptsize\textbf{gallagher21}}{\scriptsize\par}

{\scriptsize\textbf{{[}d=16{]}}}
\end{cellvarwidth}} & {\scriptsize\textbf{Best}} & {\scriptsize 0} & {\scriptsize 0} & {\scriptsize 0} & {\scriptsize 0} & {\scriptsize 0} & {\scriptsize 0} & {\scriptsize 0} & {\scriptsize 0} & {\scriptsize 0} & {\scriptsize 0}\tabularnewline
\cline{2-12}
 & {\scriptsize\textbf{Mean}} & {\scriptsize 4.99458} & {\scriptsize 1.72255} & {\scriptsize 6.80626} & {\scriptsize 10.9842} & {\scriptsize 3.85307} & {\scriptsize 8.74813} & {\scriptsize 5.86152} & {\scriptsize 4.7033} & {\scriptsize 5.23155} & {\scriptsize 5.2396}\tabularnewline
\cline{2-12}
 & {\scriptsize\textbf{Rate\%}} & {\scriptsize 50\%} & {\scriptsize 80\%} & {\scriptsize 53\%} & {\scriptsize 17\%} & {\scriptsize 43\%} & {\scriptsize 33\%} & {\scriptsize 53\%} & {\scriptsize 53\%} & {\scriptsize 30\%} & {\scriptsize 37\%}\tabularnewline
\cline{2-12}
 & {\scriptsize\textbf{SD}} & {\scriptsize 6.63609} & {\scriptsize 4.13367} & {\scriptsize 7.40059} & {\scriptsize 5.64072} & {\scriptsize 6.16112} & {\scriptsize 6.74433} & {\scriptsize 8.16359} & {\scriptsize 6.14221} & {\scriptsize 5.88774} & {\scriptsize 6.28461}\tabularnewline
\cline{2-12}
 & {\scriptsize\textbf{Time(s)}} & {\scriptsize 0.474} & {\scriptsize 0.24} & {\scriptsize 0.196} & {\scriptsize 0.115} & {\scriptsize 0.214} & {\scriptsize 0.172} & {\scriptsize 0.539} & {\scriptsize 0.206} & {\scriptsize 1.668} & {\scriptsize 1.128}\tabularnewline
\hline 
\multirow{5}{*}{\begin{cellvarwidth}[t]
\centering
{\scriptsize\textbf{griewankrosenbrock}}{\scriptsize\par}

{\scriptsize\textbf{{[}d=50{]}}}
\end{cellvarwidth}} & {\scriptsize\textbf{Best}} & {\scriptsize 0.975471} & {\scriptsize 2.99495} & {\scriptsize 1.64009} & {\scriptsize 1.99363} & {\scriptsize 16.6786} & {\scriptsize 0.545139} & {\scriptsize 2.9697} & {\scriptsize 0.756589} & {\scriptsize 7.68858} & {\scriptsize 25.3427}\tabularnewline
\cline{2-12}
 & {\scriptsize\textbf{Mean}} & {\scriptsize 2.10878} & {\scriptsize 5.82533} & {\scriptsize 2.22636} & {\scriptsize 4.73453} & {\scriptsize 24.7386} & {\scriptsize 0.773744} & {\scriptsize 3.7164} & {\scriptsize 1.60007} & {\scriptsize 10.9174} & {\scriptsize 27.4377}\tabularnewline
\cline{2-12}
 & {\scriptsize\textbf{Rate\%}} & {\scriptsize 3\%} & {\scriptsize 3\%} & {\scriptsize 3\%} & {\scriptsize 3\%} & {\scriptsize 3\%} & {\scriptsize 3\%} & {\scriptsize 3\%} & {\scriptsize 3\%} & {\scriptsize 3\%} & {\scriptsize 3\%}\tabularnewline
\cline{2-12}
 & {\scriptsize\textbf{SD}} & {\scriptsize 0.715515} & {\scriptsize 1.72956} & {\scriptsize 0.49828} & {\scriptsize 2.25592} & {\scriptsize 2.28942} & {\scriptsize 0.144207} & {\scriptsize 0.281837} & {\scriptsize 0.525703} & {\scriptsize 1.50984} & {\scriptsize 1.01623}\tabularnewline
\cline{2-12}
 & {\scriptsize\textbf{Time(s)}} & {\scriptsize 0.265} & {\scriptsize 2.152} & {\scriptsize 0.15} & {\scriptsize 0.145} & {\scriptsize 0.165} & {\scriptsize 0.164} & {\scriptsize 0.24} & {\scriptsize 0.239} & {\scriptsize 0.671} & {\scriptsize 0.602}\tabularnewline
\hline 
\multirow{5}{*}{\begin{cellvarwidth}[t]
\centering
{\scriptsize\textbf{heatexchanger}}{\scriptsize\par}

{\scriptsize\textbf{{[}d=22{]}}}
\end{cellvarwidth}} & {\scriptsize\textbf{Best}} & {\scriptsize 137715} & {\scriptsize 137682} & {\scriptsize 137682} & {\scriptsize 137682} & {\scriptsize 137682} & {\scriptsize 137715} & {\scriptsize 137715} & {\scriptsize 137715} & {\scriptsize 137682} & {\scriptsize 137682}\tabularnewline
\cline{2-12}
 & {\scriptsize\textbf{Mean}} & {\scriptsize 137715} & {\scriptsize 137687} & {\scriptsize 137713} & {\scriptsize 139785} & {\scriptsize 137682} & {\scriptsize 137715} & {\scriptsize 137715} & {\scriptsize 137715} & {\scriptsize 137682} & {\scriptsize 137682}\tabularnewline
\cline{2-12}
 & {\scriptsize\textbf{Rate\%}} & {\scriptsize 100\%} & {\scriptsize 83\%} & {\scriptsize 7\%} & {\scriptsize 33\%} & {\scriptsize 100\%} & {\scriptsize 100\%} & {\scriptsize 100\%} & {\scriptsize 100\%} & {\scriptsize 100\%} & {\scriptsize 100\%}\tabularnewline
\cline{2-12}
 & {\scriptsize\textbf{SD}} & {\scriptsize 0} & {\scriptsize 12.6233} & {\scriptsize 8.4491} & {\scriptsize 2930.89} & {\scriptsize 8.8804e-11} & {\scriptsize 0} & {\scriptsize 0} & {\scriptsize 0} & {\scriptsize 8.8804e-11} & {\scriptsize 8.8804e-11}\tabularnewline
\cline{2-12}
 & {\scriptsize\textbf{Time(s)}} & {\scriptsize 0.161} & {\scriptsize 0.305} & {\scriptsize 0.131} & {\scriptsize 0.172} & {\scriptsize 0.163} & {\scriptsize 0.164} & {\scriptsize 0.145} & {\scriptsize 0.176} & {\scriptsize 0.265} & {\scriptsize 0.355}\tabularnewline
\hline 
\multirow{5}{*}{\begin{cellvarwidth}[t]
\centering
{\scriptsize\textbf{hydrothermal}}{\scriptsize\par}

{\scriptsize\textbf{{[}d=120{]} \citep{cec2011}}}
\end{cellvarwidth}} & {\scriptsize\textbf{Best}} & {\scriptsize 141656} & {\scriptsize 141670} & {\scriptsize 141662} & {\scriptsize 231975} & {\scriptsize 141656} & {\scriptsize 145194} & {\scriptsize 141841} & {\scriptsize 141656} & {\scriptsize 141795} & {\scriptsize 141672}\tabularnewline
\cline{2-12}
 & {\scriptsize\textbf{Mean}} & {\scriptsize 141660} & {\scriptsize 141694} & {\scriptsize 141675} & {\scriptsize 351889} & {\scriptsize 141659} & {\scriptsize 151361} & {\scriptsize 142079} & {\scriptsize 141669} & {\scriptsize 141876} & {\scriptsize 141704}\tabularnewline
\cline{2-12}
 & {\scriptsize\textbf{Rate\%}} & {\scriptsize 3\%} & {\scriptsize 3\%} & {\scriptsize 3\%} & {\scriptsize 3\%} & {\scriptsize 3\%} & {\scriptsize 3\%} & {\scriptsize 3\%} & {\scriptsize 3\%} & {\scriptsize 3\%} & {\scriptsize 3\%}\tabularnewline
\cline{2-12}
 & {\scriptsize\textbf{SD}} & {\scriptsize 6.43564} & {\scriptsize 11.5562} & {\scriptsize 10.117} & {\scriptsize 67903} & {\scriptsize 2.94957} & {\scriptsize 3944.21} & {\scriptsize 299.956} & {\scriptsize 9.8093} & {\scriptsize 85.9947} & {\scriptsize 18.8546}\tabularnewline
\cline{2-12}
 & {\scriptsize\textbf{Time(s)}} & {\scriptsize 0.365} & {\scriptsize 0.411} & {\scriptsize 0.444} & {\scriptsize 0.126} & {\scriptsize 0.175} & {\scriptsize 0.169} & {\scriptsize 0.378} & {\scriptsize 0.172} & {\scriptsize 0.889} & {\scriptsize 0.717}\tabularnewline
\hline 
\multirow{5}{*}{\begin{cellvarwidth}[t]
\centering
{\scriptsize\textbf{katsuura}}{\scriptsize\par}

{\scriptsize\textbf{{[}d=50{]}}}
\end{cellvarwidth}} & {\scriptsize\textbf{Best}} & {\scriptsize 0.019147} & {\scriptsize 0} & {\scriptsize 0} & {\scriptsize 0} & {\scriptsize 0} & {\scriptsize 0.00735398} & {\scriptsize 0.539615} & {\scriptsize 0} & {\scriptsize 0} & {\scriptsize 0}\tabularnewline
\cline{2-12}
 & {\scriptsize\textbf{Mean}} & {\scriptsize 0.0712054} & {\scriptsize 6.3333e-13} & {\scriptsize 0} & {\scriptsize 0} & {\scriptsize 0} & {\scriptsize 0.471516} & {\scriptsize 0.660281} & {\scriptsize 0} & {\scriptsize 0} & {\scriptsize 0}\tabularnewline
\cline{2-12}
 & {\scriptsize\textbf{Rate\%}} & {\scriptsize 3\%} & {\scriptsize 97\%} & {\scriptsize 100\%} & {\scriptsize 100\%} & {\scriptsize 100\%} & {\scriptsize 3\%} & {\scriptsize 3\%} & {\scriptsize 100\%} & {\scriptsize 100\%} & {\scriptsize 100\%}\tabularnewline
\cline{2-12}
 & {\scriptsize\textbf{SD}} & {\scriptsize 0.0768892} & {\scriptsize 3.2851e-12} & {\scriptsize 0} & {\scriptsize 0} & {\scriptsize 0} & {\scriptsize 1.66103} & {\scriptsize 0.0674435} & {\scriptsize 0} & {\scriptsize 0} & {\scriptsize 0}\tabularnewline
\cline{2-12}
 & {\scriptsize\textbf{Time(s)}} & {\scriptsize 2.934} & {\scriptsize 2.056} & {\scriptsize 1.153} & {\scriptsize 0.489} & {\scriptsize 0.891} & {\scriptsize 0.469} & {\scriptsize 3.178} & {\scriptsize 0.836} & {\scriptsize 10.835} & {\scriptsize 6.513}\tabularnewline
\hline 
\multirow{5}{*}{\begin{cellvarwidth}[t]
\centering
{\scriptsize\textbf{levy}}{\scriptsize\par}

{\scriptsize\textbf{{[}d=100{]}}}
\end{cellvarwidth}} & {\scriptsize\textbf{Best}} & {\scriptsize 0.805754} & {\scriptsize 0.0895283} & {\scriptsize 0.358113} & {\scriptsize 2.16204} & {\scriptsize 0.626698} & {\scriptsize 2.68592} & {\scriptsize 2.97448} & {\scriptsize 0.17911} & {\scriptsize 1.70104} & {\scriptsize 0.268585}\tabularnewline
\cline{2-12}
 & {\scriptsize\textbf{Mean}} & {\scriptsize 2.14332} & {\scriptsize 0.382655} & {\scriptsize 1.2195} & {\scriptsize 8.27654} & {\scriptsize 1.28239} & {\scriptsize 3.67756} & {\scriptsize 5.40488} & {\scriptsize 1.26403} & {\scriptsize 3.22232} & {\scriptsize 1.11193}\tabularnewline
\cline{2-12}
 & {\scriptsize\textbf{Rate\%}} & {\scriptsize 7\%} & {\scriptsize 17\%} & {\scriptsize 3\%} & {\scriptsize 3\%} & {\scriptsize 7\%} & {\scriptsize 3\%} & {\scriptsize 7\%} & {\scriptsize 3\%} & {\scriptsize 3\%} & {\scriptsize 3\%}\tabularnewline
\cline{2-12}
 & {\scriptsize\textbf{SD}} & {\scriptsize 0.8117} & {\scriptsize 0.244084} & {\scriptsize 0.535975} & {\scriptsize 2.04575} & {\scriptsize 0.564979} & {\scriptsize 0.817505} & {\scriptsize 2.09253} & {\scriptsize 0.598227} & {\scriptsize 1.31638} & {\scriptsize 0.457073}\tabularnewline
\cline{2-12}
 & {\scriptsize\textbf{Time(s)}} & {\scriptsize 0.397} & {\scriptsize 6.053} & {\scriptsize 0.161} & {\scriptsize 0.112} & {\scriptsize 0.182} & {\scriptsize 0.161} & {\scriptsize 0.408} & {\scriptsize 0.165} & {\scriptsize 1.052} & {\scriptsize 0.823}\tabularnewline
\hline 
\multirow{5}{*}{\begin{cellvarwidth}[t]
\centering
{\scriptsize\textbf{messenger}}{\scriptsize\par}

{\scriptsize\textbf{{[}d=14{]}}}
\end{cellvarwidth}} & {\scriptsize\textbf{Best}} & {\scriptsize 26.7175} & {\scriptsize 26.7175} & {\scriptsize 26.7175} & {\scriptsize 26.862} & {\scriptsize 26.7175} & {\scriptsize 26.7175} & {\scriptsize 26.7175} & {\scriptsize 26.7175} & {\scriptsize 26.7175} & {\scriptsize 26.7175}\tabularnewline
\cline{2-12}
 & {\scriptsize\textbf{Mean}} & {\scriptsize 26.7175} & {\scriptsize 26.7175} & {\scriptsize 26.7175} & {\scriptsize 27.2064} & {\scriptsize 26.7175} & {\scriptsize 26.7175} & {\scriptsize 26.7175} & {\scriptsize 26.7175} & {\scriptsize 26.7175} & {\scriptsize 26.7175}\tabularnewline
\cline{2-12}
 & {\scriptsize\textbf{Rate\%}} & {\scriptsize 100\%} & {\scriptsize 100\%} & {\scriptsize 100\%} & {\scriptsize 3\%} & {\scriptsize 100\%} & {\scriptsize 100\%} & {\scriptsize 100\%} & {\scriptsize 100\%} & {\scriptsize 100\%} & {\scriptsize 100\%}\tabularnewline
\cline{2-12}
 & {\scriptsize\textbf{SD}} & {\scriptsize 7.2269e-15} & {\scriptsize 7.2269e-15} & {\scriptsize 7.2269e-15} & {\scriptsize 0.220782} & {\scriptsize 7.2269e-15} & {\scriptsize 4.0769e-13} & {\scriptsize 7.2269e-15} & {\scriptsize 7.2269e-15} & {\scriptsize 7.2269e-15} & {\scriptsize 7.2269e-15}\tabularnewline
\cline{2-12}
 & {\scriptsize\textbf{Time(s)}} & {\scriptsize 0.126} & {\scriptsize 0.147} & {\scriptsize 0.163} & {\scriptsize 0.176} & {\scriptsize 0.159} & {\scriptsize 0.169} & {\scriptsize 0.13} & {\scriptsize 0.168} & {\scriptsize 0.207} & {\scriptsize 0.303}\tabularnewline
\hline 
\multirow{5}{*}{\begin{cellvarwidth}[t]
\centering
{\scriptsize\textbf{michalewicz}}{\scriptsize\par}

{\scriptsize\textbf{{[}d=100{]}}}
\end{cellvarwidth}} & {\scriptsize\textbf{Best}} & {\scriptsize -96.0326} & {\scriptsize -77.4438} & {\scriptsize -93.8523} & {\scriptsize -79.5563} & {\scriptsize -95.7084} & {\scriptsize -99.2415} & {\scriptsize -79.5299} & {\scriptsize -98.6917} & {\scriptsize -65.359} & {\scriptsize -43.2703}\tabularnewline
\cline{2-12}
 & {\scriptsize\textbf{Mean}} & {\scriptsize -93.0301} & {\scriptsize -69.5914} & {\scriptsize -91.3666} & {\scriptsize -72.7176} & {\scriptsize -93.5333} & {\scriptsize -99.1589} & {\scriptsize -77.7664} & {\scriptsize -97.0668} & {\scriptsize -59.528} & {\scriptsize -39.8999}\tabularnewline
\cline{2-12}
 & {\scriptsize\textbf{Rate\%}} & {\scriptsize 3\%} & {\scriptsize 3\%} & {\scriptsize 3\%} & {\scriptsize 3\%} & {\scriptsize 3\%} & {\scriptsize 3\%} & {\scriptsize 3\%} & {\scriptsize 3\%} & {\scriptsize 3\%} & {\scriptsize 3\%}\tabularnewline
\cline{2-12}
 & {\scriptsize\textbf{SD}} & {\scriptsize 3.12945} & {\scriptsize 3.34186} & {\scriptsize 1.04978} & {\scriptsize 4.43642} & {\scriptsize 1.13955} & {\scriptsize 0.0714478} & {\scriptsize 0.795036} & {\scriptsize 0.964969} & {\scriptsize 2.01878} & {\scriptsize 1.40427}\tabularnewline
\cline{2-12}
 & {\scriptsize\textbf{Time(s)}} & {\scriptsize 0.783} & {\scriptsize 18.223} & {\scriptsize 0.301} & {\scriptsize 0.174} & {\scriptsize 0.29} & {\scriptsize 0.186} & {\scriptsize 0.815} & {\scriptsize 0.227} & {\scriptsize 2.615} & {\scriptsize 1.863}\tabularnewline
\hline 
\multirow{5}{*}{\begin{cellvarwidth}[t]
\centering
{\scriptsize\textbf{polyphase}}{\scriptsize\par}

{\scriptsize\textbf{{[}d=20{]} \citep{cec2011}}}
\end{cellvarwidth}} & {\scriptsize\textbf{Best}} & {\scriptsize 0.0290567} & {\scriptsize 0.0201421} & {\scriptsize 0.00571762} & {\scriptsize 0.284493} & {\scriptsize 0.023238} & {\scriptsize 0.108972} & {\scriptsize 0.314647} & {\scriptsize 0.000621833} & {\scriptsize 0.0475955} & {\scriptsize 0.0298618}\tabularnewline
\cline{2-12}
 & {\scriptsize\textbf{Mean}} & {\scriptsize 0.294367} & {\scriptsize 0.128674} & {\scriptsize 0.2011} & {\scriptsize 0.807162} & {\scriptsize 0.185771} & {\scriptsize 0.342068} & {\scriptsize 0.5693} & {\scriptsize 0.101974} & {\scriptsize 0.298147} & {\scriptsize 0.254132}\tabularnewline
\cline{2-12}
 & {\scriptsize\textbf{Rate\%}} & {\scriptsize 3\%} & {\scriptsize 3\%} & {\scriptsize 3\%} & {\scriptsize 3\%} & {\scriptsize 3\%} & {\scriptsize 3\%} & {\scriptsize 3\%} & {\scriptsize 3\%} & {\scriptsize 3\%} & {\scriptsize 3\%}\tabularnewline
\cline{2-12}
 & {\scriptsize\textbf{SD}} & {\scriptsize 0.114409} & {\scriptsize 0.100273} & {\scriptsize 0.12277} & {\scriptsize 0.293895} & {\scriptsize 0.127557} & {\scriptsize 0.0817433} & {\scriptsize 0.101133} & {\scriptsize 0.0639291} & {\scriptsize 0.164167} & {\scriptsize 0.168076}\tabularnewline
\cline{2-12}
 & {\scriptsize\textbf{Time(s)}} & {\scriptsize 0.313} & {\scriptsize 0.269} & {\scriptsize 0.182} & {\scriptsize 0.161} & {\scriptsize 0.189} & {\scriptsize 0.165} & {\scriptsize 0.473} & {\scriptsize 0.189} & {\scriptsize 0.946} & {\scriptsize 0.832}\tabularnewline
\hline 
\multirow{5}{*}{\begin{cellvarwidth}[t]
\centering
{\scriptsize\textbf{potential}}{\scriptsize\par}

{\scriptsize\textbf{{[}d=38{]}}}
\end{cellvarwidth}} & {\scriptsize\textbf{Best}} & {\scriptsize -131.443} & {\scriptsize -155.503} & {\scriptsize -160.182} & {\scriptsize -168.198} & {\scriptsize -149.112} & {\scriptsize -157.276} & {\scriptsize -110.769} & {\scriptsize -167.054} & {\scriptsize -127.985} & {\scriptsize -73.2198}\tabularnewline
\cline{2-12}
 & {\scriptsize\textbf{Mean}} & {\scriptsize -116.772} & {\scriptsize -142.71} & {\scriptsize -149.977} & {\scriptsize 361.377} & {\scriptsize -110.74} & {\scriptsize -153.238} & {\scriptsize -105.833} & {\scriptsize -160.986} & {\scriptsize -101.665} & {\scriptsize 39.4225}\tabularnewline
\cline{2-12}
 & {\scriptsize\textbf{Rate\%}} & {\scriptsize 3\%} & {\scriptsize 3\%} & {\scriptsize 3\%} & {\scriptsize 3\%} & {\scriptsize 3\%} & {\scriptsize 3\%} & {\scriptsize 3\%} & {\scriptsize 3\%} & {\scriptsize 3\%} & {\scriptsize 3\%}\tabularnewline
\cline{2-12}
 & {\scriptsize\textbf{SD}} & {\scriptsize 6.09232} & {\scriptsize 7.18836} & {\scriptsize 4.25891} & {\scriptsize 620.768} & {\scriptsize 27.1126} & {\scriptsize 2.07754} & {\scriptsize 2.85937} & {\scriptsize 3.71072} & {\scriptsize 13.6216} & {\scriptsize 84.9164}\tabularnewline
\cline{2-12}
 & {\scriptsize\textbf{Time(s)}} & {\scriptsize 0.754} & {\scriptsize 0.459} & {\scriptsize 0.824} & {\scriptsize 0.173} & {\scriptsize 0.232} & {\scriptsize 0.18} & {\scriptsize 0.637} & {\scriptsize 0.197} & {\scriptsize 2.032} & {\scriptsize 1.383}\tabularnewline
\hline 
\multirow{5}{*}{\begin{cellvarwidth}[t]
\centering
{\scriptsize\textbf{powell}}{\scriptsize\par}

{\scriptsize\textbf{{[}d=100{]}}}
\end{cellvarwidth}} & {\scriptsize\textbf{Best}} & {\scriptsize 0.000306326} & {\scriptsize 0.000406189} & {\scriptsize 8.8867e-05} & {\scriptsize 2.2919e-06} & {\scriptsize 0.000653703} & {\scriptsize 0.0703045} & {\scriptsize 0.0208072} & {\scriptsize 0} & {\scriptsize 0.0745429} & {\scriptsize 0.0121888}\tabularnewline
\cline{2-12}
 & {\scriptsize\textbf{Mean}} & {\scriptsize 0.00210306} & {\scriptsize 0.00112813} & {\scriptsize 0.000183153} & {\scriptsize 1.1710e-05} & {\scriptsize 0.00138046} & {\scriptsize 0.307694} & {\scriptsize 0.0387934} & {\scriptsize 0} & {\scriptsize 0.204422} & {\scriptsize 0.0183201}\tabularnewline
\cline{2-12}
 & {\scriptsize\textbf{Rate\%}} & {\scriptsize 3\%} & {\scriptsize 3\%} & {\scriptsize 3\%} & {\scriptsize 3\%} & {\scriptsize 3\%} & {\scriptsize 3\%} & {\scriptsize 3\%} & {\scriptsize 100\%} & {\scriptsize 3\%} & {\scriptsize 3\%}\tabularnewline
\cline{2-12}
 & {\scriptsize\textbf{SD}} & {\scriptsize 0.00330406} & {\scriptsize 0.000392613} & {\scriptsize 6.6716e-05} & {\scriptsize 5.3207e-06} & {\scriptsize 0.000354841} & {\scriptsize 0.120565} & {\scriptsize 0.00945123} & {\scriptsize 0} & {\scriptsize 0.0953043} & {\scriptsize 0.00387703}\tabularnewline
\cline{2-12}
 & {\scriptsize\textbf{Time(s)}} & {\scriptsize 0.273} & {\scriptsize 12.152} & {\scriptsize 0.153} & {\scriptsize 0.096} & {\scriptsize 0.166} & {\scriptsize 0.165} & {\scriptsize 0.282} & {\scriptsize 0.158} & {\scriptsize 0.589} & {\scriptsize 0.554}\tabularnewline
\hline 
\multirow{5}{*}{\begin{cellvarwidth}[t]
\centering
{\scriptsize\textbf{rosenbrock}}{\scriptsize\par}

{\scriptsize\textbf{{[}d=50{]}}}
\end{cellvarwidth}} & {\scriptsize\textbf{Best}} & {\scriptsize 4.29661} & {\scriptsize 17.7119} & {\scriptsize 0.0552813} & {\scriptsize 8.81167} & {\scriptsize 12.7153} & {\scriptsize 38.3447} & {\scriptsize 19.14} & {\scriptsize 14.951} & {\scriptsize 36.2386} & {\scriptsize 29.2749}\tabularnewline
\cline{2-12}
 & {\scriptsize\textbf{Mean}} & {\scriptsize 13.6938} & {\scriptsize 21.4993} & {\scriptsize 2.46021} & {\scriptsize 12.4255} & {\scriptsize 16.8552} & {\scriptsize 40.5928} & {\scriptsize 23.1997} & {\scriptsize 24.7064} & {\scriptsize 37.3863} & {\scriptsize 31.35}\tabularnewline
\cline{2-12}
 & {\scriptsize\textbf{Rate\%}} & {\scriptsize 3\%} & {\scriptsize 3\%} & {\scriptsize 3\%} & {\scriptsize 3\%} & {\scriptsize 3\%} & {\scriptsize 3\%} & {\scriptsize 3\%} & {\scriptsize 3\%} & {\scriptsize 3\%} & {\scriptsize 3\%}\tabularnewline
\cline{2-12}
 & {\scriptsize\textbf{SD}} & {\scriptsize 6.40239} & {\scriptsize 1.87024} & {\scriptsize 2.04484} & {\scriptsize 2.48772} & {\scriptsize 2.23124} & {\scriptsize 1.14897} & {\scriptsize 2.82364} & {\scriptsize 3.1607} & {\scriptsize 0.608985} & {\scriptsize 0.852551}\tabularnewline
\cline{2-12}
 & {\scriptsize\textbf{Time(s)}} & {\scriptsize 0.19} & {\scriptsize 1.831} & {\scriptsize 0.153} & {\scriptsize 0.157} & {\scriptsize 0.16} & {\scriptsize 0.164} & {\scriptsize 0.174} & {\scriptsize 0.235} & {\scriptsize 0.452} & {\scriptsize 0.404}\tabularnewline
\hline 
\multirow{5}{*}{\begin{cellvarwidth}[t]
\centering
{\scriptsize\textbf{rosetta}}{\scriptsize\par}

{\scriptsize\textbf{{[}d=22{]}}}
\end{cellvarwidth}} & {\scriptsize\textbf{Best}} & {\scriptsize 3.36872} & {\scriptsize 3.36872} & {\scriptsize 3.72444} & {\scriptsize 3.81703} & {\scriptsize 3.36872} & {\scriptsize 3.72444} & {\scriptsize 3.36872} & {\scriptsize 3.36872} & {\scriptsize 3.36872} & {\scriptsize 3.36872}\tabularnewline
\cline{2-12}
 & {\scriptsize\textbf{Mean}} & {\scriptsize 3.67701} & {\scriptsize 3.71259} & {\scriptsize 3.72444} & {\scriptsize 5.13858} & {\scriptsize 3.45172} & {\scriptsize 3.72444} & {\scriptsize 3.36872} & {\scriptsize 3.36872} & {\scriptsize 3.49915} & {\scriptsize 3.48729}\tabularnewline
\cline{2-12}
 & {\scriptsize\textbf{Rate\%}} & {\scriptsize 13\%} & {\scriptsize 3\%} & {\scriptsize 100\%} & {\scriptsize 3\%} & {\scriptsize 77\%} & {\scriptsize 80\%} & {\scriptsize 100\%} & {\scriptsize 100\%} & {\scriptsize 63\%} & {\scriptsize 67\%}\tabularnewline
\cline{2-12}
 & {\scriptsize\textbf{SD}} & {\scriptsize 0.122991} & {\scriptsize 0.0649465} & {\scriptsize 1.3550e-15} & {\scriptsize 0.77485} & {\scriptsize 0.153028} & {\scriptsize 4.3759e-08} & {\scriptsize 9.0336e-16} & {\scriptsize 9.0336e-16} & {\scriptsize 0.174353} & {\scriptsize 0.170558}\tabularnewline
\cline{2-12}
 & {\scriptsize\textbf{Time(s)}} & {\scriptsize 0.153} & {\scriptsize 0.174} & {\scriptsize 0.17} & {\scriptsize 0.172} & {\scriptsize 0.163} & {\scriptsize 0.167} & {\scriptsize 0.152} & {\scriptsize 0.174} & {\scriptsize 0.269} & {\scriptsize 0.348}\tabularnewline
\hline 
\multirow{5}{*}{\begin{cellvarwidth}[t]
\centering
{\scriptsize\textbf{rotatedrosenbrock }}{\scriptsize\par}

{\scriptsize\textbf{{[}d=50{]}}}
\end{cellvarwidth}} & {\scriptsize\textbf{Best}} & {\scriptsize 0.00604507} & {\scriptsize 17.1324} & {\scriptsize 6.8860e-09} & {\scriptsize 1.9000e-11} & {\scriptsize 1.7355e-07} & {\scriptsize 0.021825} & {\scriptsize 0.000383499} & {\scriptsize 0.000444079} & {\scriptsize 0.00248885} & {\scriptsize 0.037282}\tabularnewline
\cline{2-12}
 & {\scriptsize\textbf{Mean}} & {\scriptsize 18.4717} & {\scriptsize 24.1876} & {\scriptsize 3.76344} & {\scriptsize 9.01549} & {\scriptsize 19.4448} & {\scriptsize 18.0046} & {\scriptsize 42.9691} & {\scriptsize 27.671} & {\scriptsize 27.6186} & {\scriptsize 43.6173}\tabularnewline
\cline{2-12}
 & {\scriptsize\textbf{Rate\%}} & {\scriptsize 3\%} & {\scriptsize 3\%} & {\scriptsize 3\%} & {\scriptsize 3\%} & {\scriptsize 3\%} & {\scriptsize 3\%} & {\scriptsize 3\%} & {\scriptsize 3\%} & {\scriptsize 3\%} & {\scriptsize 3\%}\tabularnewline
\cline{2-12}
 & {\scriptsize\textbf{SD}} & {\scriptsize 17.4377} & {\scriptsize 10.0064} & {\scriptsize 3.00886} & {\scriptsize 5.33793} & {\scriptsize 23.5759} & {\scriptsize 27.9477} & {\scriptsize 42.3238} & {\scriptsize 15.0013} & {\scriptsize 33.354} & {\scriptsize 33.4042}\tabularnewline
\cline{2-12}
 & {\scriptsize\textbf{Time(s)}} & {\scriptsize 0.193} & {\scriptsize 1.9} & {\scriptsize 0.151} & {\scriptsize 0.145} & {\scriptsize 0.16} & {\scriptsize 0.168} & {\scriptsize 0.203} & {\scriptsize 0.237} & {\scriptsize 0.527} & {\scriptsize 0.418}\tabularnewline
\hline 
\multirow{5}{*}{\begin{cellvarwidth}[t]
\centering
{\scriptsize\textbf{salomon}}{\scriptsize\par}

{\scriptsize\textbf{{[}d=100{]}}}
\end{cellvarwidth}} & {\scriptsize\textbf{Best}} & {\scriptsize 0.499873} & {\scriptsize 0.799873} & {\scriptsize 0.699873} & {\scriptsize 0} & {\scriptsize 0.899873} & {\scriptsize 3.49987} & {\scriptsize 0.499873} & {\scriptsize 0} & {\scriptsize 1.19987} & {\scriptsize 0.699873}\tabularnewline
\cline{2-12}
 & {\scriptsize\textbf{Mean}} & {\scriptsize 1.16987} & {\scriptsize 1.05321} & {\scriptsize 0.959873} & {\scriptsize 0} & {\scriptsize 1.18321} & {\scriptsize 4.58654} & {\scriptsize 0.86654} & {\scriptsize 0} & {\scriptsize 1.67654} & {\scriptsize 0.859874}\tabularnewline
\cline{2-12}
 & {\scriptsize\textbf{Rate\%}} & {\scriptsize 3\%} & {\scriptsize 7\%} & {\scriptsize 3\%} & {\scriptsize 100\%} & {\scriptsize 7\%} & {\scriptsize 3\%} & {\scriptsize 3\%} & {\scriptsize 100\%} & {\scriptsize 3\%} & {\scriptsize 3\%}\tabularnewline
\cline{2-12}
 & {\scriptsize\textbf{SD}} & {\scriptsize 0.313105} & {\scriptsize 0.127937} & {\scriptsize 0.132873} & {\scriptsize 0} & {\scriptsize 0.153316} & {\scriptsize 0.644196} & {\scriptsize 0.250975} & {\scriptsize 0} & {\scriptsize 0.286095} & {\scriptsize 0.106996}\tabularnewline
\cline{2-12}
 & {\scriptsize\textbf{Time(s)}} & {\scriptsize 0.361} & {\scriptsize 13.864} & {\scriptsize 0.147} & {\scriptsize 0.094} & {\scriptsize 0.169} & {\scriptsize 0.164} & {\scriptsize 0.283} & {\scriptsize 0.165} & {\scriptsize 0.574} & {\scriptsize 0.566}\tabularnewline
\hline 
\multirow{5}{*}{\begin{cellvarwidth}[t]
\centering
{\scriptsize\textbf{schwefel}}{\scriptsize\par}

{\scriptsize\textbf{{[}d=24{]}}}
\end{cellvarwidth}} & {\scriptsize\textbf{Best}} & {\scriptsize 0.000305462} & {\scriptsize 0.000305462} & {\scriptsize 0.000305462} & {\scriptsize 3552.85} & {\scriptsize 0.000305462} & {\scriptsize 0.000305462} & {\scriptsize 0.000305462} & {\scriptsize 0.000305462} & {\scriptsize 0.000305462} & {\scriptsize 0.00207051}\tabularnewline
\cline{2-12}
 & {\scriptsize\textbf{Mean}} & {\scriptsize 0.000305462} & {\scriptsize 0.000305462} & {\scriptsize 129.625} & {\scriptsize 4825.87} & {\scriptsize 51.3236} & {\scriptsize 43.4277} & {\scriptsize 0.000305462} & {\scriptsize 0.000305462} & {\scriptsize 604.44} & {\scriptsize 2970.62}\tabularnewline
\cline{2-12}
 & {\scriptsize\textbf{Rate\%}} & {\scriptsize 53\%} & {\scriptsize 100\%} & {\scriptsize 10\%} & {\scriptsize 3\%} & {\scriptsize 63\%} & {\scriptsize 63\%} & {\scriptsize 77\%} & {\scriptsize 90\%} & {\scriptsize 3\%} & {\scriptsize 3\%}\tabularnewline
\cline{2-12}
 & {\scriptsize\textbf{SD}} & {\scriptsize 1.2576e-12} & {\scriptsize 0} & {\scriptsize 192.272} & {\scriptsize 681.49} & {\scriptsize 74.1498} & {\scriptsize 79.1968} & {\scriptsize 8.6037e-13} & {\scriptsize 6.1026e-13} & {\scriptsize 346.058} & {\scriptsize 780.031}\tabularnewline
\cline{2-12}
 & {\scriptsize\textbf{Time(s)}} & {\scriptsize 0.198} & {\scriptsize 0.24} & {\scriptsize 0.16} & {\scriptsize 0.17} & {\scriptsize 0.16} & {\scriptsize 0.169} & {\scriptsize 0.156} & {\scriptsize 0.168} & {\scriptsize 0.396} & {\scriptsize 0.437}\tabularnewline
\hline 
\multirow{5}{*}{\begin{cellvarwidth}[t]
\centering
{\scriptsize\textbf{sinusoidal}}{\scriptsize\par}

{\scriptsize\textbf{{[}d=200{]}}}
\end{cellvarwidth}} & {\scriptsize\textbf{Best}} & {\scriptsize -3.5} & {\scriptsize -3.5} & {\scriptsize -3.5} & {\scriptsize -3.5} & {\scriptsize -3.49999} & {\scriptsize -2.94397} & {\scriptsize -3.49938} & {\scriptsize -3.5} & {\scriptsize -3.34223} & {\scriptsize -3.49888}\tabularnewline
\cline{2-12}
 & {\scriptsize\textbf{Mean}} & {\scriptsize -3.49932} & {\scriptsize -3.49999} & {\scriptsize -3.5} & {\scriptsize -1.01881} & {\scriptsize -3.49993} & {\scriptsize -2.24821} & {\scriptsize -3.46135} & {\scriptsize -3.5} & {\scriptsize -3.11639} & {\scriptsize -3.49466}\tabularnewline
\cline{2-12}
 & {\scriptsize\textbf{Rate\%}} & {\scriptsize 3\%} & {\scriptsize 3\%} & {\scriptsize 7\%} & {\scriptsize 20\%} & {\scriptsize 3\%} & {\scriptsize 3\%} & {\scriptsize 3\%} & {\scriptsize 100\%} & {\scriptsize 3\%} & {\scriptsize 3\%}\tabularnewline
\cline{2-12}
 & {\scriptsize\textbf{SD}} & {\scriptsize 0.00204778} & {\scriptsize 6.7328e-06} & {\scriptsize 4.3683e-06} & {\scriptsize 1.33665} & {\scriptsize 8.0365e-05} & {\scriptsize 0.290855} & {\scriptsize 0.148488} & {\scriptsize 0} & {\scriptsize 0.44314} & {\scriptsize 0.00396205}\tabularnewline
\cline{2-12}
 & {\scriptsize\textbf{Time(s)}} & {\scriptsize 0.682} & {\scriptsize 0.333} & {\scriptsize 0.256} & {\scriptsize 0.261} & {\scriptsize 0.268} & {\scriptsize 0.19} & {\scriptsize 0.84} & {\scriptsize 0.234} & {\scriptsize 2.355} & {\scriptsize 1.627}\tabularnewline
\hline 
\multirow{5}{*}{\begin{cellvarwidth}[t]
\centering
{\scriptsize\textbf{stirredtankreactor}}{\scriptsize\par}

{\scriptsize\textbf{{[}d=20{]} \citep{cec2011}}}
\end{cellvarwidth}} & {\scriptsize\textbf{Best}} & {\scriptsize 0.534576} & {\scriptsize 0.534576} & {\scriptsize 0.534576} & {\scriptsize 0.578457} & {\scriptsize 0.534576} & {\scriptsize 0.534576} & {\scriptsize 0.534576} & {\scriptsize 0.534576} & {\scriptsize 0.534576} & {\scriptsize 0.534576}\tabularnewline
\cline{2-12}
 & {\scriptsize\textbf{Mean}} & {\scriptsize 0.534576} & {\scriptsize 0.534576} & {\scriptsize 0.534576} & {\scriptsize 0.655901} & {\scriptsize 0.534576} & {\scriptsize 0.534576} & {\scriptsize 0.534576} & {\scriptsize 0.534576} & {\scriptsize 0.534576} & {\scriptsize 0.534576}\tabularnewline
\cline{2-12}
 & {\scriptsize\textbf{Rate\%}} & {\scriptsize 100\%} & {\scriptsize 100\%} & {\scriptsize 100\%} & {\scriptsize 3\%} & {\scriptsize 100\%} & {\scriptsize 100\%} & {\scriptsize 100\%} & {\scriptsize 100\%} & {\scriptsize 100\%} & {\scriptsize 100\%}\tabularnewline
\cline{2-12}
 & {\scriptsize\textbf{SD}} & {\scriptsize 2.2584e-16} & {\scriptsize 2.2584e-16} & {\scriptsize 2.2584e-16} & {\scriptsize 0.0439165} & {\scriptsize 2.2584e-16} & {\scriptsize 2.2584e-16} & {\scriptsize 2.2584e-16} & {\scriptsize 2.2584e-16} & {\scriptsize 2.2584e-16} & {\scriptsize 2.2584e-16}\tabularnewline
\cline{2-12}
 & {\scriptsize\textbf{Time(s)}} & {\scriptsize 0.151} & {\scriptsize 0.167} & {\scriptsize 0.149} & {\scriptsize 0.171} & {\scriptsize 0.161} & {\scriptsize 0.166} & {\scriptsize 0.179} & {\scriptsize 0.174} & {\scriptsize 0.299} & {\scriptsize 0.36}\tabularnewline
\hline 
\multirow{5}{*}{\begin{cellvarwidth}[t]
\centering
{\scriptsize\textbf{tandem}}{\scriptsize\par}

{\scriptsize\textbf{{[}d=18{]}}}
\end{cellvarwidth}} & {\scriptsize\textbf{Best}} & {\scriptsize 27.6131} & {\scriptsize 27.6131} & {\scriptsize 27.6131} & {\scriptsize 27.9359} & {\scriptsize 27.6131} & {\scriptsize 27.6131} & {\scriptsize 27.6131} & {\scriptsize 27.6131} & {\scriptsize 27.6131} & {\scriptsize 27.6131}\tabularnewline
\cline{2-12}
 & {\scriptsize\textbf{Mean}} & {\scriptsize 27.6131} & {\scriptsize 27.6131} & {\scriptsize 27.6131} & {\scriptsize 28.5521} & {\scriptsize 27.6131} & {\scriptsize 27.6136} & {\scriptsize 27.6131} & {\scriptsize 27.6131} & {\scriptsize 27.6131} & {\scriptsize 27.6131}\tabularnewline
\cline{2-12}
 & {\scriptsize\textbf{Rate\%}} & {\scriptsize 100\%} & {\scriptsize 100\%} & {\scriptsize 100\%} & {\scriptsize 3\%} & {\scriptsize 100\%} & {\scriptsize 3\%} & {\scriptsize 3\%} & {\scriptsize 100\%} & {\scriptsize 100\%} & {\scriptsize 100\%}\tabularnewline
\cline{2-12}
 & {\scriptsize\textbf{SD}} & {\scriptsize 7.2269e-15} & {\scriptsize 7.2269e-15} & {\scriptsize 7.2269e-15} & {\scriptsize 0.384512} & {\scriptsize 7.2269e-15} & {\scriptsize 0.000849244} & {\scriptsize 1.4502e-05} & {\scriptsize 7.2269e-15} & {\scriptsize 2.5431e-13} & {\scriptsize 7.2269e-15}\tabularnewline
\cline{2-12}
 & {\scriptsize\textbf{Time(s)}} & {\scriptsize 0.133} & {\scriptsize 0.16} & {\scriptsize 0.159} & {\scriptsize 0.172} & {\scriptsize 0.161} & {\scriptsize 0.168} & {\scriptsize 0.129} & {\scriptsize 0.17} & {\scriptsize 0.248} & {\scriptsize 0.321}\tabularnewline
\hline 
\multirow{5}{*}{\begin{cellvarwidth}[t]
\centering
{\scriptsize\textbf{tersoffb}}{\scriptsize\par}

{\scriptsize\textbf{{[}d=30{]}}}
\end{cellvarwidth}} & {\scriptsize\textbf{Best}} & {\scriptsize -36.7227} & {\scriptsize -35.6542} & {\scriptsize -36.4017} & {\scriptsize -33.0671} & {\scriptsize -36.0879} & {\scriptsize -35.8773} & {\scriptsize -34.4524} & {\scriptsize -36.0073} & {\scriptsize -35.4998} & {\scriptsize -36.4459}\tabularnewline
\cline{2-12}
 & {\scriptsize\textbf{Mean}} & {\scriptsize -34.8526} & {\scriptsize -34.5904} & {\scriptsize -34.7363} & {\scriptsize -27.6905} & {\scriptsize -33.99} & {\scriptsize -34.9479} & {\scriptsize -33.7635} & {\scriptsize -34.0183} & {\scriptsize -33.8424} & {\scriptsize -33.5724}\tabularnewline
\cline{2-12}
 & {\scriptsize\textbf{Rate\%}} & {\scriptsize 3\%} & {\scriptsize 3\%} & {\scriptsize 3\%} & {\scriptsize 3\%} & {\scriptsize 3\%} & {\scriptsize 3\%} & {\scriptsize 3\%} & {\scriptsize 3\%} & {\scriptsize 3\%} & {\scriptsize 3\%}\tabularnewline
\cline{2-12}
 & {\scriptsize\textbf{SD}} & {\scriptsize 0.751157} & {\scriptsize 0.855449} & {\scriptsize 0.69871} & {\scriptsize 3.85757} & {\scriptsize 1.12758} & {\scriptsize 0.391396} & {\scriptsize 0.34209} & {\scriptsize 1.37822} & {\scriptsize 1.02464} & {\scriptsize 1.22774}\tabularnewline
\cline{2-12}
 & {\scriptsize\textbf{Time(s)}} & {\scriptsize 2.013} & {\scriptsize 1.387} & {\scriptsize 0.873} & {\scriptsize 0.243} & {\scriptsize 0.682} & {\scriptsize 0.454} & {\scriptsize 2.699} & {\scriptsize 0.629} & {\scriptsize 6.956} & {\scriptsize 5.091}\tabularnewline
\hline 
\multirow{5}{*}{\begin{cellvarwidth}[t]
\centering
{\scriptsize\textbf{tersoffc}}{\scriptsize\par}

{\scriptsize\textbf{{[}d=30{]}}}
\end{cellvarwidth}} & {\scriptsize\textbf{Best}} & {\scriptsize -36.2305} & {\scriptsize -37.1606} & {\scriptsize -37.6202} & {\scriptsize -34.5494} & {\scriptsize -36.0066} & {\scriptsize -35.8429} & {\scriptsize -35.2832} & {\scriptsize -35.7001} & {\scriptsize -37.1504} & {\scriptsize -35.375}\tabularnewline
\cline{2-12}
 & {\scriptsize\textbf{Mean}} & {\scriptsize -35.1432} & {\scriptsize -35.0916} & {\scriptsize -35.1291} & {\scriptsize -28.9997} & {\scriptsize -34.4414} & {\scriptsize -35.351} & {\scriptsize -34.0384} & {\scriptsize -33.9619} & {\scriptsize -34.4759} & {\scriptsize -33.8948}\tabularnewline
\cline{2-12}
 & {\scriptsize\textbf{Rate\%}} & {\scriptsize 3\%} & {\scriptsize 3\%} & {\scriptsize 3\%} & {\scriptsize 3\%} & {\scriptsize 3\%} & {\scriptsize 3\%} & {\scriptsize 3\%} & {\scriptsize 3\%} & {\scriptsize 3\%} & {\scriptsize 3\%}\tabularnewline
\cline{2-12}
 & {\scriptsize\textbf{SD}} & {\scriptsize 0.547894} & {\scriptsize 0.791315} & {\scriptsize 0.748068} & {\scriptsize 4.33677} & {\scriptsize 0.959211} & {\scriptsize 0.253253} & {\scriptsize 0.407002} & {\scriptsize 1.02789} & {\scriptsize 1.06793} & {\scriptsize 0.748882}\tabularnewline
\cline{2-12}
 & {\scriptsize\textbf{Time(s)}} & {\scriptsize 3.449} & {\scriptsize 2.074} & {\scriptsize 1.535} & {\scriptsize 0.494} & {\scriptsize 1.258} & {\scriptsize 0.826} & {\scriptsize 4.642} & {\scriptsize 1.025} & {\scriptsize 13.451} & {\scriptsize 8.976}\tabularnewline
\hline 
\multirow{5}{*}{\begin{cellvarwidth}[t]
\centering
{\scriptsize\textbf{test2n}}{\scriptsize\par}

{\scriptsize\textbf{{[}d=200{]}}}
\end{cellvarwidth}} & {\scriptsize\textbf{Best}} & {\scriptsize -7620.59} & {\scriptsize -7725.73} & {\scriptsize -7126.4} & {\scriptsize -6829.53} & {\scriptsize -7394.99} & {\scriptsize -7733.32} & {\scriptsize -7069.81} & {\scriptsize -7833.23} & {\scriptsize -6884.53} & {\scriptsize -7209.33}\tabularnewline
\cline{2-12}
 & {\scriptsize\textbf{Mean}} & {\scriptsize -7150.57} & {\scriptsize -7533.34} & {\scriptsize -6937.91} & {\scriptsize -6641.04} & {\scriptsize -7240.39} & {\scriptsize -6914.4} & {\scriptsize -6914.21} & {\scriptsize -7833.23} & {\scriptsize -6642.31} & {\scriptsize -7016.36}\tabularnewline
\cline{2-12}
 & {\scriptsize\textbf{Rate\%}} & {\scriptsize 3\%} & {\scriptsize 3\%} & {\scriptsize 3\%} & {\scriptsize 3\%} & {\scriptsize 3\%} & {\scriptsize 3\%} & {\scriptsize 3\%} & {\scriptsize 53\%} & {\scriptsize 3\%} & {\scriptsize 3\%}\tabularnewline
\cline{2-12}
 & {\scriptsize\textbf{SD}} & {\scriptsize 263.653} & {\scriptsize 99.0168} & {\scriptsize 111.766} & {\scriptsize 104.566} & {\scriptsize 107.332} & {\scriptsize 315.232} & {\scriptsize 86.5434} & {\scriptsize 0.000515372} & {\scriptsize 109.786} & {\scriptsize 85.293}\tabularnewline
\cline{2-12}
 & {\scriptsize\textbf{Time(s)}} & {\scriptsize 0.511} & {\scriptsize 0.361} & {\scriptsize 0.227} & {\scriptsize 0.191} & {\scriptsize 0.192} & {\scriptsize 0.17} & {\scriptsize 0.473} & {\scriptsize 0.175} & {\scriptsize 1.012} & {\scriptsize 0.901}\tabularnewline
\hline 
\multirow{5}{*}{\begin{cellvarwidth}[t]
\centering
{\scriptsize\textbf{test30n}}{\scriptsize\par}

{\scriptsize\textbf{{[}d=200{]}}}
\end{cellvarwidth}} & {\scriptsize\textbf{Best}} & {\scriptsize 4.4641e-08} & {\scriptsize 1.2776e-07} & {\scriptsize 1.3586e-08} & {\scriptsize 0} & {\scriptsize 3.1773e-06} & {\scriptsize 6.35398} & {\scriptsize 0.000197792} & {\scriptsize 0} & {\scriptsize 0.874805} & {\scriptsize 0.0075017}\tabularnewline
\cline{2-12}
 & {\scriptsize\textbf{Mean}} & {\scriptsize 0.0426922} & {\scriptsize 0.00797665} & {\scriptsize 0.00872158} & {\scriptsize 0.00686355} & {\scriptsize 0.0521622} & {\scriptsize 12.0216} & {\scriptsize 0.401587} & {\scriptsize 0.0131041} & {\scriptsize 2.6814} & {\scriptsize 0.107771}\tabularnewline
\cline{2-12}
 & {\scriptsize\textbf{Rate\%}} & {\scriptsize 3\%} & {\scriptsize 3\%} & {\scriptsize 3\%} & {\scriptsize 57\%} & {\scriptsize 3\%} & {\scriptsize 3\%} & {\scriptsize 3\%} & {\scriptsize 37\%} & {\scriptsize 3\%} & {\scriptsize 3\%}\tabularnewline
\cline{2-12}
 & {\scriptsize\textbf{SD}} & {\scriptsize 0.0831271} & {\scriptsize 0.0222685} & {\scriptsize 0.0108496} & {\scriptsize 0.0100529} & {\scriptsize 0.118902} & {\scriptsize 2.39491} & {\scriptsize 0.354119} & {\scriptsize 0.0286558} & {\scriptsize 1.2027} & {\scriptsize 0.116349}\tabularnewline
\cline{2-12}
 & {\scriptsize\textbf{Time(s)}} & {\scriptsize 0.61} & {\scriptsize 0.265} & {\scriptsize 0.211} & {\scriptsize 0.223} & {\scriptsize 0.235} & {\scriptsize 0.174} & {\scriptsize 0.69} & {\scriptsize 0.216} & {\scriptsize 1.852} & {\scriptsize 1.359}\tabularnewline
\hline 
\multirow{5}{*}{\begin{cellvarwidth}[t]
\centering
{\scriptsize\textbf{tnep}}{\scriptsize\par}

{\scriptsize\textbf{{[}d=7{]} \citep{cec2011}}}
\end{cellvarwidth}} & {\scriptsize\textbf{Best}} & {\scriptsize 21} & {\scriptsize 21} & {\scriptsize 21} & {\scriptsize 23} & {\scriptsize 21} & {\scriptsize 21} & {\scriptsize 21} & {\scriptsize 21} & {\scriptsize 21} & {\scriptsize 21}\tabularnewline
\cline{2-12}
 & {\scriptsize\textbf{Mean}} & {\scriptsize 21} & {\scriptsize 21} & {\scriptsize 21} & {\scriptsize 27.3333} & {\scriptsize 21} & {\scriptsize 21.5667} & {\scriptsize 21.1} & {\scriptsize 21} & {\scriptsize 21.1667} & {\scriptsize 21}\tabularnewline
\cline{2-12}
 & {\scriptsize\textbf{Rate\%}} & {\scriptsize 100\%} & {\scriptsize 100\%} & {\scriptsize 100\%} & {\scriptsize 13\%} & {\scriptsize 100\%} & {\scriptsize 67\%} & {\scriptsize 90\%} & {\scriptsize 100\%} & {\scriptsize 83\%} & {\scriptsize 100\%}\tabularnewline
\cline{2-12}
 & {\scriptsize\textbf{SD}} & {\scriptsize 0} & {\scriptsize 0} & {\scriptsize 0} & {\scriptsize 3.16591} & {\scriptsize 0} & {\scriptsize 1.10433} & {\scriptsize 0.305129} & {\scriptsize 0} & {\scriptsize 0.379049} & {\scriptsize 0}\tabularnewline
\cline{2-12}
 & {\scriptsize\textbf{Time(s)}} & {\scriptsize 0.204} & {\scriptsize 0.136} & {\scriptsize 0.138} & {\scriptsize 0.169} & {\scriptsize 0.153} & {\scriptsize 0.158} & {\scriptsize 0.254} & {\scriptsize 0.162} & {\scriptsize 0.467} & {\scriptsize 0.527}\tabularnewline
\hline 
\multirow{5}{*}{\begin{cellvarwidth}[t]
\centering
{\scriptsize\textbf{transmissionpricing}}{\scriptsize\par}

{\scriptsize\textbf{{[}d=126{]} \citep{cec2011}}}
\end{cellvarwidth}} & {\scriptsize\textbf{Best}} & {\scriptsize 4.53605} & {\scriptsize 4.53605} & {\scriptsize 4.53605} & {\scriptsize 4.54397} & {\scriptsize 4.53606} & {\scriptsize 4.58931} & {\scriptsize 4.53605} & {\scriptsize 4.53607} & {\scriptsize 4.54869} & {\scriptsize 4.5466}\tabularnewline
\cline{2-12}
 & {\scriptsize\textbf{Mean}} & {\scriptsize 4.54089} & {\scriptsize 4.53903} & {\scriptsize 4.53606} & {\scriptsize 4.71216} & {\scriptsize 4.53784} & {\scriptsize 4.83932} & {\scriptsize 4.60599} & {\scriptsize 4.55727} & {\scriptsize 4.62265} & {\scriptsize 4.6025}\tabularnewline
\cline{2-12}
 & {\scriptsize\textbf{Rate\%}} & {\scriptsize 10\%} & {\scriptsize 3\%} & {\scriptsize 3\%} & {\scriptsize 3\%} & {\scriptsize 3\%} & {\scriptsize 3\%} & {\scriptsize 3\%} & {\scriptsize 3\%} & {\scriptsize 3\%} & {\scriptsize 3\%}\tabularnewline
\cline{2-12}
 & {\scriptsize\textbf{SD}} & {\scriptsize 0.00994537} & {\scriptsize 0.00455269} & {\scriptsize 1.7330e-05} & {\scriptsize 0.103526} & {\scriptsize 0.00358582} & {\scriptsize 0.560773} & {\scriptsize 0.0499372} & {\scriptsize 0.0214215} & {\scriptsize 0.0432556} & {\scriptsize 0.0453235}\tabularnewline
\cline{2-12}
 & {\scriptsize\textbf{Time(s)}} & {\scriptsize 3.649} & {\scriptsize 2.127} & {\scriptsize 1.538} & {\scriptsize 1.014} & {\scriptsize 1.536} & {\scriptsize 0.773} & {\scriptsize 5.272} & {\scriptsize 1.001} & {\scriptsize 17.902} & {\scriptsize 10.47}\tabularnewline
\hline 
\multirow{5}{*}{\begin{cellvarwidth}[t]
\centering
{\scriptsize\textbf{trid}}{\scriptsize\par}

{\scriptsize\textbf{{[}d=40{]}}}
\end{cellvarwidth}} & {\scriptsize\textbf{Best}} & {\scriptsize -11440} & {\scriptsize -11440} & {\scriptsize -11440} & {\scriptsize -11440} & {\scriptsize -11373.5} & {\scriptsize -8588.73} & {\scriptsize -9720.71} & {\scriptsize -11440} & {\scriptsize -10532} & {\scriptsize -11399.9}\tabularnewline
\cline{2-12}
 & {\scriptsize\textbf{Mean}} & {\scriptsize -11331.6} & {\scriptsize -11439.6} & {\scriptsize -11439.9} & {\scriptsize -11440} & {\scriptsize -11230.3} & {\scriptsize -3989.34} & {\scriptsize -5355.66} & {\scriptsize -11399.3} & {\scriptsize -4500.91} & {\scriptsize -10543.8}\tabularnewline
\cline{2-12}
 & {\scriptsize\textbf{Rate\%}} & {\scriptsize 3\%} & {\scriptsize 3\%} & {\scriptsize 3\%} & {\scriptsize 100\%} & {\scriptsize 3\%} & {\scriptsize 3\%} & {\scriptsize 3\%} & {\scriptsize 3\%} & {\scriptsize 3\%} & {\scriptsize 3\%}\tabularnewline
\cline{2-12}
 & {\scriptsize\textbf{SD}} & {\scriptsize 289.641} & {\scriptsize 0.642353} & {\scriptsize 0.145549} & {\scriptsize 2.3668e-10} & {\scriptsize 113.783} & {\scriptsize 1964.33} & {\scriptsize 2929.99} & {\scriptsize 78.7711} & {\scriptsize 3426.05} & {\scriptsize 559.18}\tabularnewline
\cline{2-12}
 & {\scriptsize\textbf{Time(s)}} & {\scriptsize 0.167} & {\scriptsize 0.955} & {\scriptsize 0.155} & {\scriptsize 0.157} & {\scriptsize 0.159} & {\scriptsize 0.167} & {\scriptsize 0.291} & {\scriptsize 0.233} & {\scriptsize 0.347} & {\scriptsize 0.366}\tabularnewline
\hline 
\multirow{5}{*}{\begin{cellvarwidth}[t]
\centering
{\scriptsize\textbf{vibratingplatform}}{\scriptsize\par}

{\scriptsize\textbf{{[}d=5{]} \citep{vibratingplatform}}}
\end{cellvarwidth}} & {\scriptsize\textbf{Best}} & {\scriptsize 0.105406} & {\scriptsize 0.105406} & {\scriptsize 0.105406} & {\scriptsize 0.174408} & {\scriptsize 0.105406} & {\scriptsize 0.105406} & {\scriptsize 0.105406} & {\scriptsize 0.105406} & {\scriptsize 0.105406} & {\scriptsize 0.105406}\tabularnewline
\cline{2-12}
 & {\scriptsize\textbf{Mean}} & {\scriptsize 0.105406} & {\scriptsize 0.105406} & {\scriptsize 0.105406} & {\scriptsize 0.254461} & {\scriptsize 0.105406} & {\scriptsize 0.105407} & {\scriptsize 0.105409} & {\scriptsize 0.105406} & {\scriptsize 0.10541} & {\scriptsize 0.105406}\tabularnewline
\cline{2-12}
 & {\scriptsize\textbf{Rate\%}} & {\scriptsize 3\%} & {\scriptsize 3\%} & {\scriptsize 3\%} & {\scriptsize 3\%} & {\scriptsize 3\%} & {\scriptsize 3\%} & {\scriptsize 3\%} & {\scriptsize 3\%} & {\scriptsize 3\%} & {\scriptsize 3\%}\tabularnewline
\cline{2-12}
 & {\scriptsize\textbf{SD}} & {\scriptsize 3.7548e-07} & {\scriptsize 4.3623e-07} & {\scriptsize 1.2843e-07} & {\scriptsize 0.0615324} & {\scriptsize 3.5987e-07} & {\scriptsize 3.4385e-06} & {\scriptsize 5.3869e-06} & {\scriptsize 7.9169e-07} & {\scriptsize 6.1354e-06} & {\scriptsize 3.1740e-07}\tabularnewline
\cline{2-12}
 & {\scriptsize\textbf{Time(s)}} & {\scriptsize 0.106} & {\scriptsize 0.144} & {\scriptsize 0.162} & {\scriptsize 0.18} & {\scriptsize 0.157} & {\scriptsize 0.172} & {\scriptsize 0.107} & {\scriptsize 0.163} & {\scriptsize 0.245} & {\scriptsize 0.284}\tabularnewline
\hline 
\multirow{5}{*}{\begin{cellvarwidth}[t]
\centering
{\scriptsize\textbf{weatherirrigation }}{\scriptsize\par}

{\scriptsize\textbf{{[}d=30{]} \citep{weatherirrigation}}}
\end{cellvarwidth}} & {\scriptsize\textbf{Best}} & {\scriptsize 2736.36} & {\scriptsize 2736.36} & {\scriptsize 2737.8} & {\scriptsize 2743.85} & {\scriptsize 2736.36} & {\scriptsize 2751.62} & {\scriptsize 2741.22} & {\scriptsize 2736.36} & {\scriptsize 2736.36} & {\scriptsize 2736.36}\tabularnewline
\cline{2-12}
 & {\scriptsize\textbf{Mean}} & {\scriptsize 2745.94} & {\scriptsize 2740.61} & {\scriptsize 2744.38} & {\scriptsize 2783.55} & {\scriptsize 2751.4} & {\scriptsize 2817.49} & {\scriptsize 2773.37} & {\scriptsize 2741.69} & {\scriptsize 2763.76} & {\scriptsize 2750.55}\tabularnewline
\cline{2-12}
 & {\scriptsize\textbf{Rate\%}} & {\scriptsize 3\%} & {\scriptsize 20\%} & {\scriptsize 3\%} & {\scriptsize 3\%} & {\scriptsize 13\%} & {\scriptsize 3\%} & {\scriptsize 3\%} & {\scriptsize 20\%} & {\scriptsize 3\%} & {\scriptsize 3\%}\tabularnewline
\cline{2-12}
 & {\scriptsize\textbf{SD}} & {\scriptsize 8.72179} & {\scriptsize 4.04687} & {\scriptsize 5.79116} & {\scriptsize 25.4044} & {\scriptsize 14.0046} & {\scriptsize 29.1568} & {\scriptsize 17.7335} & {\scriptsize 5.56177} & {\scriptsize 20.5874} & {\scriptsize 10.5355}\tabularnewline
\cline{2-12}
 & {\scriptsize\textbf{Time(s)}} & {\scriptsize 0.459} & {\scriptsize 0.336} & {\scriptsize 0.191} & {\scriptsize 0.094} & {\scriptsize 0.214} & {\scriptsize 0.174} & {\scriptsize 0.537} & {\scriptsize 0.212} & {\scriptsize 1.616} & {\scriptsize 1.1}\tabularnewline
\hline 
\multirow{5}{*}{\begin{cellvarwidth}[t]
\centering
{\scriptsize\textbf{weierstrass}}{\scriptsize\par}

{\scriptsize\textbf{{[}d=70{]}}}
\end{cellvarwidth}} & {\scriptsize\textbf{Best}} & {\scriptsize 0.0444884} & {\scriptsize 0.00735326} & {\scriptsize 0.0308165} & {\scriptsize 0} & {\scriptsize 0.196602} & {\scriptsize 14.0741} & {\scriptsize 1.66031} & {\scriptsize 0} & {\scriptsize 2.39297} & {\scriptsize 0.166803}\tabularnewline
\cline{2-12}
 & {\scriptsize\textbf{Mean}} & {\scriptsize 3.11574} & {\scriptsize 0.2273} & {\scriptsize 0.944241} & {\scriptsize 0} & {\scriptsize 0.935228} & {\scriptsize 17.1846} & {\scriptsize 3.91315} & {\scriptsize 0} & {\scriptsize 4.75133} & {\scriptsize 0.776516}\tabularnewline
\cline{2-12}
 & {\scriptsize\textbf{Rate\%}} & {\scriptsize 3\%} & {\scriptsize 3\%} & {\scriptsize 3\%} & {\scriptsize 100\%} & {\scriptsize 3\%} & {\scriptsize 3\%} & {\scriptsize 3\%} & {\scriptsize 100\%} & {\scriptsize 3\%} & {\scriptsize 3\%}\tabularnewline
\cline{2-12}
 & {\scriptsize\textbf{SD}} & {\scriptsize 2.21603} & {\scriptsize 0.251044} & {\scriptsize 0.769111} & {\scriptsize 0} & {\scriptsize 0.505221} & {\scriptsize 1.82463} & {\scriptsize 1.14477} & {\scriptsize 0} & {\scriptsize 1.41729} & {\scriptsize 0.432902}\tabularnewline
\cline{2-12}
 & {\scriptsize\textbf{Time(s)}} & {\scriptsize 7.171} & {\scriptsize 5.336} & {\scriptsize 2.923} & {\scriptsize 1.647} & {\scriptsize 2.88} & {\scriptsize 1.505} & {\scriptsize 9.978} & {\scriptsize 2.064} & {\scriptsize 34.916} & {\scriptsize 20.088}\tabularnewline
\hline 
\multirow{5}{*}{\begin{cellvarwidth}[t]
\centering
{\scriptsize\textbf{whitley}}{\scriptsize\par}

{\scriptsize\textbf{{[}d=40{]}}}
\end{cellvarwidth}} & {\scriptsize\textbf{Best}} & {\scriptsize 170.351} & {\scriptsize 263.406} & {\scriptsize 263.305} & {\scriptsize 735.916} & {\scriptsize 508.102} & {\scriptsize 615.033} & {\scriptsize 133.182} & {\scriptsize 96.4938} & {\scriptsize 152.941} & {\scriptsize 575.902}\tabularnewline
\cline{2-12}
 & {\scriptsize\textbf{Mean}} & {\scriptsize 382.149} & {\scriptsize 414.062} & {\scriptsize 488.063} & {\scriptsize 735.916} & {\scriptsize 639.163} & {\scriptsize 1098.28} & {\scriptsize 517.969} & {\scriptsize 396.182} & {\scriptsize 593.577} & {\scriptsize 656.457}\tabularnewline
\cline{2-12}
 & {\scriptsize\textbf{Rate\%}} & {\scriptsize 3\%} & {\scriptsize 3\%} & {\scriptsize 3\%} & {\scriptsize 100\%} & {\scriptsize 3\%} & {\scriptsize 3\%} & {\scriptsize 3\%} & {\scriptsize 3\%} & {\scriptsize 3\%} & {\scriptsize 3\%}\tabularnewline
\cline{2-12}
 & {\scriptsize\textbf{SD}} & {\scriptsize 146.202} & {\scriptsize 90.2828} & {\scriptsize 129.892} & {\scriptsize 4.6252e-13} & {\scriptsize 50.325} & {\scriptsize 146.727} & {\scriptsize 166.89} & {\scriptsize 122.316} & {\scriptsize 173.816} & {\scriptsize 34.5759}\tabularnewline
\cline{2-12}
 & {\scriptsize\textbf{Time(s)}} & {\scriptsize 1.78} & {\scriptsize 1.929} & {\scriptsize 0.729} & {\scriptsize 0.339} & {\scriptsize 0.682} & {\scriptsize 0.334} & {\scriptsize 2.362} & {\scriptsize 0.6} & {\scriptsize 8.195} & {\scriptsize 4.926}\tabularnewline
\hline 
\multirow{5}{*}{\begin{cellvarwidth}[t]
\centering
{\scriptsize\textbf{wirelesscoverage}}{\scriptsize\par}

{\scriptsize\textbf{{[}d=6{]}}}
\end{cellvarwidth}} & {\scriptsize\textbf{Best}} & {\scriptsize 0.946351} & {\scriptsize 0.946351} & {\scriptsize 0.946351} & {\scriptsize 0.96} & {\scriptsize 0.946351} & {\scriptsize 0.946351} & {\scriptsize 0.946351} & {\scriptsize 0.946351} & {\scriptsize 0.946351} & {\scriptsize 0.946351}\tabularnewline
\cline{2-12}
 & {\scriptsize\textbf{Mean}} & {\scriptsize 0.946439} & {\scriptsize 0.946361} & {\scriptsize 0.946403} & {\scriptsize 1.01894} & {\scriptsize 0.946432} & {\scriptsize 0.947578} & {\scriptsize 0.946868} & {\scriptsize 0.946351} & {\scriptsize 0.946523} & {\scriptsize 0.946432}\tabularnewline
\cline{2-12}
 & {\scriptsize\textbf{Rate\%}} & {\scriptsize 33\%} & {\scriptsize 97\%} & {\scriptsize 50\%} & {\scriptsize 3\%} & {\scriptsize 73\%} & {\scriptsize 3\%} & {\scriptsize 7\%} & {\scriptsize 97\%} & {\scriptsize 43\%} & {\scriptsize 73\%}\tabularnewline
\cline{2-12}
 & {\scriptsize\textbf{SD}} & {\scriptsize 0.00042572} & {\scriptsize 5.5557e-05} & {\scriptsize 0.000114485} & {\scriptsize 0.0435639} & {\scriptsize 0.000136866} & {\scriptsize 0.00117242} & {\scriptsize 0.00072299} & {\scriptsize 1.0713e-06} & {\scriptsize 0.000153368} & {\scriptsize 0.000136866}\tabularnewline
\cline{2-12}
 & {\scriptsize\textbf{Time(s)}} & {\scriptsize 7.885} & {\scriptsize 4.239} & {\scriptsize 3.139} & {\scriptsize 0.928} & {\scriptsize 3.237} & {\scriptsize 1.543} & {\scriptsize 10.487} & {\scriptsize 2.66} & {\scriptsize 35.607} & {\scriptsize 21.123}\tabularnewline
\hline 
\end{longtable}
\end{landscape}
\par\end{center}

Table \ref{tab:comparison1} reports Best, Mean, Rate\%, SD, and execution
time for all ten methods on each of the forty-five problems. Aggregated
across the full set, SPARQ attains a success rate of 100\% on 15 problems,
or 33.3\% of the total. Its average success rate is 46.6\%, the highest
figure among the ten methods. The next-best average success rates
are 31.4\% for mLSHADE-RL, 30.4\% for EA4Eig, and 27.6\% for UDE3.
ARQ2, LM-CMA-ES, SFCDE, TridentDE, and NL-SHADE-RSP each average below
25\%. The magnitude of this advantage is clearly visible on individual
problems. On alpine1 (d=100), bucherastrigin (d=50), salomon (d=100),
and weierstrass (d=70), SPARQ and LM-CMA-ES jointly reach a mean error
of 0 at 100\% success. The remaining eight methods are confined to
3 to 7\% success. Their mean errors on bucherastrigin, for instance,
range from 18.46 for EA4Eig to 87.45 for TridentDE. On powell (d=100)
and sinusoidal (d=200), SPARQ alone reaches 100\% success, with a
mean error of 0 and -3.5 respectively. Every other method is held
to 3 to 20\% success on these two problems. LM-CMA-ES in particular
sees its sinusoidal mean collapse to -1.02. On the real-world load-dispatch
and engineering problems, namely eld2 through eld5, hydrothermal,
tersoffb, tersoffc, and transmissionpricing, every method in the comparison
records success rates no higher than 3 to 10\%. SPARQ is no exception
here. This pattern indicates that the difficulty is intrinsic to these
landscapes rather than a shortcoming of any single method. In terms
of computational cost, SPARQ's cumulative execution time across all
forty-five problems is 22.04 seconds. This is the second-lowest total
among the ten methods, behind only LM-CMA-ES at 13.47 seconds. ARQ2
requires 60.95 seconds, SFCDE 78.01 seconds, EA4Eig 102.45 seconds,
UDE3 157.60 seconds, and TridentDE 254.43 seconds. TridentDE alone
requires roughly eleven times SPARQ's total runtime for the same benchmark
set.

\begin{table}[H]
\caption{Best-Value Ranking of SPARQ and Nine Competing Metaheuristics Across
the Full Benchmark Suite\label{tab:bestRankin1}}

\begin{centering}
\begin{tabular}{|c|c|c|c|c|c|c|c|c|c|c|}
\hline 
{\tiny\textbf{PROBLEM}} & {\tiny\textbf{arq2}} & {\tiny\textbf{ea4eig}} & {\tiny\textbf{jso}} & {\tiny\textbf{lmcmaes}} & {\tiny\textbf{mlshaderl}} & {\tiny\textbf{nlshadelbc}} & {\tiny\textbf{sfcde}} & {\tiny\textbf{sparq}} & {\tiny\textbf{tridentde}} & {\tiny\textbf{ude3}}\tabularnewline
\hline 
\hline 
{\tiny\textbf{alpine1}} & {\tiny 1} & {\tiny 6} & {\tiny 1} & {\tiny 1} & {\tiny 1} & {\tiny 10} & {\tiny 7} & {\tiny 1} & {\tiny 9} & {\tiny 8}\tabularnewline
\hline 
{\tiny\textbf{antennaarray}} & {\tiny 1} & {\tiny 1} & {\tiny 1} & {\tiny 10} & {\tiny 1} & {\tiny 1} & {\tiny 1} & {\tiny 1} & {\tiny 1} & {\tiny 1}\tabularnewline
\hline 
{\tiny\textbf{antennaula}} & {\tiny 1} & {\tiny 1} & {\tiny 1} & {\tiny 10} & {\tiny 1} & {\tiny 1} & {\tiny 1} & {\tiny 1} & {\tiny 1} & {\tiny 1}\tabularnewline
\hline 
{\tiny\textbf{bucherastrigin}} & {\tiny 3} & {\tiny 6} & {\tiny 5} & {\tiny 1} & {\tiny 8} & {\tiny 7} & {\tiny 4} & {\tiny 1} & {\tiny 10} & {\tiny 9}\tabularnewline
\hline 
{\tiny\textbf{cosinemixture}} & {\tiny 1} & {\tiny 1} & {\tiny 1} & {\tiny 1} & {\tiny 1} & {\tiny 1} & {\tiny 1} & {\tiny 1} & {\tiny 1} & {\tiny 1}\tabularnewline
\hline 
{\tiny\textbf{ded1}} & {\tiny 4} & {\tiny 5} & {\tiny 1} & {\tiny 8} & {\tiny 2} & {\tiny 10} & {\tiny 7} & {\tiny 3} & {\tiny 9} & {\tiny 6}\tabularnewline
\hline 
{\tiny\textbf{ded2}} & {\tiny 4} & {\tiny 3} & {\tiny 1} & {\tiny 9} & {\tiny 6} & {\tiny 10} & {\tiny 8} & {\tiny 5} & {\tiny 7} & {\tiny 2}\tabularnewline
\hline 
{\tiny\textbf{eld1}} & {\tiny 1} & {\tiny 1} & {\tiny 1} & {\tiny 1} & {\tiny 1} & {\tiny 1} & {\tiny 1} & {\tiny 1} & {\tiny 1} & {\tiny 1}\tabularnewline
\hline 
{\tiny\textbf{eld2}} & {\tiny 1} & {\tiny 2} & {\tiny 8} & {\tiny 10} & {\tiny 3} & {\tiny 8} & {\tiny 6} & {\tiny 5} & {\tiny 4} & {\tiny 7}\tabularnewline
\hline 
{\tiny\textbf{eld3}} & {\tiny 1} & {\tiny 5} & {\tiny 8} & {\tiny 10} & {\tiny 1} & {\tiny 9} & {\tiny 7} & {\tiny 5} & {\tiny 1} & {\tiny 1}\tabularnewline
\hline 
{\tiny\textbf{eld4}} & {\tiny 2} & {\tiny 5} & {\tiny 7} & {\tiny 10} & {\tiny 1} & {\tiny 9} & {\tiny 8} & {\tiny 5} & {\tiny 4} & {\tiny 3}\tabularnewline
\hline 
{\tiny\textbf{eld5}} & {\tiny 5} & {\tiny 3} & {\tiny 1} & {\tiny 10} & {\tiny 2} & {\tiny 9} & {\tiny 8} & {\tiny 4} & {\tiny 7} & {\tiny 6}\tabularnewline
\hline 
{\tiny\textbf{fmsynth}} & {\tiny 1} & {\tiny 1} & {\tiny 1} & {\tiny 10} & {\tiny 8} & {\tiny 7} & {\tiny 6} & {\tiny 1} & {\tiny 8} & {\tiny 5}\tabularnewline
\hline 
{\tiny\textbf{gallagher101}} & {\tiny 1} & {\tiny 1} & {\tiny 1} & {\tiny 7} & {\tiny 7} & {\tiny 9} & {\tiny 1} & {\tiny 9} & {\tiny 1} & {\tiny 1}\tabularnewline
\hline 
{\tiny\textbf{gallagher21}} & {\tiny 1} & {\tiny 1} & {\tiny 1} & {\tiny 1} & {\tiny 1} & {\tiny 1} & {\tiny 1} & {\tiny 1} & {\tiny 1} & {\tiny 1}\tabularnewline
\hline 
{\tiny\textbf{griewankrosenbrock}} & {\tiny 3} & {\tiny 7} & {\tiny 4} & {\tiny 5} & {\tiny 9} & {\tiny 1} & {\tiny 6} & {\tiny 2} & {\tiny 8} & {\tiny 10}\tabularnewline
\hline 
{\tiny\textbf{heatexchanger}} & {\tiny 7} & {\tiny 1} & {\tiny 1} & {\tiny 1} & {\tiny 1} & {\tiny 7} & {\tiny 7} & {\tiny 7} & {\tiny 1} & {\tiny 1}\tabularnewline
\hline 
{\tiny\textbf{hydrothermal}} & {\tiny 2} & {\tiny 5} & {\tiny 4} & {\tiny 10} & {\tiny 1} & {\tiny 9} & {\tiny 8} & {\tiny 3} & {\tiny 7} & {\tiny 6}\tabularnewline
\hline 
{\tiny\textbf{katsuura}} & {\tiny 9} & {\tiny 1} & {\tiny 1} & {\tiny 1} & {\tiny 1} & {\tiny 8} & {\tiny 10} & {\tiny 1} & {\tiny 1} & {\tiny 1}\tabularnewline
\hline 
{\tiny\textbf{levy}} & {\tiny 6} & {\tiny 1} & {\tiny 4} & {\tiny 8} & {\tiny 5} & {\tiny 9} & {\tiny 10} & {\tiny 2} & {\tiny 7} & {\tiny 3}\tabularnewline
\hline 
{\tiny\textbf{messenger}} & {\tiny 1} & {\tiny 1} & {\tiny 1} & {\tiny 10} & {\tiny 1} & {\tiny 1} & {\tiny 1} & {\tiny 1} & {\tiny 1} & {\tiny 1}\tabularnewline
\hline 
{\tiny\textbf{michalewicz}} & {\tiny 3} & {\tiny 8} & {\tiny 5} & {\tiny 6} & {\tiny 4} & {\tiny 1} & {\tiny 7} & {\tiny 2} & {\tiny 9} & {\tiny 10}\tabularnewline
\hline 
{\tiny\textbf{polyphase}} & {\tiny 5} & {\tiny 3} & {\tiny 2} & {\tiny 9} & {\tiny 4} & {\tiny 8} & {\tiny 10} & {\tiny 1} & {\tiny 7} & {\tiny 6}\tabularnewline
\hline 
{\tiny\textbf{potential}} & {\tiny 7} & {\tiny 5} & {\tiny 3} & {\tiny 1} & {\tiny 6} & {\tiny 4} & {\tiny 9} & {\tiny 2} & {\tiny 8} & {\tiny 10}\tabularnewline
\hline 
{\tiny\textbf{powell }} & {\tiny 4} & {\tiny 5} & {\tiny 3} & {\tiny 2} & {\tiny 6} & {\tiny 9} & {\tiny 8} & {\tiny 1} & {\tiny 10} & {\tiny 7}\tabularnewline
\hline 
{\tiny\textbf{rosenbrock}} & {\tiny 2} & {\tiny 6} & {\tiny 1} & {\tiny 3} & {\tiny 4} & {\tiny 10} & {\tiny 7} & {\tiny 5} & {\tiny 9} & {\tiny 8}\tabularnewline
\hline 
{\tiny\textbf{rosetta}} & {\tiny 1} & {\tiny 1} & {\tiny 8} & {\tiny 10} & {\tiny 1} & {\tiny 8} & {\tiny 1} & {\tiny 1} & {\tiny 1} & {\tiny 1}\tabularnewline
\hline 
{\tiny\textbf{rotatedrosenbrock}} & {\tiny 7} & {\tiny 10} & {\tiny 2} & {\tiny 1} & {\tiny 3} & {\tiny 8} & {\tiny 4} & {\tiny 5} & {\tiny 6} & {\tiny 9}\tabularnewline
\hline 
{\tiny\textbf{salomon}} & {\tiny 3} & {\tiny 7} & {\tiny 5} & {\tiny 1} & {\tiny 8} & {\tiny 10} & {\tiny 4} & {\tiny 1} & {\tiny 9} & {\tiny 5}\tabularnewline
\hline 
{\tiny\textbf{schwefel}} & {\tiny 1} & {\tiny 1} & {\tiny 1} & {\tiny 10} & {\tiny 1} & {\tiny 1} & {\tiny 1} & {\tiny 1} & {\tiny 1} & {\tiny 9}\tabularnewline
\hline 
{\tiny\textbf{sinusoidal}} & {\tiny 3} & {\tiny 5} & {\tiny 4} & {\tiny 1} & {\tiny 6} & {\tiny 10} & {\tiny 7} & {\tiny 1} & {\tiny 9} & {\tiny 8}\tabularnewline
\hline 
{\tiny\textbf{stirredtankreactor}} & {\tiny 1} & {\tiny 1} & {\tiny 1} & {\tiny 10} & {\tiny 1} & {\tiny 1} & {\tiny 1} & {\tiny 1} & {\tiny 1} & {\tiny 1}\tabularnewline
\hline 
{\tiny\textbf{tandem}} & {\tiny 1} & {\tiny 1} & {\tiny 1} & {\tiny 10} & {\tiny 1} & {\tiny 9} & {\tiny 8} & {\tiny 1} & {\tiny 1} & {\tiny 1}\tabularnewline
\hline 
{\tiny\textbf{tersoffb}} & {\tiny 1} & {\tiny 7} & {\tiny 3} & {\tiny 10} & {\tiny 4} & {\tiny 6} & {\tiny 9} & {\tiny 5} & {\tiny 8} & {\tiny 2}\tabularnewline
\hline 
{\tiny\textbf{tersoffc}} & {\tiny 4} & {\tiny 2} & {\tiny 1} & {\tiny 10} & {\tiny 5} & {\tiny 6} & {\tiny 9} & {\tiny 7} & {\tiny 3} & {\tiny 8}\tabularnewline
\hline 
{\tiny\textbf{test2n}} & {\tiny 4} & {\tiny 3} & {\tiny 7} & {\tiny 10} & {\tiny 5} & {\tiny 2} & {\tiny 8} & {\tiny 1} & {\tiny 9} & {\tiny 6}\tabularnewline
\hline 
{\tiny\textbf{test30n}} & {\tiny 4} & {\tiny 5} & {\tiny 3} & {\tiny 1} & {\tiny 6} & {\tiny 10} & {\tiny 7} & {\tiny 1} & {\tiny 9} & {\tiny 8}\tabularnewline
\hline 
{\tiny\textbf{tnep}} & {\tiny 1} & {\tiny 1} & {\tiny 1} & {\tiny 10} & {\tiny 1} & {\tiny 1} & {\tiny 1} & {\tiny 1} & {\tiny 1} & {\tiny 1}\tabularnewline
\hline 
{\tiny\textbf{transmissionpricing}} & {\tiny 1} & {\tiny 4} & {\tiny 2} & {\tiny 7} & {\tiny 5} & {\tiny 10} & {\tiny 3} & {\tiny 6} & {\tiny 9} & {\tiny 8}\tabularnewline
\hline 
{\tiny\textbf{trid}} & {\tiny 4} & {\tiny 3} & {\tiny 1} & {\tiny 1} & {\tiny 7} & {\tiny 10} & {\tiny 9} & {\tiny 5} & {\tiny 8} & {\tiny 6}\tabularnewline
\hline 
{\tiny\textbf{vibratingplatform}} & {\tiny 7} & {\tiny 1} & {\tiny 1} & {\tiny 10} & {\tiny 1} & {\tiny 1} & {\tiny 9} & {\tiny 7} & {\tiny 1} & {\tiny 1}\tabularnewline
\hline 
{\tiny\textbf{weatherirrigation}} & {\tiny 1} & {\tiny 1} & {\tiny 7} & {\tiny 9} & {\tiny 1} & {\tiny 10} & {\tiny 8} & {\tiny 1} & {\tiny 1} & {\tiny 1}\tabularnewline
\hline 
{\tiny\textbf{weierstrass}} & {\tiny 5} & {\tiny 3} & {\tiny 4} & {\tiny 1} & {\tiny 7} & {\tiny 10} & {\tiny 8} & {\tiny 1} & {\tiny 9} & {\tiny 6}\tabularnewline
\hline 
{\tiny\textbf{whitley}} & {\tiny 4} & {\tiny 6} & {\tiny 5} & {\tiny 10} & {\tiny 7} & {\tiny 9} & {\tiny 2} & {\tiny 1} & {\tiny 3} & {\tiny 8}\tabularnewline
\hline 
{\tiny\textbf{wirelesscoverage}} & {\tiny 1} & {\tiny 1} & {\tiny 1} & {\tiny 10} & {\tiny 1} & {\tiny 1} & {\tiny 1} & {\tiny 1} & {\tiny 1} & {\tiny 1}\tabularnewline
\hline 
\end{tabular}
\par\end{centering}
\end{table}

Table \ref{tab:bestRankin1} reports the ranking of each method according
to its best-value result. It covers the same ten methods and forty-five
problems. SPARQ attains rank 1, whether shared or exclusive, on 25
of the 45 problems. This corresponds to 55.6\%, the highest share
among all ten methods. If only cases where SPARQ is the sole method
at rank 1 are counted, with no ties, this figure falls to 4 of 45,
or 8.9\%. By comparison, jSO achieves 5 of 45 exclusive first places,
and ARQ2 achieves 3 of 45. This indicates that a substantial part
of SPARQ's rank-1 share arises from problems where several methods
jointly reach the global optimum, rather than from unique dominance.
The mean rank across all forty-five problems is 2.64 for SPARQ. This
is the lowest, and therefore best, mean rank of the ten methods. It
sits ahead of jSO at 2.80 and ARQ2 at 2.91. At the other extreme,
LM-CMA-ES and NL-SHADE-RSP average 6.38 and 6.29 respectively. They
register the worst rank, 10th, on 20 problems (44.4\%) and 11 problems
(24.4\%) respectively.

\begin{table}[H]
\caption{Mean-Value Ranking of SPARQ and Nine Competing Metaheuristics Across
the Full Benchmark Suite\label{tab:meanRanking1}}

\centering{}%
\begin{tabular}{|c|c|c|c|c|c|c|c|c|c|c|}
\hline 
{\tiny\textbf{PROBLEM}} & {\tiny\textbf{arq2}} & {\tiny\textbf{ea4eig}} & {\tiny\textbf{jso}} & {\tiny\textbf{lmcmaes}} & {\tiny\textbf{mlshaderl}} & {\tiny\textbf{nlshadelbc}} & {\tiny\textbf{sfcde}} & {\tiny\textbf{sparq}} & {\tiny\textbf{tridentde}} & {\tiny\textbf{ude3}}\tabularnewline
\hline 
\hline 
{\tiny\textbf{alpine1}} & {\tiny 7} & {\tiny 5} & {\tiny 3} & {\tiny 1} & {\tiny 4} & {\tiny 10} & {\tiny 8} & {\tiny 1} & {\tiny 9} & {\tiny 6}\tabularnewline
\hline 
{\tiny\textbf{antennaarray}} & {\tiny 8} & {\tiny 1} & {\tiny 1} & {\tiny 10} & {\tiny 1} & {\tiny 6} & {\tiny 9} & {\tiny 1} & {\tiny 7} & {\tiny 1}\tabularnewline
\hline 
{\tiny\textbf{antennaula}} & {\tiny 1} & {\tiny 1} & {\tiny 1} & {\tiny 10} & {\tiny 1} & {\tiny 1} & {\tiny 1} & {\tiny 1} & {\tiny 9} & {\tiny 1}\tabularnewline
\hline 
{\tiny\textbf{bucherastrigin}} & {\tiny 4} & {\tiny 3} & {\tiny 6} & {\tiny 1} & {\tiny 7} & {\tiny 9} & {\tiny 5} & {\tiny 1} & {\tiny 10} & {\tiny 8}\tabularnewline
\hline 
{\tiny\textbf{cosinemixture}} & {\tiny 1} & {\tiny 1} & {\tiny 1} & {\tiny 1} & {\tiny 7} & {\tiny 8} & {\tiny 10} & {\tiny 1} & {\tiny 9} & {\tiny 1}\tabularnewline
\hline 
{\tiny\textbf{ded1}} & {\tiny 5} & {\tiny 4} & {\tiny 1} & {\tiny 8} & {\tiny 2} & {\tiny 9} & {\tiny 7} & {\tiny 3} & {\tiny 10} & {\tiny 6}\tabularnewline
\hline 
{\tiny\textbf{ded2}} & {\tiny 4} & {\tiny 1} & {\tiny 2} & {\tiny 9} & {\tiny 3} & {\tiny 10} & {\tiny 8} & {\tiny 5} & {\tiny 7} & {\tiny 6}\tabularnewline
\hline 
{\tiny\textbf{eld1}} & {\tiny 1} & {\tiny 2} & {\tiny 3} & {\tiny 10} & {\tiny 8} & {\tiny 4} & {\tiny 5} & {\tiny 6} & {\tiny 9} & {\tiny 7}\tabularnewline
\hline 
{\tiny\textbf{eld2}} & {\tiny 2} & {\tiny 4} & {\tiny 8} & {\tiny 10} & {\tiny 1} & {\tiny 9} & {\tiny 5} & {\tiny 6} & {\tiny 7} & {\tiny 3}\tabularnewline
\hline 
{\tiny\textbf{eld3}} & {\tiny 1} & {\tiny 3} & {\tiny 8} & {\tiny 10} & {\tiny 2} & {\tiny 9} & {\tiny 6} & {\tiny 5} & {\tiny 7} & {\tiny 4}\tabularnewline
\hline 
{\tiny\textbf{eld4}} & {\tiny 2} & {\tiny 4} & {\tiny 5} & {\tiny 10} & {\tiny 1} & {\tiny 9} & {\tiny 8} & {\tiny 6} & {\tiny 7} & {\tiny 3}\tabularnewline
\hline 
{\tiny\textbf{eld5}} & {\tiny 5} & {\tiny 2} & {\tiny 1} & {\tiny 10} & {\tiny 4} & {\tiny 9} & {\tiny 8} & {\tiny 3} & {\tiny 7} & {\tiny 6}\tabularnewline
\hline 
{\tiny\textbf{fmsynth}} & {\tiny 6} & {\tiny 2} & {\tiny 1} & {\tiny 10} & {\tiny 8} & {\tiny 3} & {\tiny 5} & {\tiny 4} & {\tiny 9} & {\tiny 7}\tabularnewline
\hline 
{\tiny\textbf{gallagher101}} & {\tiny 1} & {\tiny 3} & {\tiny 5} & {\tiny 10} & {\tiny 2} & {\tiny 9} & {\tiny 8} & {\tiny 7} & {\tiny 6} & {\tiny 4}\tabularnewline
\hline 
{\tiny\textbf{gallagher21}} & {\tiny 4} & {\tiny 1} & {\tiny 8} & {\tiny 10} & {\tiny 2} & {\tiny 9} & {\tiny 7} & {\tiny 3} & {\tiny 5} & {\tiny 6}\tabularnewline
\hline 
{\tiny\textbf{griewankrosenbrock}} & {\tiny 3} & {\tiny 7} & {\tiny 4} & {\tiny 6} & {\tiny 9} & {\tiny 1} & {\tiny 5} & {\tiny 2} & {\tiny 8} & {\tiny 10}\tabularnewline
\hline 
{\tiny\textbf{heatexchanger}} & {\tiny 6} & {\tiny 4} & {\tiny 5} & {\tiny 10} & {\tiny 1} & {\tiny 6} & {\tiny 6} & {\tiny 6} & {\tiny 1} & {\tiny 1}\tabularnewline
\hline 
{\tiny\textbf{hydrothermal}} & {\tiny 2} & {\tiny 5} & {\tiny 4} & {\tiny 10} & {\tiny 1} & {\tiny 9} & {\tiny 8} & {\tiny 3} & {\tiny 7} & {\tiny 6}\tabularnewline
\hline 
{\tiny\textbf{katsuura}} & {\tiny 8} & {\tiny 1} & {\tiny 1} & {\tiny 1} & {\tiny 1} & {\tiny 9} & {\tiny 10} & {\tiny 1} & {\tiny 1} & {\tiny 1}\tabularnewline
\hline 
{\tiny\textbf{levy}} & {\tiny 6} & {\tiny 1} & {\tiny 3} & {\tiny 10} & {\tiny 5} & {\tiny 8} & {\tiny 9} & {\tiny 4} & {\tiny 7} & {\tiny 2}\tabularnewline
\hline 
{\tiny\textbf{messenger}} & {\tiny 1} & {\tiny 1} & {\tiny 1} & {\tiny 10} & {\tiny 1} & {\tiny 1} & {\tiny 1} & {\tiny 1} & {\tiny 1} & {\tiny 1}\tabularnewline
\hline 
{\tiny\textbf{michalewicz}} & {\tiny 4} & {\tiny 8} & {\tiny 5} & {\tiny 7} & {\tiny 3} & {\tiny 1} & {\tiny 6} & {\tiny 2} & {\tiny 9} & {\tiny 10}\tabularnewline
\hline 
{\tiny\textbf{polyphase}} & {\tiny 6} & {\tiny 2} & {\tiny 4} & {\tiny 10} & {\tiny 3} & {\tiny 8} & {\tiny 9} & {\tiny 1} & {\tiny 7} & {\tiny 5}\tabularnewline
\hline 
{\tiny\textbf{potential}} & {\tiny 5} & {\tiny 4} & {\tiny 3} & {\tiny 10} & {\tiny 6} & {\tiny 2} & {\tiny 7} & {\tiny 1} & {\tiny 8} & {\tiny 9}\tabularnewline
\hline 
{\tiny\textbf{powell}} & {\tiny 6} & {\tiny 4} & {\tiny 3} & {\tiny 2} & {\tiny 5} & {\tiny 10} & {\tiny 8} & {\tiny 1} & {\tiny 9} & {\tiny 7}\tabularnewline
\hline 
{\tiny\textbf{rosenbrock}} & {\tiny 3} & {\tiny 5} & {\tiny 1} & {\tiny 2} & {\tiny 4} & {\tiny 10} & {\tiny 6} & {\tiny 7} & {\tiny 9} & {\tiny 8}\tabularnewline
\hline 
{\tiny\textbf{rosetta}} & {\tiny 6} & {\tiny 7} & {\tiny 8} & {\tiny 10} & {\tiny 3} & {\tiny 9} & {\tiny 1} & {\tiny 1} & {\tiny 5} & {\tiny 4}\tabularnewline
\hline 
{\tiny\textbf{rotatedrosenbrock}} & {\tiny 4} & {\tiny 6} & {\tiny 1} & {\tiny 2} & {\tiny 5} & {\tiny 3} & {\tiny 9} & {\tiny 8} & {\tiny 7} & {\tiny 10}\tabularnewline
\hline 
{\tiny\textbf{salomon}} & {\tiny 7} & {\tiny 6} & {\tiny 5} & {\tiny 1} & {\tiny 8} & {\tiny 10} & {\tiny 4} & {\tiny 1} & {\tiny 9} & {\tiny 3}\tabularnewline
\hline 
{\tiny\textbf{schwefel}} & {\tiny 1} & {\tiny 1} & {\tiny 7} & {\tiny 10} & {\tiny 6} & {\tiny 5} & {\tiny 1} & {\tiny 1} & {\tiny 8} & {\tiny 9}\tabularnewline
\hline 
{\tiny\textbf{sinusoidal}} & {\tiny 5} & {\tiny 3} & {\tiny 2} & {\tiny 10} & {\tiny 4} & {\tiny 9} & {\tiny 7} & {\tiny 1} & {\tiny 8} & {\tiny 6}\tabularnewline
\hline 
{\tiny\textbf{stirredtankreactor}} & {\tiny 1} & {\tiny 1} & {\tiny 1} & {\tiny 10} & {\tiny 1} & {\tiny 1} & {\tiny 1} & {\tiny 1} & {\tiny 1} & {\tiny 1}\tabularnewline
\hline 
{\tiny\textbf{tandem}} & {\tiny 1} & {\tiny 1} & {\tiny 1} & {\tiny 10} & {\tiny 1} & {\tiny 9} & {\tiny 8} & {\tiny 1} & {\tiny 1} & {\tiny 1}\tabularnewline
\hline 
{\tiny\textbf{tersoffb}} & {\tiny 2} & {\tiny 4} & {\tiny 3} & {\tiny 10} & {\tiny 6} & {\tiny 1} & {\tiny 8} & {\tiny 5} & {\tiny 7} & {\tiny 9}\tabularnewline
\hline 
{\tiny\textbf{tersoffc}} & {\tiny 2} & {\tiny 4} & {\tiny 3} & {\tiny 10} & {\tiny 6} & {\tiny 1} & {\tiny 7} & {\tiny 8} & {\tiny 5} & {\tiny 9}\tabularnewline
\hline 
{\tiny\textbf{test2n}} & {\tiny 4} & {\tiny 2} & {\tiny 6} & {\tiny 10} & {\tiny 3} & {\tiny 7} & {\tiny 8} & {\tiny 1} & {\tiny 9} & {\tiny 5}\tabularnewline
\hline 
{\tiny\textbf{test30n}} & {\tiny 5} & {\tiny 2} & {\tiny 3} & {\tiny 1} & {\tiny 6} & {\tiny 10} & {\tiny 8} & {\tiny 4} & {\tiny 9} & {\tiny 7}\tabularnewline
\hline 
{\tiny\textbf{tnep}} & {\tiny 1} & {\tiny 1} & {\tiny 1} & {\tiny 10} & {\tiny 1} & {\tiny 9} & {\tiny 7} & {\tiny 1} & {\tiny 8} & {\tiny 1}\tabularnewline
\hline 
{\tiny\textbf{transmissionpricing}} & {\tiny 4} & {\tiny 3} & {\tiny 1} & {\tiny 9} & {\tiny 2} & {\tiny 10} & {\tiny 7} & {\tiny 5} & {\tiny 8} & {\tiny 6}\tabularnewline
\hline 
{\tiny\textbf{trid}} & {\tiny 5} & {\tiny 3} & {\tiny 2} & {\tiny 1} & {\tiny 6} & {\tiny 10} & {\tiny 8} & {\tiny 4} & {\tiny 9} & {\tiny 7}\tabularnewline
\hline 
{\tiny\textbf{vibratingplatform}} & {\tiny 2} & {\tiny 5} & {\tiny 1} & {\tiny 10} & {\tiny 4} & {\tiny 7} & {\tiny 8} & {\tiny 6} & {\tiny 9} & {\tiny 3}\tabularnewline
\hline 
{\tiny\textbf{weatherirrigation}} & {\tiny 4} & {\tiny 1} & {\tiny 3} & {\tiny 9} & {\tiny 6} & {\tiny 10} & {\tiny 8} & {\tiny 2} & {\tiny 7} & {\tiny 5}\tabularnewline
\hline 
{\tiny\textbf{weierstrass}} & {\tiny 7} & {\tiny 3} & {\tiny 6} & {\tiny 1} & {\tiny 5} & {\tiny 10} & {\tiny 8} & {\tiny 1} & {\tiny 9} & {\tiny 4}\tabularnewline
\hline 
{\tiny\textbf{whitley}} & {\tiny 1} & {\tiny 3} & {\tiny 4} & {\tiny 9} & {\tiny 7} & {\tiny 10} & {\tiny 5} & {\tiny 2} & {\tiny 6} & {\tiny 8}\tabularnewline
\hline 
{\tiny\textbf{wirelesscoverage}} & {\tiny 6} & {\tiny 2} & {\tiny 3} & {\tiny 10} & {\tiny 4} & {\tiny 9} & {\tiny 8} & {\tiny 1} & {\tiny 7} & {\tiny 4}\tabularnewline
\hline 
\end{tabular}
\end{table}

Table \ref{tab:meanRanking1}reports the corresponding ranking based
on mean value. This is a stricter criterion, since it rewards consistency
across repeated runs rather than a single strong result. Under this
criterion, SPARQ ranks first on 20 of 45 problems, or 44.4\%. This
is a lower share than the 55.6\% obtained under the best-value criterion.
Its exclusive, untied first-place count rises slightly to 6 of 45,
compared with 7 of 45 for jSO. Despite this lower first-place frequency,
SPARQ's mean rank across the benchmark set is 3.02. This is marginally
better than EA4Eig at 3.04, and it is the best mean rank among all
ten methods. SPARQ never registers the worst rank, 10th, on any problem.
The divergence between the best-value and mean-value criteria is most
pronounced on rotatedrosenbrock. There, SPARQ's rank worsens from
5th to 8th. On tersoffc it worsens from 7th to 8th. On potential,
by contrast, it improves from 2nd to 1st.

\subsection{Comparative Results on Large-Scale, High-Dimensional Benchmark Problems}

This subsection presents the outcome of the comparative evaluation
on the large-scale portion of the benchmark suite. SPARQ is compared
against five optimizers developed specifically for high-dimensional
search: CCDG2 \citep{CCDG2_1,CCDG2_2}, LM-CMA-ES \citep{LMCMAES},
MMES \citep{MMES}, RMES \citep{RMES}, and VkD-CMA-ES \citep{VkD-CMA-ESc}.
The comparison covers seventeen benchmark problems, each evaluated
at two dimensions, d = 1000 and d = 10,000, for a total of thirty-four
problem-dimension combinations. The problem set spans separable functions
of low intrinsic difficulty alongside strongly non-separable and multimodal
landscapes, allowing both the accuracy and the computational scalability
of each method to be assessed as the number of variables grows by
an order of magnitude. As in the previous subsection, every method
was executed independently over a fixed number of repeated runs under
an identical evaluation budget. Three complementary tables report
the results. Table \ref{tab:comparison2} presents detailed statistics
for each method and problem-dimension combination, namely the best
value found, the mean value, the success rate, the standard deviation,
and the execution time. Table \ref{tab:bestRankin2} reports the ranking
of all six methods according to the best value achieved, and Table
\ref{tab:meanRanking2} reports the corresponding ranking according
to the mean value.
\begin{center}
\begin{landscape}
\begin{longtable}[c]{|c|c|c|c|c|c|c|c|}
\caption{Detailed Performance Comparison of SPARQ Against Large-Scale Optimization
Methods on High-Dimensional Problems (d = 1000, 10,000)\label{tab:comparison2}}
\tabularnewline
\hline 
\textbf{PROBLEM} & \textbf{METRIC} & \textbf{ccdg2} & \textbf{lmcmaes} & \textbf{mmes} & \textbf{rmes} & \textbf{sparq} & \textbf{vkdcmaes}\tabularnewline
\endhead
\hline 
\multirow{5}{*}{\begin{cellvarwidth}[t]
\centering
\textbf{alpine1}

\textbf{{[}d=1000{]}}
\end{cellvarwidth}} & \textbf{Best} & 0 & 0 & 0 & 0 & 0 & 0\tabularnewline
\cline{2-8}
 & \textbf{Mean} & 0 & 0 & 0 & 0 & 0 & 0\tabularnewline
\cline{2-8}
 & \textbf{Rate\%} & 100\% & 100\% & 100\% & 100\% & 100\% & 100\%\tabularnewline
\cline{2-8}
 & \textbf{SD} & 0 & 0 & 0 & 0 & 0 & 0\tabularnewline
\cline{2-8}
 & \textbf{Time(s)} & 2.243 & 1.41 & 0.736 & 0.746 & 0.853 & 5.142\tabularnewline
\hline 
\multirow{5}{*}{\begin{cellvarwidth}[t]
\centering
\textbf{alpine1}

\textbf{{[}d=10,000{]}}
\end{cellvarwidth}} & \textbf{Best} & 0 & 0 & 0 & 0 & 0 & 27033.8\tabularnewline
\cline{2-8}
 & \textbf{Mean} & 0 & 0 & 0 & 0 & 0 & 27246.9\tabularnewline
\cline{2-8}
 & \textbf{Rate\%} & 100\% & 100\% & 100\% & 100\% & 100\% & 3\%\tabularnewline
\cline{2-8}
 & \textbf{SD} & 0 & 0 & 0 & 0 & 0 & 105.421\tabularnewline
\cline{2-8}
 & \textbf{Time(s)} & 1528.92 & 37.223 & 12.18 & 11.334 & 13.855 & 136.395\tabularnewline
\hline 
\multirow{5}{*}{\begin{cellvarwidth}[t]
\centering
\textbf{bucherastrigin}

\textbf{{[}d=1000{]}}
\end{cellvarwidth}} & \textbf{Best} & 0 & 0 & 0 & 0 & 0 & 0\tabularnewline
\cline{2-8}
 & \textbf{Mean} & 0 & 0 & 0 & 0 & 0 & 0\tabularnewline
\cline{2-8}
 & \textbf{Rate\%} & 100\% & 100\% & 100\% & 100\% & 100\% & 100\%\tabularnewline
\cline{2-8}
 & \textbf{SD} & 0 & 0 & 0 & 0 & 0 & 0\tabularnewline
\cline{2-8}
 & \textbf{Time(s)} & 8.306 & 2.834 & 2.195 & 2.221 & 2.27 & 5.264\tabularnewline
\hline 
\multirow{5}{*}{\begin{cellvarwidth}[t]
\centering
\textbf{bucherastrigin}

\textbf{{[}d=10,000{]}}
\end{cellvarwidth}} & \textbf{Best} & 0 & 0 & 0 & 0 & 0 & 2.0878e+06\tabularnewline
\cline{2-8}
 & \textbf{Mean} & 0 & 0 & 0 & 0 & 0 & 2.3368e+06\tabularnewline
\cline{2-8}
 & \textbf{Rate\%} & 100\% & 100\% & 100\% & 100\% & 100\% & 3\%\tabularnewline
\cline{2-8}
 & \textbf{SD} & 0 & 0 & 0 & 0 & 0 & 102010\tabularnewline
\cline{2-8}
 & \textbf{Time(s)} & 1702.33 & 58.263 & 35.485 & 34.966 & 27.656 & 157.158\tabularnewline
\hline 
\multirow{5}{*}{\begin{cellvarwidth}[t]
\centering
\textbf{cigar}

\textbf{{[}d=1000{]}}
\end{cellvarwidth}} & \textbf{Best} & 0 & 0 & 0 & 0 & 0 & 0\tabularnewline
\cline{2-8}
 & \textbf{Mean} & 0 & 0 & 0 & 0 & 0 & 0\tabularnewline
\cline{2-8}
 & \textbf{Rate\%} & 100\% & 100\% & 100\% & 100\% & 100\% & 100\%\tabularnewline
\cline{2-8}
 & \textbf{SD} & 0 & 0 & 0 & 0 & 0 & 0\tabularnewline
\cline{2-8}
 & \textbf{Time(s)} & 1.614 & 1.295 & 0.625 & 0.594 & 0.817 & 4.905\tabularnewline
\hline 
\multirow{5}{*}{\begin{cellvarwidth}[t]
\centering
\textbf{cigar}

\textbf{{[}d=10,000{]}}
\end{cellvarwidth}} & \textbf{Best} & 0 & 0 & 0 & 0 & 0 & 7.2592e+10\tabularnewline
\cline{2-8}
 & \textbf{Mean} & 0 & 0 & 0 & 0 & 0 & 7.3695e+10\tabularnewline
\cline{2-8}
 & \textbf{Rate\%} & 100\% & 100\% & 100\% & 100\% & 100\% & 3\%\tabularnewline
\cline{2-8}
 & \textbf{SD} & 0 & 0 & 0 & 0 & 0 & 4.2344e+08\tabularnewline
\cline{2-8}
 & \textbf{Time(s)} & 1542.51 & 34.293 & 11.212 & 8.851 & 13.589 & 142.792\tabularnewline
\hline 
\multirow{5}{*}{\begin{cellvarwidth}[t]
\centering
\textbf{cosinemixture}

\textbf{{[}d=1000{]}}
\end{cellvarwidth}} & \textbf{Best} & -100 & -100 & -100 & -100 & -100 & -100\tabularnewline
\cline{2-8}
 & \textbf{Mean} & -100 & -100 & -100 & -100 & -100 & -100\tabularnewline
\cline{2-8}
 & \textbf{Rate\%} & 100\% & 100\% & 100\% & 100\% & 100\% & 100\%\tabularnewline
\cline{2-8}
 & \textbf{SD} & 4.3361e-14 & 4.3361e-14 & 4.3361e-14 & 4.3361e-14 & 4.3361e-14 & 4.3361e-14\tabularnewline
\cline{2-8}
 & \textbf{Time(s)} & 2.343 & 1.419 & 0.759 & 0.745 & 0.803 & 5.059\tabularnewline
\hline 
\multirow{5}{*}{\begin{cellvarwidth}[t]
\centering
\textbf{cosinemixture}

\textbf{{[}d=10,000{]}}
\end{cellvarwidth}} & \textbf{Best} & -1000 & -1000 & -1000 & -1000 & -1000 & 2997.85\tabularnewline
\cline{2-8}
 & \textbf{Mean} & -1000 & -1000 & -1000 & -1000 & -1000 & 3036.13\tabularnewline
\cline{2-8}
 & \textbf{Rate\%} & 100\% & 100\% & 100\% & 100\% & 100\% & 3\%\tabularnewline
\cline{2-8}
 & \textbf{SD} & 0 & 0 & 0 & 0 & 0 & 17.5634\tabularnewline
\cline{2-8}
 & \textbf{Time(s)} & 1678.44 & 38.904 & 12.448 & 11.935 & 11.272 & 134.271\tabularnewline
\hline 
\multirow{5}{*}{\begin{cellvarwidth}[t]
\centering
\textbf{differentpowers}

\textbf{{[}d=1000{]}}
\end{cellvarwidth}} & \textbf{Best} & 0 & 0 & 0 & 0 & 0 & 0\tabularnewline
\cline{2-8}
 & \textbf{Mean} & 0 & 0 & 0 & 0 & 0 & 0\tabularnewline
\cline{2-8}
 & \textbf{Rate\%} & 100\% & 100\% & 100\% & 100\% & 100\% & 100\%\tabularnewline
\cline{2-8}
 & \textbf{SD} & 0 & 0 & 0 & 0 & 0 & 0\tabularnewline
\cline{2-8}
 & \textbf{Time(s)} & 3.726 & 1.807 & 1.121 & 1.101 & 1.199 & 5.11\tabularnewline
\hline 
\multirow{5}{*}{\begin{cellvarwidth}[t]
\centering
\textbf{differentpowers}

\textbf{{[}d=10,000{]}}
\end{cellvarwidth}} & \textbf{Best} & 0 & 0 & 0 & 0 & 0 & 3.7419e+06\tabularnewline
\cline{2-8}
 & \textbf{Mean} & 0 & 0 & 0 & 0 & 0 & 3.8879e+06\tabularnewline
\cline{2-8}
 & \textbf{Rate\%} & 100\% & 100\% & 100\% & 100\% & 100\% & 3\%\tabularnewline
\cline{2-8}
 & \textbf{SD} & 0 & 0 & 0 & 0 & 0 & 62082\tabularnewline
\cline{2-8}
 & \textbf{Time(s)} & 1770.91 & 42.456 & 15.265 & 14.604 & 15.986 & 142.936\tabularnewline
\hline 
\multirow{5}{*}{\begin{cellvarwidth}[t]
\centering
\textbf{diracproblem}

\textbf{{[}d=1000{]}}
\end{cellvarwidth}} & \textbf{Best} & 0 & 0 & 0 & 0 & 0 & 0\tabularnewline
\cline{2-8}
 & \textbf{Mean} & 0 & 0 & 0 & 0 & 0 & 0\tabularnewline
\cline{2-8}
 & \textbf{Rate\%} & 100\% & 100\% & 100\% & 100\% & 100\% & 100\%\tabularnewline
\cline{2-8}
 & \textbf{SD} & 0 & 0 & 0 & 0 & 0 & 0\tabularnewline
\cline{2-8}
 & \textbf{Time(s)} & 2.136 & 1.298 & 0.623 & 0.606 & 0.921 & 4.941\tabularnewline
\hline 
\multirow{5}{*}{\begin{cellvarwidth}[t]
\centering
\textbf{diracproblem}

\textbf{{[}d=10,000{]}}
\end{cellvarwidth}} & \textbf{Best} & 0 & 0 & 0 & 0 & 0 & 1\tabularnewline
\cline{2-8}
 & \textbf{Mean} & 0 & 0 & 0 & 0 & 0 & 1\tabularnewline
\cline{2-8}
 & \textbf{Rate\%} & 100\% & 100\% & 100\% & 100\% & 100\% & 100\%\tabularnewline
\cline{2-8}
 & \textbf{SD} & 0 & 0 & 0 & 0 & 0 & 0\tabularnewline
\cline{2-8}
 & \textbf{Time(s)} & 875.296 & 36.69 & 11.133 & 8.779 & 12.794 & 130.536\tabularnewline
\hline 
\multirow{5}{*}{\begin{cellvarwidth}[t]
\centering
\textbf{ellipsoidal}

\textbf{{[}d=1000{]}}
\end{cellvarwidth}} & \textbf{Best} & 0 & 0 & 0 & 0 & 0 & 0\tabularnewline
\cline{2-8}
 & \textbf{Mean} & 0 & 0 & 0 & 0 & 0 & 0\tabularnewline
\cline{2-8}
 & \textbf{Rate\%} & 100\% & 100\% & 100\% & 100\% & 100\% & 100\%\tabularnewline
\cline{2-8}
 & \textbf{SD} & 0 & 0 & 0 & 0 & 0 & 0\tabularnewline
\cline{2-8}
 & \textbf{Time(s)} & 1.672 & 1.312 & 0.611 & 0.596 & 0.798 & 4.976\tabularnewline
\hline 
\multirow{5}{*}{\begin{cellvarwidth}[t]
\centering
\textbf{ellipsoidal}

\textbf{{[}d=10,000{]}}
\end{cellvarwidth}} & \textbf{Best} & 0 & 0 & 0 & 0 & 0 & 4.9311e+09\tabularnewline
\cline{2-8}
 & \textbf{Mean} & 0 & 0 & 0 & 0 & 0 & 5.1532e+09\tabularnewline
\cline{2-8}
 & \textbf{Rate\%} & 100\% & 100\% & 100\% & 100\% & 100\% & 3\%\tabularnewline
\cline{2-8}
 & \textbf{SD} & 0 & 0 & 0 & 0 & 0 & 8.7994e+07\tabularnewline
\cline{2-8}
 & \textbf{Time(s)} & 1577.81 & 38.552 & 11.663 & 9.046 & 14.957 & 140.431\tabularnewline
\hline 
\multirow{5}{*}{\begin{cellvarwidth}[t]
\centering
\textbf{exponential}

\textbf{{[}d=1000{]}}
\end{cellvarwidth}} & \textbf{Best} & 0 & 0 & 0 & 0 & 0 & 0\tabularnewline
\cline{2-8}
 & \textbf{Mean} & 0 & 0 & 0 & 0 & 0 & 0\tabularnewline
\cline{2-8}
 & \textbf{Rate\%} & 100\% & 100\% & 100\% & 100\% & 100\% & 100\%\tabularnewline
\cline{2-8}
 & \textbf{SD} & 0 & 0 & 0 & 0 & 0 & 0\tabularnewline
\cline{2-8}
 & \textbf{Time(s)} & 2.165 & 1.319 & 0.618 & 0.608 & 0.749 & 4.981\tabularnewline
\hline 
\multirow{5}{*}{\textbf{exponential {[}d=10,000{]}}} & \textbf{Best} & 0 & 0 & 0 & 0 & 0 & 1\tabularnewline
\cline{2-8}
 & \textbf{Mean} & 0 & 0 & 0 & 0 & 0 & 1\tabularnewline
\cline{2-8}
 & \textbf{Rate\%} & 100\% & 100\% & 100\% & 100\% & 100\% & 100\%\tabularnewline
\cline{2-8}
 & \textbf{SD} & 0 & 0 & 0 & 0 & 0 & 0\tabularnewline
\cline{2-8}
 & \textbf{Time(s)} & 789.594 & 38.367 & 11.381 & 8.536 & 10.738 & 128.231\tabularnewline
\hline 
\multirow{5}{*}{\begin{cellvarwidth}[t]
\centering
\textbf{lunacekbirastrigin}

\textbf{{[}d=1000{]}}
\end{cellvarwidth}} & \textbf{Best} & 10450.5 & 5566.78 & 6187.62 & 6614.44 & 2299.42 & 19102.3\tabularnewline
\cline{2-8}
 & \textbf{Mean} & 11893.7 & 6069.35 & 7237.64 & 7622.02 & 3098.66 & 20598\tabularnewline
\cline{2-8}
 & \textbf{Rate\%} & 3\% & 3\% & 3\% & 3\% & 3\% & 3\%\tabularnewline
\cline{2-8}
 & \textbf{SD} & 511.293 & 167.9 & 324.119 & 348.468 & 345.159 & 569.204\tabularnewline
\cline{2-8}
 & \textbf{Time(s)} & 8.359 & 1.534 & 0.865 & 0.834 & 1.145 & 5.212\tabularnewline
\hline 
\multirow{5}{*}{\begin{cellvarwidth}[t]
\centering
\textbf{lunacekbirastrigin}

\textbf{{[}d=10,000{]}}
\end{cellvarwidth}} & \textbf{Best} & 134332 & 131825 & 77351.2 & 89482.3 & 89005.2 & 241229\tabularnewline
\cline{2-8}
 & \textbf{Mean} & 136484 & 135368 & 79718.8 & 92535 & 92694.6 & 242922\tabularnewline
\cline{2-8}
 & \textbf{Rate\%} & 3\% & 3\% & 3\% & 3\% & 3\% & 3\%\tabularnewline
\cline{2-8}
 & \textbf{SD} & 970.906 & 1568.31 & 1008.61 & 1348.43 & 1469.83 & 841.349\tabularnewline
\cline{2-8}
 & \textbf{Time(s)} & 1746.11 & 39.801 & 13.175 & 11.84 & 14.625 & 136.589\tabularnewline
\hline 
\multirow{5}{*}{\begin{cellvarwidth}[t]
\centering
\textbf{powell}

\textbf{{[}d=1000{]}}
\end{cellvarwidth}} & \textbf{Best} & 1729.48 & 0.310681 & 0.0503792 & 0.0627767 & 0 & 7578.12\tabularnewline
\cline{2-8}
 & \textbf{Mean} & 2053.37 & 0.491502 & 0.0629321 & 0.0878404 & 0 & 7578.12\tabularnewline
\cline{2-8}
 & \textbf{Rate\%} & 3\% & 3\% & 3\% & 3\% & 100\% & 100\%\tabularnewline
\cline{2-8}
 & \textbf{SD} & 244.11 & 0.111928 & 0.00578802 & 0.0134842 & 0 & 0\tabularnewline
\cline{2-8}
 & \textbf{Time(s)} & 5.871 & 1.284 & 0.618 & 0.607 & 0.815 & 5.006\tabularnewline
\hline 
\multirow{5}{*}{\begin{cellvarwidth}[t]
\centering
\textbf{powell}

\textbf{{[}d=10,000{]}}
\end{cellvarwidth}} & \textbf{Best} & 75750.9 & 75781.2 & 2267.94 & 8860.09 & 0 & 6.9552e+06\tabularnewline
\cline{2-8}
 & \textbf{Mean} & 75750.9 & 75781.2 & 2436.05 & 10050.7 & 0 & 7.7562e+06\tabularnewline
\cline{2-8}
 & \textbf{Rate\%} & 93\% & 100\% & 3\% & 3\% & 100\% & 3\%\tabularnewline
\cline{2-8}
 & \textbf{SD} & 0.000116568 & 0 & 78.3414 & 488.163 & 0 & 350669\tabularnewline
\cline{2-8}
 & \textbf{Time(s)} & 2509.59 & 35.757 & 11.597 & 9.191 & 13.856 & 132.32\tabularnewline
\hline 
\multirow{5}{*}{\begin{cellvarwidth}[t]
\centering
\textbf{rastrigin}

\textbf{{[}d=1000{]}}
\end{cellvarwidth}} & \textbf{Best} & 0 & 0 & 0 & 0 & 0 & 0\tabularnewline
\cline{2-8}
 & \textbf{Mean} & 0 & 0 & 0 & 0 & 0 & 0\tabularnewline
\cline{2-8}
 & \textbf{Rate\%} & 100\% & 100\% & 100\% & 100\% & 100\% & 100\%\tabularnewline
\cline{2-8}
 & \textbf{SD} & 0 & 0 & 0 & 0 & 0 & 0\tabularnewline
\cline{2-8}
 & \textbf{Time(s)} & 2.889 & 1.486 & 0.808 & 0.78 & 0.96 & 5.089\tabularnewline
\hline 
\multirow{5}{*}{\begin{cellvarwidth}[t]
\centering
\textbf{rastrigin}

\textbf{{[}d=10,000{]}}
\end{cellvarwidth}} & \textbf{Best} & 0 & 0 & 0 & 0 & 0 & 168428\tabularnewline
\cline{2-8}
 & \textbf{Mean} & 0 & 0 & 0 & 0 & 0 & 169487\tabularnewline
\cline{2-8}
 & \textbf{Rate\%} & 100\% & 100\% & 100\% & 100\% & 100\% & 3\%\tabularnewline
\cline{2-8}
 & \textbf{SD} & 0 & 0 & 0 & 0 & 0 & 394.995\tabularnewline
\cline{2-8}
 & \textbf{Time(s)} & 1630.46 & 39.001 & 13.052 & 11.606 & 14.462 & 141.418\tabularnewline
\hline 
\multirow{5}{*}{\begin{cellvarwidth}[t]
\centering
\textbf{rotatedrosenbrock}

\textbf{{[}d=1000{]}}
\end{cellvarwidth}} & \textbf{Best} & 419550 & 1859.02 & 1427.36 & 1134.42 & 990.947 & 1.4071e+06\tabularnewline
\cline{2-8}
 & \textbf{Mean} & 499437 & 2285.84 & 1843.56 & 1323.44 & 996.372 & 1.4071e+06\tabularnewline
\cline{2-8}
 & \textbf{Rate\%} & 3\% & 3\% & 3\% & 3\% & 3\% & 100\%\tabularnewline
\cline{2-8}
 & \textbf{SD} & 42352.3 & 207.652 & 209.122 & 94.9028 & 1.86206 & 0\tabularnewline
\cline{2-8}
 & \textbf{Time(s)} & 15.656 & 1.319 & 0.643 & 0.612 & 0.852 & 5.007\tabularnewline
\hline 
\multirow{5}{*}{\begin{cellvarwidth}[t]
\centering
\textbf{rotatedrosenbrock}

\textbf{{[}d=10,000{]}}
\end{cellvarwidth}} & \textbf{Best} & 1.4079e+07 & 1.4084e+07 & 103692 & 107959 & 9961.05 & 1.1383e+09\tabularnewline
\cline{2-8}
 & \textbf{Mean} & 1.4079e+07 & 1.4084e+07 & 109670 & 112869 & 9983.2 & 1.1627e+09\tabularnewline
\cline{2-8}
 & \textbf{Rate\%} & 63\% & 100\% & 3\% & 3\% & 3\% & 3\%\tabularnewline
\cline{2-8}
 & \textbf{SD} & 2.13849 & 0 & 2834.64 & 3074.61 & 14.6511 & 1.0213e+07\tabularnewline
\cline{2-8}
 & \textbf{Time(s)} & 2699.14 & 36.57 & 13.065 & 11.634 & 18.34 & 140.815\tabularnewline
\hline 
\multirow{5}{*}{\begin{cellvarwidth}[t]
\centering
\textbf{salomon}

\textbf{{[}d=1000{]}}
\end{cellvarwidth}} & \textbf{Best} & 0 & 0 & 0 & 0 & 0 & 0\tabularnewline
\cline{2-8}
 & \textbf{Mean} & 0 & 0 & 0 & 0 & 0 & 0\tabularnewline
\cline{2-8}
 & \textbf{Rate\%} & 100\% & 100\% & 100\% & 100\% & 100\% & 100\%\tabularnewline
\cline{2-8}
 & \textbf{SD} & 0 & 0 & 0 & 0 & 0 & 0\tabularnewline
\cline{2-8}
 & \textbf{Time(s)} & 26.613 & 1.278 & 0.627 & 0.612 & 0.657 & 4.986\tabularnewline
\hline 
\multirow{5}{*}{\begin{cellvarwidth}[t]
\centering
\textbf{salomon}

\textbf{{[}d=10,000{]}}
\end{cellvarwidth}} & \textbf{Best} & 0 & 0 & 0 & 0 & 0 & 540.428\tabularnewline
\cline{2-8}
 & \textbf{Mean} & 0 & 0 & 0 & 0 & 0 & 543.982\tabularnewline
\cline{2-8}
 & \textbf{Rate\%} & 100\% & 100\% & 100\% & 100\% & 100\% & 3\%\tabularnewline
\cline{2-8}
 & \textbf{SD} & 0 & 0 & 0 & 0 & 0 & 1.44276\tabularnewline
\cline{2-8}
 & \textbf{Time(s)} & 1607.43 & 33.701 & 11.967 & 8.413 & 9.817 & 131.777\tabularnewline
\hline 
\multirow{5}{*}{\begin{cellvarwidth}[t]
\centering
\textbf{sphere}

\textbf{{[}d=1000{]}}
\end{cellvarwidth}} & \textbf{Best} & 0 & 0 & 0 & 0 & 0 & 0\tabularnewline
\cline{2-8}
 & \textbf{Mean} & 0 & 0 & 0 & 0 & 0 & 0\tabularnewline
\cline{2-8}
 & \textbf{Rate\%} & 100\% & 100\% & 100\% & 100\% & 100\% & 100\%\tabularnewline
\cline{2-8}
 & \textbf{SD} & 0 & 0 & 0 & 0 & 0 & 0\tabularnewline
\cline{2-8}
 & \textbf{Time(s)} & 1.621 & 1.278 & 0.621 & 0.593 & 0.796 & 4.95\tabularnewline
\hline 
\multirow{5}{*}{\begin{cellvarwidth}[t]
\centering
\textbf{sphere}

\textbf{{[}d=10,000{]}}
\end{cellvarwidth}} & \textbf{Best} & 0 & 0 & 0 & 0 & 0 & 2.9038e+07\tabularnewline
\cline{2-8}
 & \textbf{Mean} & 0 & 0 & 0 & 0 & 0 & 2.9480e+07\tabularnewline
\cline{2-8}
 & \textbf{Rate\%} & 100\% & 100\% & 100\% & 100\% & 100\% & 3\%\tabularnewline
\cline{2-8}
 & \textbf{SD} & 0 & 0 & 0 & 0 & 0 & 169617\tabularnewline
\cline{2-8}
 & \textbf{Time(s)} & 1581.58 & 35.094 & 12.069 & 9.123 & 13.391 & 133.305\tabularnewline
\hline 
\multirow{5}{*}{\begin{cellvarwidth}[t]
\centering
\textbf{stepellipsoidal}

\textbf{{[}d=1000{]}}
\end{cellvarwidth}} & \textbf{Best} & 0 & 0 & 0 & 0 & 0 & 0\tabularnewline
\cline{2-8}
 & \textbf{Mean} & 0 & 0 & 0 & 0 & 0 & 0\tabularnewline
\cline{2-8}
 & \textbf{Rate\%} & 100\% & 100\% & 100\% & 100\% & 100\% & 100\%\tabularnewline
\cline{2-8}
 & \textbf{SD} & 0 & 0 & 0 & 0 & 0 & 0\tabularnewline
\cline{2-8}
 & \textbf{Time(s)} & 1.906 & 1.391 & 0.682 & 0.656 & 0.801 & 5.043\tabularnewline
\hline 
\multirow{5}{*}{\begin{cellvarwidth}[t]
\centering
\textbf{stepellipsoidal}

\textbf{{[}d=10,000{]}}
\end{cellvarwidth}} & \textbf{Best} & 0 & 0 & 0 & 0 & 0 & 1.5546e+06\tabularnewline
\cline{2-8}
 & \textbf{Mean} & 0 & 0 & 0 & 0 & 0 & 1.5902e+06\tabularnewline
\cline{2-8}
 & \textbf{Rate\%} & 100\% & 100\% & 100\% & 100\% & 100\% & 3\%\tabularnewline
\cline{2-8}
 & \textbf{SD} & 0 & 0 & 0 & 0 & 0 & 12750.3\tabularnewline
\cline{2-8}
 & \textbf{Time(s)} & 1365 & 38.333 & 13.859 & 10.538 & 10.578 & 137.022\tabularnewline
\hline 
\multirow{5}{*}{\begin{cellvarwidth}[t]
\centering
\textbf{vincent}

\textbf{{[}d=1000{]}}
\end{cellvarwidth}} & \textbf{Best} & -0.40637 & -0.899759 & -1 & -1 & -0.999548 & -0.0202873\tabularnewline
\cline{2-8}
 & \textbf{Mean} & -0.396153 & -0.700881 & -1 & -1 & -0.976253 & 0.00595358\tabularnewline
\cline{2-8}
 & \textbf{Rate\%} & 3\% & 3\% & 40\% & 13\% & 3\% & 3\%\tabularnewline
\cline{2-8}
 & \textbf{SD} & 0.005192 & 0.0936933 & 1.6216e-08 & 2.3654e-07 & 0.0153607 & 0.0109219\tabularnewline
\cline{2-8}
 & \textbf{Time(s)} & 5.736 & 1.849 & 1.09 & 1.082 & 1.306 & 5.453\tabularnewline
\hline 
\multirow{5}{*}{\begin{cellvarwidth}[t]
\centering
\textbf{vincent}

\textbf{{[}d=10,000{]}}
\end{cellvarwidth}} & \textbf{Best} & -0.163397 & -0.191599 & -0.909896 & -0.871154 & -0.714049 & 0.0294198\tabularnewline
\cline{2-8}
 & \textbf{Mean} & -0.154748 & -0.153332 & -0.891337 & -0.855477 & -0.614299 & 0.0368099\tabularnewline
\cline{2-8}
 & \textbf{Rate\%} & 3\% & 3\% & 3\% & 3\% & 3\% & 3\%\tabularnewline
\cline{2-8}
 & \textbf{SD} & 0.00364364 & 0.0124093 & 0.00912978 & 0.00786513 & 0.0368712 & 0.00350821\tabularnewline
\cline{2-8}
 & \textbf{Time(s)} & 1158.81 & 43.701 & 16.806 & 15.609 & 16.423 & 142.12\tabularnewline
\hline 
\multirow{5}{*}{\begin{cellvarwidth}[t]
\centering
\textbf{zakharov}

\textbf{{[}d=1000{]}}
\end{cellvarwidth}} & \textbf{Best} & 4891.83 & 38331.5 & 14956.8 & 9384.52 & 0 & 31780.9\tabularnewline
\cline{2-8}
 & \textbf{Mean} & 7411.23 & 4.8237e+18 & 18096.4 & 10612.6 & 0 & 40341.7\tabularnewline
\cline{2-8}
 & \textbf{Rate\%} & 3\% & 3\% & 3\% & 3\% & 100\% & 3\%\tabularnewline
\cline{2-8}
 & \textbf{SD} & 2005.16 & 1.7063e+19 & 2229.57 & 655.413 & 0 & 3646.85\tabularnewline
\cline{2-8}
 & \textbf{Time(s)} & 27.395 & 1.31 & 0.647 & 0.63 & 0.699 & 4.972\tabularnewline
\hline 
\multirow{5}{*}{\begin{cellvarwidth}[t]
\centering
\textbf{zakharov}

\textbf{{[}d=10,000{]}}
\end{cellvarwidth}} & \textbf{Best} & 181819 & 5.3179e+29 & 457099 & 142882 & 0 & 416320\tabularnewline
\cline{2-8}
 & \textbf{Mean} & 212905 & 1.8168e+30 & 2.6472e+28 & 147899 & 0 & 467117\tabularnewline
\cline{2-8}
 & \textbf{Rate\%} & 3\% & 3\% & 3\% & 3\% & 100\% & 3\%\tabularnewline
\cline{2-8}
 & \textbf{SD} & 14567.4 & 3.7902e+29 & 1.7509e+28 & 2590.91 & 0 & 25296.3\tabularnewline
\cline{2-8}
 & \textbf{Time(s)} & 3400.22 & 37.846 & 11.877 & 9.248 & 9.805 & 136.596\tabularnewline
\hline 
\end{longtable}
\end{landscape}
\par\end{center}

Table \ref{tab:comparison2} shows that, aggregated across all thirty-four
cases, SPARQ achieves a success rate of 100\% on 28 combinations,
or 82.4\%. Its average success rate is 82.9\%, the highest among the
six methods. The closest competitor by average success rate is LM-CMA-ES
at 77.2\%, followed by CCDG2 at 75.9\%, MMES at 72.6\%, and RMES at
71.8\%. VkD-CMA-ES trails markedly at 48.6\%. On twelve separable,
low-difficulty problems, all six methods reach an exact mean error
of 0 at 100\% success. Here the differentiation rests entirely on
runtime. On bucherastrigin (d=10,000), CCDG2 requires 1702.33 seconds,
against 58.26 for LM-CMA-ES, 35.49 for MMES, 34.97 for RMES, and 27.65
for SPARQ. On alpine1 (d=10,000), CCDG2 requires 1528.92 seconds against
13.85 for SPARQ, a difference exceeding two orders of magnitude. Summed
across all thirty-four combinations, CCDG2's cumulative runtime reaches
29,284.4 seconds, against 258.6 for SPARQ, making CCDG2 roughly 113
times slower overall. On the harder, non-separable problems the accuracy
gap becomes substantial. On zakharov (d=10,000), SPARQ attains a mean
error of 0 at 100\% success, while LM-CMA-ES registers $1.82\times10^{30}$,
MMES $2.65\times10^{28}$, VkD-CMA-ES $4.67\times10^{5}$, CCDG2 $2.13\times10^{5}$,
and RMES $1.48\times10^{5}$, all at 3\% success. On powell (d=10,000),
SPARQ again reaches 0 at 100\% success, against $7.56\times10^{6}$
for VkD-CMA-ES, 75,781 for LM-CMA-ES (itself at 100\% success but
converging to a far inferior value), 75,751 for CCDG2 at 93\% success,
10,051 for RMES, and 2436 for MMES. A notable nuance emerges on powell
(d=1000) and rotatedrosenbrock at both dimensions, where VkD-CMA-ES
also reports a 100\% success rate with a standard deviation of exactly
0, yet with mean values many orders of magnitude worse than SPARQ's:
7578.1 on powell (d=1000), and $1.407\times10^{6}$ and $1.163\times10^{9}$
on rotatedrosenbrock. This indicates that VkD-CMA-ES converges highly
consistently to a markedly suboptimal point, rather than failing to
converge at all. The one problem on which SPARQ is clearly outperformed
is vincent. At d=1000, MMES reaches the exact optimum, with a mean
of -1 at 40\% success, and RMES likewise reaches -1 at 13\% success,
against SPARQ's mean of -0.976 at 3\% success. This gap narrows in
relative terms at d=10,000 but persists, with MMES reaching -0.891,
RMES -0.855, and SPARQ -0.614.

\begin{table}[H]
\caption{Best-Value Ranking of SPARQ and Five Large-Scale Optimizers at High
Dimensionality\label{tab:bestRankin2}}

\begin{centering}
\begin{tabular}{|c|c|c|c|c|c|c|c|}
\hline 
PROBLEM & DIM & ccdg2 & lmcmaes & mmes & rmes & sparq & vkdcmaes\tabularnewline
\hline 
\hline 
alpine1 & 1000 & 1 & 1 & 1 & 1 & 1 & 1\tabularnewline
\hline 
alpine1 & 10000 & 1 & 1 & 1 & 1 & 1 & 6\tabularnewline
\hline 
bucherastrigin & 1000 & 1 & 1 & 1 & 1 & 1 & 1\tabularnewline
\hline 
bucherastrigin  & 10000 & 1 & 1 & 1 & 1 & 1 & 6\tabularnewline
\hline 
cigar & 1000 & 1 & 1 & 1 & 1 & 1 & 1\tabularnewline
\hline 
cigar & 10000 & 1 & 1 & 1 & 1 & 1 & 6\tabularnewline
\hline 
cosinemixture & 1000 & 1 & 1 & 1 & 1 & 1 & 1\tabularnewline
\hline 
cosinemixture & 10000 & 1 & 1 & 1 & 1 & 1 & 6\tabularnewline
\hline 
differentpowers & 1000 & 1 & 1 & 1 & 1 & 1 & 1\tabularnewline
\hline 
differentpowers & 10000 & 1 & 1 & 1 & 1 & 1 & 6\tabularnewline
\hline 
diracproblem & 1000 & 1 & 1 & 1 & 1 & 1 & 1\tabularnewline
\hline 
diracproblem & 10000 & 1 & 1 & 1 & 1 & 1 & 6\tabularnewline
\hline 
ellipsoidal & 1000 & 1 & 1 & 1 & 1 & 1 & 1\tabularnewline
\hline 
ellipsoidal & 10000 & 1 & 1 & 1 & 1 & 1 & 6\tabularnewline
\hline 
exponential & 1000 & 1 & 1 & 1 & 1 & 1 & 1\tabularnewline
\hline 
exponential & 10000 & 1 & 1 & 1 & 1 & 1 & 6\tabularnewline
\hline 
lunacekbirastrigin & 1000 & 5 & 2 & 3 & 4 & 1 & 6\tabularnewline
\hline 
lunacekbirastrigin & 10000 & 5 & 4 & 1 & 3 & 2 & 6\tabularnewline
\hline 
powell & 1000 & 5 & 4 & 2 & 3 & 1 & 6\tabularnewline
\hline 
powell & 10000 & 4 & 5 & 2 & 3 & 1 & 6\tabularnewline
\hline 
rastrigin & 1000 & 1 & 1 & 1 & 1 & 1 & 1\tabularnewline
\hline 
rastrigin & 10000 & 1 & 1 & 1 & 1 & 1 & 6\tabularnewline
\hline 
rotatedrosenbrock & 1000 & 5 & 4 & 3 & 2 & 1 & 6\tabularnewline
\hline 
rotatedrosenbrock & 10000 & 4 & 5 & 2 & 3 & 1 & 6\tabularnewline
\hline 
salomon & 1000 & 1 & 1 & 1 & 1 & 1 & 1\tabularnewline
\hline 
salomon & 10000 & 1 & 1 & 1 & 1 & 1 & 6\tabularnewline
\hline 
sphere & 1000 & 1 & 1 & 1 & 1 & 1 & 1\tabularnewline
\hline 
sphere & 10000 & 1 & 1 & 1 & 1 & 1 & 6\tabularnewline
\hline 
stepellipsoidal & 1000 & 1 & 1 & 1 & 1 & 1 & 1\tabularnewline
\hline 
stepellipsoidal & 10000 & 1 & 1 & 1 & 1 & 1 & 6\tabularnewline
\hline 
vincent & 1000 & 5 & 4 & 1 & 1 & 3 & 6\tabularnewline
\hline 
vincent & 10000 & 5 & 4 & 1 & 2 & 3 & 6\tabularnewline
\hline 
zakharov & 1000 & 2 & 6 & 4 & 3 & 1 & 5\tabularnewline
\hline 
zakharov & 10000 & 3 & 6 & 5 & 2 & 1 & 4\tabularnewline
\hline 
\end{tabular}
\par\end{centering}
\end{table}

Table \ref{tab:bestRankin2} reports the best-value ranking of the
same six large-scale methods across the thirty-four problem-dimension
combinations. SPARQ achieves rank 1, shared or exclusive, in 31 of
34 cases, or 91.2\%. This is the highest share among all six methods.
Its exclusive, untied first-place count is 7 of 34, or 20.6\%. This
is higher than any competitor. MMES attains 2 of 34 exclusively, and
all others attain 0. SPARQ's mean rank across the thirty-four cases
is 1.15, the best of the six methods. MMES follows at 1.41 and RMES
at 1.47. There are two exceptions to SPARQ's first-place standing.
On lunacekbirastrigin (d=10,000) it ranks 2nd. On vincent, at both
dimensions, it ranks 3rd. VkD-CMA-ES, by contrast, shares rank 1 with
the other five methods on every problem at d=1000. It drops to rank
6, the worst position, on every single problem at d=10,000. Overall
it registers the worst rank in 20 of 34 cases, or 58.8\%. This pattern
indicates a pronounced scalability limitation at the highest dimensionality
tested.

\begin{table}[H]
\caption{Mean-Value Ranking of SPARQ and Five Large-Scale Optimizers at High
Dimensionality\label{tab:meanRanking2}}

\centering{}%
\begin{tabular}{|c|c|c|c|c|c|c|c|}
\hline 
PROBLEM & DIM & ccdg2 & lmcmaes & mmes & rmes & sparq & vkdcmaes\tabularnewline
\hline 
\hline 
alpine1 & 1000 & 1 & 1 & 1 & 1 & 1 & 1\tabularnewline
\hline 
alpine1 & 10000 & 1 & 1 & 1 & 1 & 1 & 6\tabularnewline
\hline 
bucherastrigin & 1000 & 1 & 1 & 1 & 1 & 1 & 1\tabularnewline
\hline 
bucherastrigin & 10000 & 1 & 1 & 1 & 1 & 1 & 6\tabularnewline
\hline 
cigar & 1000 & 1 & 1 & 1 & 1 & 1 & 1\tabularnewline
\hline 
cigar & 10000 & 1 & 1 & 1 & 1 & 1 & 6\tabularnewline
\hline 
cosinemixture & 1000 & 1 & 1 & 1 & 1 & 1 & 1\tabularnewline
\hline 
cosinemixture & 10000 & 1 & 1 & 1 & 1 & 1 & 6\tabularnewline
\hline 
differentpowers & 1000 & 1 & 1 & 1 & 1 & 1 & 1\tabularnewline
\hline 
differentpowers & 10000 & 1 & 1 & 1 & 1 & 1 & 6\tabularnewline
\hline 
diracproblem & 1000 & 1 & 1 & 1 & 1 & 1 & 1\tabularnewline
\hline 
diracproblem & 10000 & 1 & 1 & 1 & 1 & 1 & 6\tabularnewline
\hline 
ellipsoidal & 1000 & 1 & 1 & 1 & 1 & 1 & 1\tabularnewline
\hline 
ellipsoidal & 10000 & 1 & 1 & 1 & 1 & 1 & 6\tabularnewline
\hline 
exponential & 1000 & 1 & 1 & 1 & 1 & 1 & 1\tabularnewline
\hline 
exponential & 10000 & 1 & 1 & 1 & 1 & 1 & 6\tabularnewline
\hline 
lunacekbirastrigin & 1000 & 5 & 2 & 3 & 4 & 1 & 6\tabularnewline
\hline 
lunacekbirastrigin & 10000 & 5 & 4 & 1 & 2 & 3 & 6\tabularnewline
\hline 
powell & 1000 & 5 & 4 & 2 & 3 & 1 & 6\tabularnewline
\hline 
powell & 10000 & 4 & 5 & 2 & 3 & 1 & 6\tabularnewline
\hline 
rastrigin & 1000 & 1 & 1 & 1 & 1 & 1 & 1\tabularnewline
\hline 
rastrigin & 10000 & 1 & 1 & 1 & 1 & 1 & 6\tabularnewline
\hline 
rotatedrosenbrock & 1000 & 5 & 4 & 3 & 2 & 1 & 6\tabularnewline
\hline 
rotatedrosenbrock & 10000 & 4 & 5 & 2 & 3 & 1 & 6\tabularnewline
\hline 
salomon & 1000 & 1 & 1 & 1 & 1 & 1 & 1\tabularnewline
\hline 
salomon & 10000 & 1 & 1 & 1 & 1 & 1 & 6\tabularnewline
\hline 
sphere & 1000 & 1 & 1 & 1 & 1 & 1 & 1\tabularnewline
\hline 
sphere & 10000 & 1 & 1 & 1 & 1 & 1 & 6\tabularnewline
\hline 
stepellipsoidal & 1000 & 1 & 1 & 1 & 1 & 1 & 1\tabularnewline
\hline 
stepellipsoidal & 10000 & 1 & 1 & 1 & 1 & 1 & 6\tabularnewline
\hline 
vincent & 1000 & 5 & 4 & 1 & 2 & 3 & 6\tabularnewline
\hline 
vincent & 10000 & 4 & 5 & 1 & 2 & 3 & 6\tabularnewline
\hline 
zakharov & 1000 & 2 & 6 & 4 & 3 & 1 & 5\tabularnewline
\hline 
zakharov & 10000 & 3 & 6 & 5 & 2 & 1 & 4\tabularnewline
\hline 
\end{tabular}
\end{table}

Table \ref{tab:meanRanking2} reports the corresponding mean-value
ranking for the same six methods. The resulting figures are nearly
identical to those of Table 5. SPARQ again ranks first in 31 of 34
cases, or 91.2\%. Its exclusive first-place count is again 7 of 34.
It holds the lowest mean rank of the six methods, at 1.18, against
1.41 for MMES. The same two exceptions persist. On lunacekbirastrigin
(d=10,000), SPARQ's rank is 3rd rather than the 2nd observed under
the best-value criterion. On vincent, at both dimensions, it again
ranks 3rd. The best-value and mean-value rankings nearly coincide
at this scale. This contrasts with the more pronounced divergence
observed between Tables \ref{tab:bestRankin1} and \ref{tab:meanRanking1}
at lower dimensionality. It indicates that SPARQ's behavior becomes
markedly more consistent as dimensionality, and the corresponding
evaluation budget, increases.

\subsection{Computational Complexity and Scalability with Dimension\label{subsec:Complexity}}

To assess how the computational cost of SPARQ scales with problem
dimensionality, independent of the difficulty of locating the global
optimum, five separable, unimodal or mildly multimodal benchmark functions
were selected: Ellipsoidal, Exponential, Salomon, Sphere, and Zakharov.
These functions were chosen specifically because they all share a
known global minimum of exactly zero, which allows the reported best,
mean, and standard deviation values to be verified as identically
zero across all tested dimensions, with a 100\% success rate in every
case. This removes search difficulty as a confounding factor and isolates
raw per-evaluation computational overhead as the sole quantity of
interest. Each function was evaluated at six problem dimensions, D
= 100, 2000, 4000, 6000, 8000, and 10,000, with the reported execution
time corresponding to the total wall-clock time required to complete
30 independent runs per dimension, executed in parallel mode. Figure
\ref{fig:complexity} reports the resulting execution time, in seconds,
as a function of D for each of the five functions.

\begin{figure}[H]
\begin{centering}
\includegraphics[scale=0.8]{complexity}
\par\end{centering}
\caption{Execution time of SPARQ as a function of problem dimension D (100--10,000)
for five separable benchmark functions (Ellipsoidal, Exponential,
Salomon, Sphere, Zakharov), each solved to the exact global optimum
with 100\% success rate.\label{fig:complexity}}

\end{figure}

The results show that execution time grows monotonically with dimension
for all five functions, but the growth is distinctly non-linear: the
increase from D = 100 to D = 2000 is comparatively modest, whereas
each subsequent step of 2000 dimensions contributes a progressively
larger absolute increment, with the steepest rise occurring between
D = 8000 and D = 10,000. This pattern is consistent with a per-iteration
cost that is dominated by operations scaling worse than linearly in
D, such as the periodic eigen-basis refresh (Section \ref{subsec:EigenCrossover}),
whose cost grows with the square or cube of the dimension whenever
it is triggered, rather than by the $O(D)$ cost of the core DE mutation
and crossover operators alone. Despite this super-linear trend, the
absolute spread between the fastest and slowest functions remains
narrow throughout the tested range, from under 0.2 s at D = 100 to
under 7 s at D = 10,000, indicating that SPARQ's runtime is governed
primarily by dimension rather than by the specific separable landscape
being optimized. Among the five functions, Salomon and Zakharov consistently
exhibit the lowest execution times at high dimension, while Exponential
and Ellipsoidal are consistently the most expensive, a difference
attributable to their internal evaluation cost per call rather than
to any additional search effort, since all five problems are solved
to the exact global optimum with a 100\% success rate throughout.
These findings indicate that SPARQ's computational overhead scales
predictably and remains practically manageable even at D = 10,000,
supporting its applicability to large-scale optimization tasks without
a disproportionate increase in wall-clock cost.

\section{Conclusions\label{sec:conclusions}}

This work introduced SPARQ, a differential-evolution-based optimizer
that extends the ARQ2 architecture by replacing nearly every one of
its fixed or singly-adapted mechanisms with an adaptive counterpart,
and by adding several mechanisms with no counterpart in its predecessor,
including a stagnation-gated elite polish phase, a trajectory-echo
memory, and a Solis-Wets local search probe. The experimental evaluation
covered both standard-dimensionality and large-scale problem sets,
comprising a combined total of seventy-nine problem instances and
spanning synthetic test functions alongside real-world engineering
formulations.

Across the standard-dimensionality benchmark, SPARQ achieved the highest
average success rate and the lowest mean rank among ten competing
methods, both under the best-value and the mean-value criterion. Its
advantage was most pronounced on problems where competing methods
were confined to low single-digit success rates, while on real-world
problems with intrinsically difficult landscapes, such as the economic
load dispatch family and several chemical-engineering formulations,
no method in the comparison, SPARQ included, achieved consistently
high success rates, indicating that the difficulty in these cases
is a property of the problem rather than of any particular optimizer.
A similar pattern held on the large-scale benchmark, where SPARQ was
compared against five optimizers developed specifically for high-dimensional
search: it attained the highest success rate and the lowest mean rank
of the six methods, with the one clear exception being the vincent
function, on which two competing methods reached the exact optimum
more reliably.

The results support the central premise of this work, namely that
extending nearly every component of ARQ2's architecture with an adaptive
counterpart yields a method that performs competitively on both real-world
and synthetic problems, and at both standard and large-scale dimensionality,
without being restricted to a narrow class of problem structures.
A further observation is that this breadth of adaptivity does not
come at a disproportionate computational cost: despite the substantially
larger number of mechanisms it coordinates relative to its predecessor
and to several competing methods, SPARQ's execution times remained
among the lowest recorded in both parts of the study, rather than
scaling upward with the added algorithmic complexity.

At the same time, the results also delineate the boundaries of this
approach. SPARQ's reliability advantage narrows or disappears on a
specific subset of problems, most notably vincent among the large-scale
functions and several real-world load-dispatch and material-potential
problems among the standard-dimensionality set, suggesting that the
current balance of its exploration and exploitation mechanisms is
better suited to some landscape types than others.

Several directions follow naturally from the present findings. First,
the subset of problems on which SPARQ's advantage narrows, particularly
vincent and the economic load dispatch family, warrants closer examination
to determine whether their difficulty stems from deceptive multimodality,
ill-conditioning, or a mismatch with the current elite-polish and
quarantine mechanisms, and whether a targeted adjustment could close
this gap without weakening performance elsewhere. Second, the individual
contribution of each of SPARQ's adaptive components was not isolated
in this study; a systematic ablation analysis would help identify
which mechanisms drive the bulk of the observed improvement over ARQ2
and which contribute only marginally, informing a leaner variant of
the algorithm. Third, the present evaluation was limited to box-constrained,
single-objective continuous problems; extending SPARQ to handle general
nonlinear constraints, multiple objectives, or mixed-integer variables
would broaden its applicability to a wider range of engineering design
problems. Finally, given that several of SPARQ's mechanisms, such
as the SHADE-style memory and the Thompson-sampling bandit, involve
internally tunable hyperparameters, an empirical sensitivity study
of these settings across problem classes would clarify the robustness
of the default configuration used in this work.

\appendixtitles{no}

\begin{adjustwidth}{-\extralength}{0cm}{}

\begin{thebibliography}{999}
\bibitem{problem}Wolpert, D. H., \& Macready, W. G. (1997). No free
lunch theorems for optimization. IEEE Transactions on Evolutionary
Computation, 1(1), 67--82. DOI: 10.1109/4235.585893

\bibitem{SolisAndWets}Solis, F. J., \& Wets, R. J.-B. (1981). Minimization
by random search techniques. Mathematics of Operations Research, 6(1),
19--30. DOI: 10.1287/moor.6.1.19

\bibitem{SA}Kirkpatrick, S., Gelatt, C. D., \& Vecchi, M. P. (1983).
Optimization by simulated annealing. Science, 220(4598), 671--680.
DOI: 10.1126/science.220.4598.671

\bibitem{GA}Holland, J. H. (1975). Adaptation in Natural and Artificial
Systems. University of Michigan Press.

\bibitem{GiantArmadillo}Kyrou, G., Charilogis, V., \& Tsoulos, I.
G. (2024). Improving the Giant-Armadillo Optimization Method. Analytics,
3(2), 225--240. DOI: 10.3390/analytics3020013

\bibitem{PSO}Kennedy, J., \& Eberhart, R. (1995). Particle swarm
optimization. Proceedings of the IEEE International Conference on
Neural Networks, 4, 1942--1948. DOI: 10.1109/ICNN.1995.488968

\bibitem{CMAES}Hansen, N., \& Ostermeier, A. (2001). Completely derandomized
self-adaptation in evolution strategies. Evolutionary Computation,
9(2), 159--195. DOI: 10.1162/106365601750190398

\bibitem{LMCMAES}Loshchilov, I. (2014). A computationally efficient
limited memory CMA-ES for large scale optimization. Proceedings of
the Genetic and Evolutionary Computation Conference (GECCO), 397--404.
DOI: 10.1145/2576768.2598294

\bibitem{InitialDistribution}Charilogis, V., Tsoulos, I. G., \& Stavrou,
V. N. (2023). An Intelligent Technique for Initial Distribution of
Genetic Algorithms. Axioms, 12(10), 980. DOI: 10.3390/axioms12100980

\bibitem{DE}Storn, R., \& Price, K. (1997). Differential evolution
--- A simple and efficient heuristic for global optimization over
continuous spaces. Journal of Global Optimization, 11(4), 341--359.
DOI: 10.1023/A:1008202821328

\bibitem{DE2}Das, S., \& Suganthan, P. N. (2011). Differential evolution:
A survey of the state-of-the-art. IEEE Transactions on Evolutionary
Computation, 15(1), 4--31. DOI: 10.1109/TEVC.2010.2059031

\bibitem{JADE}Zhang, J., \& Sanderson, A. C. (2009). JADE: Adaptive
differential evolution with optional external archive. IEEE Transactions
on Evolutionary Computation, 13(5), 945--958. DOI: 10.1109/TEVC.2009.2014613

\bibitem{SHADE}Tanabe, R., \& Fukunaga, A. (2013). Success-history
based parameter adaptation for differential evolution. Proceedings
of the IEEE Congress on Evolutionary Computation (CEC), 71--78. DOI:
10.1109/CEC.2013.6557555

\bibitem{LSHADE}Tanabe, R., \& Fukunaga, A. S. (2014). Improving
the search performance of SHADE using linear population size reduction.
Proceedings of the IEEE Congress on Evolutionary Computation (CEC),
1658--1665. DOI: 10.1109/CEC.2014.6900380

\bibitem{DEd}Brest, J., \& Maučec, M. S. (2008). Population size
reduction for the differential evolution algorithm. Applied Intelligence,
29(3), 228--247.

\bibitem{JSO}Brest, J., Maučec, M. S., \& Bošković, B. (2017). Single
objective real-parameter optimization: Algorithm jSO. Proceedings
of the IEEE Congress on Evolutionary Computation (CEC), 1311--1318.
DOI: 10.1109/CEC.2017.7969456

\bibitem{LSHADE2}Stanovov, V., Akhmedova, S., \& Semenkin, E. (2018).
LSHADE Algorithm with Rank-Based Selective Pressure Strategy for Solving
CEC 2017 Benchmark Problems. Proceedings of the IEEE Congress on Evolutionary
Computation (CEC), 1--8. DOI: 10.1109/CEC.2018.8477977

\bibitem{NL-SHADE-RSP}Stanovov, V., Akhmedova, S., \& Semenkin, E.
(2021). NL-SHADE-RSP algorithm with adaptive archive and selective
pressure for CEC 2021 numerical optimization. Proceedings of the IEEE
Congress on Evolutionary Computation (CEC), 809--816. DOI: 10.1109/CEC45853.2021.9504959

\bibitem{CoBiDE}Wang, Y., Li, H.-X., Huang, T., \& Li, L. (2014).
Differential evolution based on covariance matrix learning and bimodal
distribution parameter setting. Applied Soft Computing, 18, 232--247.

\bibitem{DEp}Bujok, P., \& Tvrdík, J. (2017). Enhanced individual-dependent
differential evolution with population size adaptation. Proceedings
of the IEEE Congress on Evolutionary Computation (CEC).

\bibitem{SPS-L-SHADE-EIG}Guo, S.-M., Tsai, J.-H., Yang, C.-C., \&
Hsu, P.-H. (2015). A self-optimization approach for L-SHADE incorporated
with eigenvector-based crossover and successful-parent-selecting framework
on CEC 2015 benchmark set. Proceedings of the IEEE Congress on Evolutionary
Computation (CEC).

\bibitem{EI4Eig}Bujok, P., \& Kolenovský, P. (2022). Eigen crossover
in cooperative model of evolutionary algorithms applied to CEC 2022
single objective numerical optimisation. Proceedings of the IEEE Congress
on Evolutionary Computation (CEC), 1--8. DOI: 10.1109/CEC55065.2022.9870433

\bibitem{obl}Tizhoosh, H. R. (2005). Opposition-based learning: A
new scheme for machine intelligence. Proceedings of the International
Conference on Computational Intelligence for Modelling, Control and
Automation (CIMCA), 695--701. DOI: 10.1109/CIMCA.2005.1631345

\bibitem{obde}Rahnamayan, S., Tizhoosh, H. R., \& Salama, M. M. A.
(2008). Opposition-based differential evolution. IEEE Transactions
on Evolutionary Computation, 12(1), 64--79. DOI: 10.1109/TEVC.2007.894200

\bibitem{Mantegna}Mantegna, R. N. (1994). Fast, accurate algorithm
for numerical simulation of Lévy stable stochastic processes. Physical
Review E, 49(5), 4677--4683. DOI: 10.1103/PhysRevE.49.4677

\bibitem{Thompson}Thompson, W. R. (1933). On the likelihood that
one unknown probability exceeds another in view of the evidence of
two samples. Biometrika, 25(3/4), 285--294. DOI: 10.1093/biomet/25.3-4.285

\bibitem{ARQ}Charilogis, V., Tsoulos, I. G., Gianni, A. M., \& Tsalikakis,
D. (2025). ARQ: A cohesive optimization design for stable performance
on noisy landscapes. Applied Sciences, 15(22), 12180. DOI: 10.3390/app152212180

\bibitem{ARQ2}Charilogis, V., Tsoulos, I. G., \& Gianni, A. M. (2026).
ARQ2: Toward stability-aware hybrid optimization on complex and noisy
search problems. Symmetry, 18(5), 844. DOI: 10.3390/sym18050844

\bibitem{optimsolution}Charilogis, V., Tsoulos, I. G., \& Gianni,
A. M. (2026). OptimSolution: A Cross-Platform Framework for Benchmarking,
Sensitivity, and Complexity Analysis of Continuous Optimisation Methods.
Software, 5(3), 28. DOI: 10.3390/software5030028

\bibitem{de_CEC2011}Elsayed, S. M., Sarker, R. A., \& Essam, D. L.
(2011). Differential evolution with multiple strategies for solving
CEC2011 real-world numerical optimization problems. In 2011 IEEE Congress
on Evolutionary Computation (pp. 1041-1048). IEEE. https://doi.org/10.1109/CEC.2011.5949732

\bibitem{mLSHADE-RL}Chauhan, D., Trivedi, A., \& Shivani. (2024).
A Multi-operator Ensemble LSHADE with Restart and Local Search Mechanisms
for Single-objective Optimization. arXiv:2409.15994. DOI: 10.48550/arXiv.2409.15994

\bibitem{NL-SHADE-LBC}Stanovov, V., Akhmedova, S., \& Semenkin, E.
(2022). NL-SHADE-LBC algorithm with linear parameter adaptation bias
change for CEC 2022 Numerical Optimization. Proceedings of the IEEE
Congress on Evolutionary Computation (CEC), 1--8. DOI: 10.1109/CEC55065.2022.9870295

\bibitem{SFCDE}Wang, Z., Li, D., Yu, J., Abdel-Salam, M., Hu, G.,
Houssein, E. H., Abdel Samee, N., \& Zhong, R. (2026). Competitive
differential evolution with success-failure adaptation mechanism:
performance benchmarking and application in eggplant disease diagnosis.
International Journal of Machine Learning and Cybernetics, 17, 155.
DOI: 10.1007/s13042-025-02931-3

\bibitem{TridentDE}Charilogis, V., Tsoulos, I. G., \& Gianni, A.
M. (2025). TRIDENT-DE: Triple-Operator Differential Evolution with
Adaptive Restarts and Greedy Refinement. Future Internet, 17(11),
488. DOI: 10.3390/fi17110488

\bibitem{UDE3}Trivedi, A., \& Chauhan, D. (2024). UDE-III: An Enhanced
Unified Differential Evolution Algorithm for Constrained Optimization
Problems. arXiv:2410.03992. DOI: 10.48550/arXiv.2410.03992

\bibitem{cec2011}Das, S., \& Suganthan, P. N. (2011). Problem Definitions
and Evaluation Criteria for CEC 2011 Competition on Testing Evolutionary
Algorithms on Real World Optimization Problems. Technical Report,
Jadavpur University, Kolkata, India, and Nanyang Technological University,
Singapore.

\bibitem{vibratingplatform}Thanedar, P. B., \& Vanderplaats, G. N.
(1995). Survey of discrete variable optimization for structural design.
Journal of Structural Engineering, 121(2), 301--305. DOI: 10.1061/(ASCE)0733-9445(1995)121:2(301)

\bibitem{weatherirrigation}Charilogis, V., Tsoulos, I. G., Gianni,
A. M., \& Tsalikakis, D. (2026). A Novel Weather-Aware Irrigation
Scheduling Benchmark for Continuous Global Optimization. Mathematics,
14(11), 2018. DOI: 10.3390/math14112018

\bibitem{CCDG2_1}Omidvar, M. N., Li, X., Mei, Y., \& Yao, X. (2014).
Cooperative Co-evolution with Differential Grouping for Large Scale
Optimization. IEEE Transactions on Evolutionary Computation, 18(3),
378--393. DOI: 10.1109/TEVC.2013.2281543

\bibitem{CCDG2_2}Omidvar, M. N., Yang, M., Mei, Y., Li, X., \& Yao,
X. (2017). DG2: A Faster and More Accurate Differential Grouping for
Large-Scale Black-Box Optimization. IEEE Transactions on Evolutionary
Computation, 21(6), 929--942. DOI: 10.1109/TEVC.2017.2694221

\bibitem{MMES}He, X., Zheng, Z., \& Zhou, Y. (2021). MMES: Mixture
Model based Evolution Strategy for Large-Scale Optimization. IEEE
Transactions on Evolutionary Computation, 25(2), 320--333. DOI: 10.1109/TEVC.2020.3034769

\bibitem{RMES}Li, Z., \& Zhang, Q. (2018). A Simple Yet Efficient
Evolution Strategy for Large-Scale Black-Box Optimization. IEEE Transactions
on Evolutionary Computation, 22(5), 637--646. DOI: 10.1109/TEVC.2017.2765682

\bibitem{VkD-CMA-ESc}Akimoto, Y., \& Hansen, N. (2016). Projection-based
restricted covariance matrix adaptation for high dimension. Proceedings
of the Genetic and Evolutionary Computation Conference (GECCO), 197--204.
DOI: 10.1145/2908812.2908863

\end{thebibliography}

\end{adjustwidth}{}
\end{document}
